% El mundo rural, la movilidad y el desarrollo local — Espagnol 1ère A — Cours signé OnBuch+
% Programme officiel MINESEC. Compiler avec : tectonic EA10-mundo-rural-movilidad-desarrollo-local.tex
\newcommand{\DOCMATIERE}{Espagnol}
\newcommand{\DOCNIVEAU}{1ère A}
\newcommand{\DOCMODULE}{Módulo 3 : Citoyenneté}
\newcommand{\DOCLECON}{EA10}
\newcommand{\DOCTITRE}{El mundo rural, la movilidad y el desarrollo local}
% OnBuch+ — préambule « cours signés » ENRICHI (compilé avec tectonic / XeLaTeX)
% Système visuel « Pop » d'OnBuch+ : fond crème (jamais blanc), bordures encre
% épaisses, ombres « dures » décalées, titres Archivo Black en capitales.
% Version MATHÉMATIQUES : graphes (pgfplots/tikz), boîtes pédagogiques
% (définition, propriété, méthode, attention, le savais-tu, exercice résolu,
% exercice, TP, bilan), page de garde et signature OnBuch+ ancrée en bas.
%
% Macros attendues dans main.tex AVANT \input{preamble} :
%   \DOCMATIERE  \DOCNIVEAU  \DOCMODULE  \DOCLECON (numéro)  \DOCTITRE
\documentclass[11pt]{article}

\usepackage{fontspec}
\usepackage[spanish,french]{babel} % français = langue principale ; l'espagnol via \esp{..} / espagnol
\usepackage[a4paper,margin=2cm,top=3.4cm,bottom=2.8cm,headheight=1.4cm,footskip=1.3cm]{geometry}
\usepackage{amsmath,amssymb,amsfonts}
\usepackage{mathtools} % \xmapsto, \coloneqq... souvent utilisés par les modèles
\usepackage{cancel} % simplifications barrées (bilans de forces)
% \abs{z} (module, valeur absolue) et \norm{u} (norme), utilisés par les modèles
\providecommand{\abs}{}\let\abs\relax\DeclarePairedDelimiter{\abs}{\lvert}{\rvert}
\providecommand{\norm}{}\let\norm\relax\DeclarePairedDelimiter{\norm}{\lVert}{\rVert}
\usepackage[table]{xcolor} % [table] : fournit aussi \rowcolors (alternance de lignes)
\usepackage{pagecolor}
\usepackage{tikz}
\usetikzlibrary{arrows.meta,positioning,calc,shapes.geometric,shapes.misc,shapes.arrows,decorations.pathmorphing,decorations.pathreplacing,decorations.markings,patterns,shadows,mindmap,trees,fit,backgrounds,matrix,intersections,angles,quotes}
\usepackage{pgfplots}
\usepackage[version=4]{mhchem} % équations nucléaires \ce{^{235}_{92}U}
\usepackage[european,siunitx]{circuitikz} % schémas électriques
\pgfplotsset{compat=1.18}
\usepgfplotslibrary{fillbetween,groupplots} % aires entre courbes (calcul intégral)
\usepackage{enumitem}
\usepackage{fancyhdr}
% [most] charge toutes les bibliothèques tcolorbox (listings, minted, raster,
% documentation...) et gonfle inutilement la mémoire TeX sur un document long
% (~300 boîtes) : on ne charge que ce qui est réellement utilisé.
\usepackage[skins,breakable]{tcolorbox}
\usepackage{graphicx}
\usepackage{colortbl}
\usepackage{booktabs}
\usepackage{tabularx}
\usepackage{diagbox} % case d'en-tête diagonale des tableaux à double entrée
% Colonnes extensibles centrée / gauche / droite (C, L, R), souvent utilisées par les modèles
\newcolumntype{C}{>{\centering\arraybackslash}X}
\newcolumntype{L}{>{\raggedright\arraybackslash}X}
\newcolumntype{R}{>{\raggedleft\arraybackslash}X}
\usepackage{array}
\usepackage{multicol}
\usepackage{siunitx}
% parse-numbers=false : accepte aussi \SI{2,5 \times 10^{-3}}{...}
\sisetup{locale=FR,output-decimal-marker={,},group-minimum-digits=4,parse-numbers=false}
\DeclareSIUnit\ohm{\ensuremath{\symup{\Omega}}} % U+2126 absent de la police
\DeclareSIUnit\division{div} % réglages d'oscilloscope (ms/div, V/div)
\DeclareSIUnit\tr{tr} % tours (vitesse de rotation, ex. \SI{1500}{\tr\per\minute})
\DeclareSIUnit\molar{M} % concentration molaire (ex. \SI{3}{\molar}, \SI{1}{\micro\molar})
\DeclareSIUnit\an{an} % année (le modèle confond parfois avec \year, un compteur TeX) — ex. \SI{5}{\kilo\meter\per\an}
\DeclareSIUnit\FCFA{FCFA} % franc CFA (ex. \SI{500}{\FCFA\per\kilo\gram})
\DeclareSIUnit\degreeBrix{\textdegree Brix} % teneur en sucre d'un sirop/jus (réfractométrie)
\DeclareSIUnit\month{mois} % le modèle utilise parfois \month, un compteur TeX, comme unité de durée
\usepackage{qrcode}
% \degree : commande fréquemment utilisée par les modèles pour un angle ou une
% température en dehors de \SI{}{} (non fournie par les paquets chargés).
\providecommand{\degree}{^{\circ}}
% \qed (fin de démonstration), utilisé par les modèles sans amsthm
\providecommand{\qed}{\hfill$\square$}
% \barre : double barre des valeurs interdites dans un tableau de variations (tkz-tab)
\providecommand{\barre}{\ensuremath{\|}}
% environnement proof (amsthm non chargé) utilisé par les modèles
\ifdefined\proof\else\newenvironment{proof}[1][Démonstration]{\par\noindent\textit{#1.}\ }{\qed\par}\fi
\usepackage{lastpage}
\usepackage{hyperref}
\hypersetup{hidelinks}

% ============================================================
% Palette « Pop » OnBuch+ — mêmes hex que la classe P (plus_ui.dart)
% ============================================================
\definecolor{poporange}{HTML}{F59321}
\definecolor{popdark}{HTML}{E07A0C}
\definecolor{popcream}{HTML}{FFF6E8}
\definecolor{popink}{HTML}{1C1714}
\definecolor{poppurple}{HTML}{7A5AE0}
\definecolor{popgreen}{HTML}{2E9E6A}
\definecolor{poppink}{HTML}{E0457B}
\definecolor{popblue}{HTML}{2F7DE1}
\definecolor{popgold}{HTML}{C98A12}
\definecolor{popmuted}{HTML}{8A7F76}
\definecolor{popline}{HTML}{EADFD2}
\definecolor{popgray}{HTML}{8A7F76}
\colorlet{popgrey}{popgray}
% Alias tolérants (le modèle nomme parfois les couleurs différemment)
\colorlet{popwhite}{white}
\colorlet{popblack}{popink}
\colorlet{popred}{poppink}
\colorlet{popyellow}{popgold}
% Teintes claires (fonds de tableaux/encadrés) — le modèle nomme parfois
% les couleurs avec un suffixe L (ex. popblueL) sans les avoir définies.
\colorlet{poporangeL}{poporange!20}
\colorlet{popdarkL}{popdark!20}
\colorlet{poppurpleL}{poppurple!20}
\colorlet{popgreenL}{popgreen!20}
\colorlet{poppinkL}{poppink!20}
\colorlet{popblueL}{popblue!20}
\colorlet{popgoldL}{popgold!20}
\colorlet{popredL}{popred!20}
\colorlet{popgrayL}{popgray!20}
% Teintes claires pour fonds de tableaux / zones
\colorlet{popblueL}{popblue!10!white}
\colorlet{poporangeL}{poporange!14!white}
\colorlet{popgreenL}{popgreen!12!white}
\colorlet{poppurpleL}{poppurple!12!white}
\colorlet{poppinkL}{poppink!12!white}
\colorlet{popgoldL}{popgold!14!white}

\pagecolor{popcream}

% ============================================================
% Typographie
% ============================================================
\defaultfontfeatures{Ligatures=TeX}
\setmainfont{PlusJakartaSans-Medium.ttf}[Path=../../../fonts/,BoldFont=PlusJakartaSans-Bold.ttf]
% Maths en sans-serif (Fira Math) avec chiffres et lettres droites de Plus
% Jakarta, pour que les formules se fondent dans le texte.
\usepackage[math-style=ISO,bold-style=ISO]{unicode-math}
\setmathfont{FiraMath-Regular.otf}[Path=../../../fonts/]
\setmathfont{PlusJakartaSans-Medium.ttf}[Path=../../../fonts/,range={up/num,up/{Latin,latin},\mathup}]
\usepackage{icomma} % virgule décimale sans espace en mode math
% Caractères mathématiques Unicode absents des polices de texte (Plus Jakarta,
% Archivo Black) : rendus par leur équivalent mathématique, en texte comme en
% formule — sinon ils s'affichent en carré vide (ex. « Arithmétique dans ℤ »).
\usepackage{newunicodechar}
\newunicodechar{ℝ}{\ensuremath{\mathbb{R}}}
\newunicodechar{ℤ}{\ensuremath{\mathbb{Z}}}
\newunicodechar{ℂ}{\ensuremath{\mathbb{C}}}
\newunicodechar{ℕ}{\ensuremath{\mathbb{N}}}
\newunicodechar{ℚ}{\ensuremath{\mathbb{Q}}}
\newunicodechar{𝔻}{\ensuremath{\mathbb{D}}}
\newunicodechar{α}{\ensuremath{\alpha}}
\newunicodechar{β}{\ensuremath{\beta}}
\newunicodechar{γ}{\ensuremath{\gamma}}
\newunicodechar{δ}{\ensuremath{\delta}}
\newunicodechar{ε}{\ensuremath{\varepsilon}}
\newunicodechar{θ}{\ensuremath{\theta}}
\newunicodechar{λ}{\ensuremath{\lambda}}
\newunicodechar{μ}{\ensuremath{\mu}}
\newunicodechar{σ}{\ensuremath{\sigma}}
\newunicodechar{τ}{\ensuremath{\tau}}
\newunicodechar{φ}{\ensuremath{\varphi}}
\newunicodechar{ψ}{\ensuremath{\psi}}
\newunicodechar{ω}{\ensuremath{\omega}}
\newunicodechar{Σ}{\ensuremath{\Sigma}}
% majuscules produites par \MakeUppercase dans les titres (α-aminés -> Α)
\newunicodechar{Α}{\ensuremath{\alpha}}
\newunicodechar{Β}{\ensuremath{\beta}}
\newunicodechar{Γ}{\ensuremath{\Gamma}}
\newunicodechar{Δ}{\ensuremath{\Delta}}
\newunicodechar{Ε}{\ensuremath{\varepsilon}}
\newunicodechar{Θ}{\ensuremath{\Theta}}
\newunicodechar{Λ}{\ensuremath{\Lambda}}
\newunicodechar{Μ}{\ensuremath{\mu}}
\newunicodechar{Τ}{\ensuremath{\tau}}
\newunicodechar{Φ}{\ensuremath{\Phi}}
\newunicodechar{Ψ}{\ensuremath{\Psi}}
\newunicodechar{Ω}{\ensuremath{\Omega}}
\newunicodechar{↦}{\ensuremath{\mapsto}}
\newunicodechar{∈}{\ensuremath{\in}}
\newunicodechar{∉}{\ensuremath{\notin}}
\newunicodechar{∧}{\ensuremath{\wedge}}
\newunicodechar{⇔}{\ensuremath{\Leftrightarrow}}
\newunicodechar{⟺}{\ensuremath{\Longleftrightarrow}}
\newunicodechar{⇒}{\ensuremath{\Rightarrow}}
\newunicodechar{⟹}{\ensuremath{\Longrightarrow}}
\newunicodechar{∘}{\ensuremath{\circ}}
\newunicodechar{∪}{\ensuremath{\cup}}
\newunicodechar{∩}{\ensuremath{\cap}}
\newunicodechar{≡}{\ensuremath{\equiv}}
\newunicodechar{⊂}{\ensuremath{\subset}}
\newunicodechar{⊥}{\ensuremath{\perp}}
\newunicodechar{∃}{\ensuremath{\exists}}
\newunicodechar{∀}{\ensuremath{\forall}}
\newunicodechar{∛}{\ensuremath{\sqrt[3]{\,}}}
\newunicodechar{∜}{\ensuremath{\sqrt[4]{\,}}}
\newunicodechar{⃗}{\ensuremath{{}^{\to}}}
\newunicodechar{ⁿ}{\ifmmode^{n}\else\textsuperscript{\ensuremath{n}}\fi}
\newunicodechar{ˣ}{\ifmmode^{x}\else\textsuperscript{\ensuremath{x}}\fi}
\newunicodechar{ᵇ}{\ifmmode^{b}\else\textsuperscript{\ensuremath{b}}\fi}
\newunicodechar{ʳ}{\ifmmode^{r}\else\textsuperscript{\ensuremath{r}}\fi}
\newunicodechar{ᵘ}{\ifmmode^{u}\else\textsuperscript{\ensuremath{u}}\fi}
\newunicodechar{ʸ}{\ifmmode^{y}\else\textsuperscript{\ensuremath{y}}\fi}
\newunicodechar{ᵃ}{\ifmmode^{a}\else\textsuperscript{\ensuremath{a}}\fi}
\newunicodechar{ᵉ}{\ifmmode^{e}\else\textsuperscript{\ensuremath{e}}\fi}
\newunicodechar{ᵏ}{\ifmmode^{k}\else\textsuperscript{\ensuremath{k}}\fi}
\newunicodechar{ᵖ}{\ifmmode^{p}\else\textsuperscript{\ensuremath{p}}\fi}
\newunicodechar{ᴺ}{\ifmmode^{N}\else\textsuperscript{\ensuremath{N}}\fi}
\newunicodechar{ᶜ}{\ifmmode^{c}\else\textsuperscript{\ensuremath{c}}\fi}
\newunicodechar{ʰ}{\ifmmode^{h}\else\textsuperscript{\ensuremath{h}}\fi}
\newunicodechar{ᵠ}{\ifmmode^{q}\else\textsuperscript{\ensuremath{q}}\fi}
\newunicodechar{ₙ}{\ifmmode_{n}\else\textsubscript{\ensuremath{n}}\fi}
\newunicodechar{ₐ}{\ifmmode_{a}\else\textsubscript{\ensuremath{a}}\fi}
\newunicodechar{ᵢ}{\ifmmode_{i}\else\textsubscript{\ensuremath{i}}\fi}
\newunicodechar{ᵣ}{\ifmmode_{r}\else\textsubscript{\ensuremath{r}}\fi}
\newunicodechar{ₖ}{\ifmmode_{k}\else\textsubscript{\ensuremath{k}}\fi}
\newunicodechar{ᵦ}{\ifmmode_{\beta}\else\textsubscript{\ensuremath{\beta}}\fi}
\newunicodechar{ₓ}{\ifmmode_{x}\else\textsubscript{\ensuremath{x}}\fi}
\newfontfamily\popdisplay{ArchivoBlack-Regular.ttf}[Path=../../../fonts/]
\newfontfamily\poplogo{SpaceGrotesk-Bold.ttf}[Path=../../../fonts/]
\newfontfamily\popsemi{PlusJakartaSans-SemiBold.ttf}[Path=../../../fonts/]
\newfontfamily\popxbold{PlusJakartaSans-ExtraBold.ttf}[Path=../../../fonts/]
\newfontfamily\popxboldls{PlusJakartaSans-ExtraBold.ttf}[Path=../../../fonts/,LetterSpace=3.0]

\setlist[enumerate]{leftmargin=1.7em,itemsep=5pt,topsep=4pt}
\setlist[itemize]{leftmargin=1.7em,itemsep=4pt,topsep=4pt,label={\color{poporange}\textbullet}}
\setlength{\parindent}{0pt}
\setlength{\parskip}{5pt}
\renewcommand{\arraystretch}{1.25}

% pgfplots : style Pop par défaut
\pgfplotsset{
  every axis/.append style={axis line style={popink,line width=1pt},tick style={popink},
    grid style={popline,line width=0.5pt},label style={font=\small},tick label style={font=\footnotesize}},
  popcurve/.style={poporange,line width=1.8pt,no markers,smooth},
}
% Même style disponible dans un \draw TikZ simple (hors pgfplots)
\tikzset{popcurve/.style={poporange,line width=1.8pt}}
\tikzset{popgrid/.style={popline,very thin}}
\colorlet{popcurve}{poporange} % le nom du style est parfois utilisé comme couleur

% ============================================================
% Pill / squiggle
% ============================================================
\newcommand{\poppill}[3]{%
  \tikz[baseline=-0.55ex]\node[draw=popink,line width=1.2pt,rounded corners=2.6pt,%
    fill=#2,inner xsep=6.5pt,inner ysep=3pt]{%
    \popxboldls\fontsize{7}{7}\selectfont\color{#3}\MakeUppercase{#1}};%
}
\newcommand{\popsquiggle}[1]{%
  \tikz[baseline=-0.4ex]{
    \draw[#1,line width=2.6pt,line cap=round]
      plot [smooth] coordinates {(0,0.09) (0.34,0.01) (0.68,0.21) (1.02,0.01) (1.36,0.10)};
  }%
}
% Mot-clé surligné (marqueur orange sous le texte)
\newcommand{\cle}[1]{{\popxbold\color{popdark}#1}}

% ============================================================
% En-tête / pied
% ============================================================
\pagestyle{fancy}
\fancyhf{}
\renewcommand{\headrulewidth}{0pt}
\renewcommand{\footrulewidth}{0pt}
\lhead{\raisebox{-0.15ex}{\poplogo\fontsize{13}{13}\selectfont\color{popink}OnBuch\color{poporange}+}}
\rhead{\poppill{\DOCMATIERE\ \textbullet\ \DOCNIVEAU\ \textbullet\ Leçon \DOCLECON}{white}{popink}}
\lfoot{\popxbold\fontsize{7.6}{7.6}\selectfont\color{popmuted}\MakeUppercase{Cours signé OnBuch+}\ \textbullet\ {\popsemi\DOCTITRE}}
\rfoot{\tikz[baseline=-0.6ex]\node[rounded corners=8.5pt,fill=popink,minimum height=17pt,minimum width=17pt,inner xsep=5pt,inner sep=0]{\popxbold\fontsize{8}{8}\selectfont\color{popcream}\thepage\,/\,\pageref*{LastPage}};}

% ============================================================
% Titres
% ============================================================
\newcommand{\chaptitre}[2]{%
  \begin{tcolorbox}[enhanced,colback=poporange,colframe=popink,boxrule=2.6pt,arc=11pt,
      drop shadow=popink,left=15pt,right=15pt,top=12pt,bottom=11pt,before skip=3pt,after skip=18pt]
    \poppill{#2}{popink}{popcream}\par\vspace{9pt}
    {\raggedright\hyphenpenalty=10000\exhyphenpenalty=10000\popdisplay\fontsize{22}{24}\selectfont\color{popcream}\MakeUppercase{#1}\par}
  \end{tcolorbox}
}
% Section numérotée : pastille orange avec le numéro + titre Archivo
\newcounter{popsec}
\newcommand{\coursec}[1]{%
  \stepcounter{popsec}\setcounter{popsub}{0}%
  \par\vspace{16pt}\noindent
  \begin{minipage}{\linewidth}
  \tikz[baseline=-0.7ex]\node[draw=popink,line width=1.6pt,fill=poporange,rounded corners=4pt,minimum size=22pt,inner sep=0]{\popdisplay\fontsize{12}{12}\selectfont\color{popcream}\thepopsec};%
  \hspace{8pt}{\popdisplay\fontsize{15}{17}\selectfont\color{popink}\MakeUppercase{#1}}\par
  \vspace{4pt}\noindent{\color{poporange}\rule{36pt}{3.6pt}}
  \end{minipage}\par\nopagebreak\vspace{8pt}
}
\newcounter{popsub}
\newcommand{\courssub}[1]{%
  \stepcounter{popsub}%
  \par\vspace{9pt}\noindent{\popxbold\fontsize{11.5}{13}\selectfont\color{popink}{\color{poporange}\thepopsec.\thepopsub}\hspace{6pt}#1}\par\nopagebreak\vspace{4pt}
}

% ============================================================
% Boîtes pédagogiques — même recette : fond blanc, bordure épaisse colorée,
% ombre dure encre, badge plein. Titre optionnel [#1].
% ============================================================
\tcbset{popbase/.style={breakable,enhanced,colback=white,boxrule=2.3pt,arc=9pt,
  shadow={3pt}{-3pt}{0pt}{popink},left=14pt,right=14pt,top=11pt,bottom=11pt,before skip=11pt,after skip=15pt,
  fontupper=\popsemi\fontsize{10.6}{14.5}\selectfont\color{popink},
  fontlower=\popsemi\fontsize{10.6}{14.5}\selectfont\color{popink}}}
% Coche dessinée (le glyphe ✓ n'existe pas dans les polices du document)
\newcommand{\popcheck}[1]{\tikz[baseline=-0.5ex]\draw[#1,line width=1.6pt,line cap=round,line join=round](-0.13,0.0)--(-0.04,-0.09)--(0.13,0.1);}
\providecommand{\coche}{\popcheck{popgreen}}
\providecommand{\resultat}[1]{\par\smallskip\noindent\fcolorbox{popgreen}{popgreenL}{\parbox{\dimexpr\linewidth-2\fboxsep-2\fboxrule\relax}{\textbf{Résultat.} #1}}\par\smallskip}
\newcommand{\popboxhead}[3]{\poppill{#1}{#2}{white}\ifx\relax#3\relax\else\hspace{6pt}{\popxbold\fontsize{10.6}{12}\selectfont\color{popink}#3}\fi\par\vspace{7pt}}

\newtcolorbox{aretenir}[1][]{popbase,colframe=popgreen,before upper={\popboxhead{À retenir}{popgreen}{#1}}}
% Texte en espagnol (document de lecture, dialogue, poème, extrait) : cadre
% d'étude. Un titre optionnel (source, auteur) s'affiche dans l'en-tête.
\newtcolorbox{textebox}[1][]{popbase,colframe=popblue,before upper={\popboxhead{Texto}{popblue}{#1}\begin{otherlanguage}{spanish}},after upper={\end{otherlanguage}}}
% Marques ✓ / ✗ (forme correcte / fautive) souvent utilisées par les modèles
\providecommand{\cross}{\textcolor{poppink}{\ensuremath{\times}}}
\providecommand{\xmark}{\cross}\providecommand{\crossmark}{\cross}\providecommand{\cmark}{\ensuremath{\checkmark}}
% Ligne de réponse / justification à compléter (inventée par les modèles)
\providecommand{\justification}[1]{\par\noindent\textit{Justification~:} \dotfill\par}
% \textipa (package tipa non chargé) : la police ne couvre pas l'API -> simple passage
\providecommand{\textipa}[1]{#1}
% Soulignement (ulem non chargé) : \uline, \ul
\providecommand{\uline}[1]{\underline{#1}}\providecommand{\ul}[1]{\underline{#1}}
% Ordinaux français écrits en macro par les modèles : 1\ière, 3\ième
\newcommand{\ière}{\textsuperscript{re}}\newcommand{\ième}{\textsuperscript{e}}
% Mot ou phrase en espagnol à l'intérieur d'un paragraphe français
\newcommand{\esp}[1]{\ifmmode\text{\textit{#1}}\else\begingroup\selectlanguage{spanish}\itshape #1\endgroup\fi}
\newtcolorbox{exemplebox}[1][]{popbase,colframe=poporange,before upper={\popboxhead{Exemple}{poporange}{#1}}}
\newtcolorbox{definition}[1][]{popbase,colframe=popblue,before upper={\popboxhead{Définition}{popblue}{#1}}}
\newtcolorbox{propriete}[1][]{popbase,colframe=poppurple,before upper={\popboxhead{Propriété}{poppurple}{#1}}}
\newtcolorbox{methode}[1][]{popbase,colframe=popgold,before upper={\popboxhead{Méthode}{popgold}{#1}}}
\newtcolorbox{attention}[1][]{popbase,colframe=poppink,before upper={\popboxhead{Attention}{poppink}{#1}}}
\newtcolorbox{experience}[1][]{popbase,colframe=popblue,colback=popblueL,before upper={\popboxhead{Expérience}{popblue}{#1}}}
% « Le savais-tu ? » : carte violette pleine (contexte, culture, Cameroun)
\newtcolorbox{savaistu}[1][]{popbase,colback=poppurple,colframe=popink,
  fontupper=\popsemi\fontsize{10.4}{14.2}\selectfont\color{white},
  before upper={\poppill{Le savais-tu\,?}{white}{poppurple}\ifx\relax#1\relax\else\hspace{6pt}{\popxbold\color{white}#1}\fi\par\vspace{7pt}}}
% Exercice résolu : énoncé en haut, \tcblower puis la solution sur fond crème
\newcounter{popexo}\newcounter{popexr}
\newtcolorbox{exoresolu}[1][]{popbase,colframe=poporange,colbacklower=popcream,
  segmentation style={popink,line width=1.2pt,dashed},
  before upper={\stepcounter{popexr}\popboxhead{Exercice résolu \thepopexr}{poporange}{#1}},
  before lower={\poppill{Solution}{popink}{popcream}\par\vspace{6pt}}}
% Exercice (non résolu en place) — la correction est dans la partie Corrigés
\newtcolorbox{exercice}[1][]{popbase,colframe=popink,
  before upper={\stepcounter{popexo}\popboxhead{Exercice \thepopexo}{popink}{#1}}}
% Corrigé
\newtcolorbox{corrige}[1][]{popbase,colframe=popgreen,colback=popgreenL,
  before upper={\popboxhead{Corrigé}{popgreen}{#1}}}
% Objectifs / prérequis / bilan
\newtcolorbox{objectifs}{popbase,colframe=popink,colback=poporangeL,
  before upper={\popboxhead{Objectifs de la leçon}{popink}{}}}
\newtcolorbox{prerequis}{popbase,colframe=popink,colback=popblueL,
  before upper={\popboxhead{Prérequis}{popblue}{}}}
\newtcolorbox{bilan}{popbase,colframe=popink,colback=white,boxrule=2.8pt,
  before upper={\popboxhead{Fiche bilan}{poporange}{L'essentiel à connaître pour l'examen}}}
% Figure encadrée avec légende
\newtcolorbox{popfigure}[1][]{enhanced,breakable=false,colback=white,colframe=popline,boxrule=1.4pt,arc=8pt,
  left=10pt,right=10pt,top=10pt,bottom=8pt,before skip=10pt,after skip=12pt,halign=center,
  lower separated=false,halign lower=center,
  fontlower=\popsemi\fontsize{9}{11.5}\selectfont\color{popmuted},
  #1}
\newcommand{\legende}[1]{\par\vspace{6pt}{\popsemi\fontsize{9}{11.5}\selectfont\color{popmuted}#1}}

% ============================================================
% Signature OnBuch+
% ============================================================
\newcommand{\poplogoinline}[1]{{\poplogo\fontsize{#1}{#1}\selectfont\color{popink}OnBuch\color{poporange}+}}
\newcommand{\signatureonbuch}{%
  \begin{tcolorbox}[enhanced,colback=popink,colframe=popink,boxrule=2.2pt,arc=10pt,
      drop shadow=poporange,left=14pt,right=14pt,top=10pt,bottom=10pt,before skip=0pt,after skip=0pt]
    \tikz[baseline=-0.7ex]\node[circle,fill=popgreen,draw=popcream,line width=1.3pt,minimum size=22pt,inner sep=0]{\popcheck{white}};%
    \hspace{9pt}\begin{minipage}[c]{0.62\linewidth}
      {\popxbold\fontsize{12}{13}\selectfont\color{popcream}Signé\ \textbullet\ {\poplogo OnBuch\color{poporange}+}}\par\vspace{2pt}
      {\raggedright\popsemi\fontsize{8}{10.5}\selectfont\color{popline}\MakeUppercase{Équipe pédagogique OnBuch+}\\\MakeUppercase{Conforme au programme officiel MINESEC}\par}
    \end{minipage}\hfill
    {\poplogo\fontsize{22}{22}\selectfont\color{popcream}OnBuch\color{poporange}+}
  \end{tcolorbox}%
}

% ============================================================
% Page de garde
% #1 titre · #2 matière · #3 niveau · #4 module · #5 n° leçon · #6 durée ·
% #7 description / objectifs (texte ou itemize)
% ============================================================
\newcommand{\pagegarde}[7]{%
  \thispagestyle{empty}
  \newgeometry{a4paper,margin=1.7cm,top=1.6cm,bottom=1.4cm}%
  \begin{tikzpicture}[remember picture,overlay]
    \fill[poporange!18] ([xshift=-3.2cm,yshift=-3.6cm]current page.north east) circle (4.6cm);
    \fill[poppurple!14] ([xshift=2.4cm,yshift=7.6cm]current page.south west) circle (3.8cm);
  \end{tikzpicture}%
  {\poplogo\fontsize{30}{30}\selectfont\color{popink}OnBuch\color{poporange}+}\hfill
  \raisebox{0.6ex}{\poppill{Cours signé \textbullet\ Premium}{popink}{popcream}}\par
  \vspace{1pt}\popsquiggle{poppurple}\par
  \vspace{8pt}
  \begin{tcolorbox}[enhanced,colback=poporange,colframe=popink,boxrule=2.8pt,arc=13pt,
      drop shadow=popink,left=18pt,right=18pt,top=12pt,bottom=13pt,after skip=10pt]
    \poppill{#2 \textbullet\ #3}{popink}{popcream}\hspace{5pt}\poppill{#4}{white}{popink}\par\vspace{8pt}
    {\popdisplay\fontsize{14}{14}\selectfont\color{popink}LEÇON #5}\par\vspace{4pt}
    {\raggedright\hyphenpenalty=10000\exhyphenpenalty=10000\popdisplay\fontsize{25}{27}\selectfont\color{popcream}\MakeUppercase{#1}\par}
    \vspace{8pt}
    {\popxbold\fontsize{9.5}{9.5}\selectfont\color{popink}\MakeUppercase{Durée conseillée : #6}}
  \end{tcolorbox}
  \begin{tcolorbox}[enhanced,colback=white,colframe=popink,boxrule=2.2pt,arc=10pt,
      drop shadow=popink,left=16pt,right=16pt,top=10pt,bottom=10pt,after skip=8pt]
    \poppill{Dans cette leçon}{poppurple}{white}\par\vspace{5pt}
    \popsemi\fontsize{9.3}{12.6}\selectfont\color{popink}#7
  \end{tcolorbox}
  \noindent\begin{minipage}[t]{0.487\linewidth}
    \begin{tcolorbox}[enhanced,colback=popgreen,colframe=popink,boxrule=2.2pt,arc=10pt,
        drop shadow=popink,left=11pt,right=11pt,top=10pt,bottom=10pt,height=3.1cm,valign=center]
      \tcbox[colback=white,colframe=popink,boxrule=1.3pt,arc=5pt,boxsep=0pt,nobeforeafter,
        left=3pt,right=3pt,top=3pt,bottom=3pt,valign=center]{\qrcode[height=1.9cm]{https://play.google.com/store/apps/details?id=com.luvvix.onbuch&hl=fr}}%
      \hspace{8pt}\begin{minipage}[c]{0.5\linewidth}
        {\popdisplay\fontsize{11}{12.5}\selectfont\color{white}\MakeUppercase{Télécharge OnBuch}}\par\vspace{4pt}
        {\raggedright\popsemi\fontsize{8.4}{11}\selectfont\color{white}Gratuit \textbullet\ Google Play\\Tuteur IA, annales, concours, résultats\par}
      \end{minipage}
    \end{tcolorbox}
  \end{minipage}\hfill
  \begin{minipage}[t]{0.487\linewidth}
    \begin{tcolorbox}[enhanced,colback=poppurple,colframe=popink,boxrule=2.2pt,arc=10pt,
        drop shadow=popink,left=11pt,right=11pt,top=10pt,bottom=10pt,height=3.1cm,valign=center]
      \tikz[baseline=-0.7ex]\node[circle,fill=white,draw=popink,line width=1.3pt,minimum size=1.6cm,inner sep=0]{\popdisplay\fontsize{24}{24}\selectfont\color{poppurple}+};%
      \hspace{8pt}\begin{minipage}[c]{0.6\linewidth}
        {\popdisplay\fontsize{11}{12.5}\selectfont\color{white}\MakeUppercase{Passe à OnBuch+}}\par\vspace{4pt}
        {\raggedright\popsemi\fontsize{8.4}{11}\selectfont\color{white}Tous les cours signés, Tuteur IA illimité, fiches PDF — onglet OnBuch+ de l'app\par}
      \end{minipage}
    \end{tcolorbox}
  \end{minipage}
  \vspace{8pt}
  \popcontenu
  \vfill
  \signatureonbuch
  \restoregeometry
  \newpage
}

% Rangée « Ce cours contient » (page de garde)
\newcommand{\poptile}[3]{% #1 couleur #2 gros texte #3 légende
  \begin{tcolorbox}[enhanced,colback=#1,colframe=popink,boxrule=2pt,arc=9pt,shadow={2.5pt}{-2.5pt}{0pt}{popink},
      left=6pt,right=6pt,top=9pt,bottom=9pt,halign=center,valign=center,height=1.75cm,width=\linewidth,nobeforeafter]
    {\popdisplay\fontsize{12.5}{13}\selectfont\color{white}\MakeUppercase{#2}}\par\vspace{3pt}
    {\centering\popsemi\fontsize{7.8}{9.5}\selectfont\color{white}#3\par}
  \end{tcolorbox}}
\newcommand{\popcontenu}{%
  \par\noindent{\popxboldls\fontsize{7.5}{7.5}\selectfont\color{popmuted}CE COURS SIGNÉ CONTIENT}\par\vspace{7pt}
  \noindent\begin{minipage}[t]{0.235\linewidth}\poptile{popblue}{Cours}{complet, illustré, pas à pas}\end{minipage}\hfill
  \begin{minipage}[t]{0.235\linewidth}\poptile{poporange}{Méthodes}{et exercices résolus}\end{minipage}\hfill
  \begin{minipage}[t]{0.235\linewidth}\poptile{poppink}{Exercices}{type examen + corrigés}\end{minipage}\hfill
  \begin{minipage}[t]{0.235\linewidth}\poptile{popgreen}{Fiche bilan}{l'essentiel à retenir}\end{minipage}\par}

% Sommaire
\newtcolorbox{sommaire}{enhanced,colback=white,colframe=popink,boxrule=2.2pt,arc=10pt,
  drop shadow=popink,left=18pt,right=18pt,top=13pt,bottom=13pt,before skip=6pt,after skip=20pt,
  before upper={\poppill{Sommaire}{poppurple}{white}\par\vspace{9pt}\popsemi\fontsize{10.6}{14.5}\selectfont\color{popink}}}

% Signature de fin de cours — poussée tout en bas de la dernière page
\newcommand{\findecours}{%
  \par\vfill\nopagebreak
  \begin{center}\popsquiggle{poporange}\end{center}\vspace{4pt}
  \signatureonbuch
}
\AtBeginDocument{\def\llbracket{\mathopen{[\mkern-3.5mu[}}\def\rrbracket{\mathclose{]\mkern-3.5mu]}}} % intervalles d'entiers (glyphe ⟦ absent de la police)
% Glyphes absents de Fira Math : redéfinis à partir de glyphes disponibles
\AtBeginDocument{%
  \def\setminus{\mathbin{\backslash}}\let\smallsetminus\setminus
  \def\vdots{\mathord{\vbox{\baselineskip3.2pt\lineskiplimit0pt\kern4pt\hbox{.}\hbox{.}\hbox{.}}}}%
  \def\ddots{\mathinner{\mkern1mu\raise6.5pt\hbox{.}\mkern2mu\raise3.6pt\hbox{.}\mkern2mu\raise0.7pt\hbox{.}\mkern1mu}}%
  \def\triangle{\mathord{\scalebox{0.9}{$\symup{\Delta}$}}}%
  \def\nabla{\mathord{\raisebox{\depth}{\scalebox{1}[-1]{$\symup{\Delta}$}}}}%
  \def\checkmark{\popcheck{popdark}}%
  \def\star{\mathbin{\ast}}\def\bigstar{\mathord{\ast}}%
  \def\diamond{\mathbin{\scalebox{0.6}{$\Diamond$}}}%
  \def\bigcup{\mathop{\mathchoice{\raisebox{-0.3em}{\scalebox{1.7}{$\cup$}}}{\scalebox{1.25}{$\cup$}}{\cup}{\cup}}\displaylimits}%
  \def\bigcap{\mathop{\mathchoice{\raisebox{-0.3em}{\scalebox{1.7}{$\cap$}}}{\scalebox{1.25}{$\cap$}}{\cap}{\cap}}\displaylimits}%
  \def\lesseqqgtr{\mathrel{\lessgtr}}\def\bigoplus{\mathop{\oplus}\displaylimits}%
}
\providecommand{\ssi}{si et seulement si} % abréviation fréquente
\AtBeginDocument{\providecommand{\wideparen}{\overparen}} % arc de cercle

%% FIGFIX-BEGIN v4 — correctifs globaux des figures (docs/onbuch-generation/tools/figfix)
\usepackage{adjustbox}\usepackage{etoolbox}
% toute figure TikZ/pgfplots plus large que la ligne est réduite à la largeur disponible
\BeforeBeginEnvironment{tikzpicture}{\begin{adjustbox}{max width=\linewidth}}
\AfterEndEnvironment{tikzpicture}{\end{adjustbox}}
% légendes pgfplots lisibles par-dessus les courbes
\pgfplotsset{every axis/.append style={legend style={fill=white,fill opacity=0.92,text opacity=1,draw=black!15,inner sep=3pt}}}
% étiquettes d'arêtes (syntaxe edge["..."]) avec fond blanc
\tikzset{every edge quotes/.append style={fill=white,inner sep=1.2pt}}
%% FIGFIX-END
\begin{document}

\pagegarde{El mundo rural, la movilidad y el desarrollo local}{Espagnol}{1ère A}{Módulo 3 : Citoyenneté}{EA10}{2 h}{Cette leçon mixte du module Citoyenneté permet de comprendre un texte sur le monde rural, l'exode rural et le développement local, d'acquérir le lexique thématique et de maîtriser l'emploi combiné du prétérit indéfini et de l'imparfait pour raconter le passé.
\vspace{4pt}
\begin{multicols}{2}\small
\begin{itemize}
\item À la fin de cette leçon, je sais lire et comprendre un texte sur le monde rural et le développement local
\item À la fin de cette leçon, je sais utiliser le lexique du champ, de l'agriculture, de la mobilité et du développement local
\item À la fin de cette leçon, je sais expliquer les causes et les conséquences de l'exode rural
\item À la fin de cette leçon, je sais citer des initiatives de développement local au Cameroun et dans le monde hispanophone
\item À la fin de cette leçon, je sais conjuguer le prétérit indéfini et l'imparfait des verbes réguliers et irréguliers fréquents
\item À la fin de cette leçon, je sais choisir entre indéfini et imparfait selon la valeur de l'action (ponctuelle ou habituelle/descriptive)
\end{itemize}\end{multicols}}

\begin{sommaire}
\begin{multicols}{2}
\begin{enumerate}
\item Texte support et découverte
\item Compréhension du texte
\item Lexique thématique
\item Grammaire : prétérit indéfini et imparfait
\item Méthode du commentaire et commentaire modèle
\item Production guidée et entraînement à la traduction
\item Activité d'intégration
\item Méthodes et exercices résolus
\item Exercices
\item Corrigés des exercices
\item Fiche bilan
\end{enumerate}
\end{multicols}
\end{sommaire}

\begin{objectifs}
\begin{itemize}
\item À la fin de cette leçon, je sais lire et comprendre un texte sur le monde rural et le développement local
\item À la fin de cette leçon, je sais utiliser le lexique du champ, de l'agriculture, de la mobilité et du développement local
\item À la fin de cette leçon, je sais expliquer les causes et les conséquences de l'exode rural
\item À la fin de cette leçon, je sais citer des initiatives de développement local au Cameroun et dans le monde hispanophone
\item À la fin de cette leçon, je sais conjuguer le prétérit indéfini et l'imparfait des verbes réguliers et irréguliers fréquents
\item À la fin de cette leçon, je sais choisir entre indéfini et imparfait selon la valeur de l'action (ponctuelle ou habituelle/descriptive)
\item À la fin de cette leçon, je sais raconter au passé une expérience rurale vécue ou imaginée
\item À la fin de cette leçon, je sais traduire des phrases du français vers l'espagnol et inversement en respectant l'aspect verbal
\end{itemize}
\end{objectifs}

\begin{prerequis}
\begin{itemize}
\item Le présent de l'indicatif des verbes réguliers et irréguliers (Seconde)
\item Le lexique de la ville et des métiers (Seconde)
\item Les marqueurs temporels du passé : ayer, antes, cuando era niño (Seconde)
\item La structure de la phrase simple et les connecteurs (porque, pero, entonces)
\end{itemize}
\end{prerequis}

\newpage

\coursec{Texte support et découverte}

Chaque année, des milliers de jeunes quittent les villages de l'Ouest ou du Nord du Cameroun pour Douala et Yaoundé, à la recherche d'un emploi et d'une vie meilleure. Mais que deviennent les champs, les récoltes et les anciens restés au village ? En Espagne aussi, des régions entières se vident : on parle de la \esp{España vaciada}, « l'Espagne vidée ». Pourtant, des coopératives et des initiatives locales redonnent vie au monde rural. Comment raconter ces transformations au passé en espagnol ? Cette leçon nous donne les mots et les temps pour le faire.

\textbf{Problématique :} comment un texte espagnol simple raconte-t-il le départ des habitants d'un village et le retour de l'espoir grâce au développement local ? Quels mots et quels temps du passé utilise-t-il ?

\courssub{Lecture globale silencieuse du texte}

Avant toute explication, lisons le texte en silence, une première fois, sans dictionnaire. L'objectif n'est pas de tout comprendre, mais de saisir l'idée générale : de quoi parle-t-on ? Qui sont les personnages ? Que s'est-il passé ?

\begin{textebox}[El éxodo rural y la esperanza del campo]
Cuando yo era niño, mi abuelo trabajaba en el campo. Cada mañana se levantaba temprano y cuidaba sus cosechas de maíz y de café. Pero en 2010, mis tíos dejaron el pueblo y se fueron a la ciudad, porque allí había más trabajo y mejores carreteras. El pueblo se quedó casi vacío. Sin embargo, el año pasado los jóvenes crearon una cooperativa agrícola. Con el apoyo del municipio, construyeron una carretera nueva y vendieron sus productos en el mercado regional. Hoy, el desarrollo local devuelve la vida al campo.
\end{textebox}

\textbf{Glossaire (mots nouveaux) :}

\begin{center}
\begin{tabularx}{\linewidth}{|X|X|}
\hline
\rowcolor{popblueL} \textbf{Español} & \textbf{Français} \\
\hline
la cosecha & la récolte \\
\hline
se fueron (de \esp{irse}) & ils sont partis \\
\hline
vacío, -a & vide \\
\hline
el apoyo & le soutien \\
\hline
devuelve (de \esp{devolver}) & redonne, rend \\
\hline
se levantaba (de \esp{levantarse}) & se levait \\
\hline
el pueblo & le village \\
\hline
sin embargo & cependant, pourtant \\
\hline
\end{tabularx}
\end{center}

\textbf{Traduction du texte :} « Quand j'étais enfant, mon grand-père travaillait aux champs. Chaque matin, il se levait tôt et s'occupait de ses récoltes de maïs et de café. Mais en 2010, mes oncles ont quitté le village et sont partis à la ville, parce que là-bas il y avait plus de travail et de meilleures routes. Le village est resté presque vide. Cependant, l'année dernière, les jeunes ont créé une coopérative agricole. Avec le soutien de la commune, ils ont construit une route nouvelle et ont vendu leurs produits au marché régional. Aujourd'hui, le développement local redonne la vie à la campagne. »

\begin{methode}[Comment lire un texte espagnol]
\begin{enumerate}
\item Lire le \cle{titre} et la source : ils annoncent le thème.
\item Faire une \cle{première lecture globale}, sans dictionnaire, pour saisir l'idée générale.
\item Souligner les \cle{mots transparents} (mots proches du français : \esp{cooperativa}, \esp{regional}, \esp{municipio}).
\item Faire une \cle{deuxième lecture} avec le glossaire, phrase par phrase.
\item \cle{Reformuler} l'idée principale avec ses propres mots, d'abord en français, puis en espagnol.
\end{enumerate}
\end{methode}

\courssub{Lecture à voix haute : thème, type de texte, destinataire}

L'enseignant lit maintenant le texte à voix haute. Écoutons la prononciation : \esp{cosecha} se prononce « ko-sé-tcha », \esp{carretera} « ka-rré-té-ra » (avec un r roulé), \esp{jovenes} « rro-bé-nès » (le j se prononce comme un r grasseyé). Après l'écoute, répondons à trois questions de repérage.

\begin{exoresolu}[Repérage : thème, type de texte, destinataire]
\textbf{Énoncé.} Après l'écoute du texte, réponds aux questions : 1) De quel thème parle le texte ? 2) Quel est le type de ce texte (récit, description, lettre, article) ? 3) À qui s'adresse-t-il ?
\tcblower
\textbf{Solution détaillée.}
\begin{enumerate}
\item \textbf{Thème :} le texte parle de l'\esp{éxodo rural} (l'exode rural, le départ des villageois vers la ville) et du \esp{desarrollo local} (le développement local) qui redonne vie au village.
\item \textbf{Type de texte :} c'est un \cle{récit personnel} (un témoignage) : le narrateur dit \esp{yo} (« je ») et raconte des souvenirs de famille au passé, avec un début (\esp{cuando yo era niño}), un problème (\esp{mis tíos dejaron el pueblo}) et une fin heureuse (\esp{devuelve la vida al campo}).
\item \textbf{Destinataire :} le texte s'adresse à tout lecteur, en particulier aux jeunes, pour montrer que le retour au village et les initiatives locales sont possibles. C'est un texte à visée \cle{optimiste} et \cle{encourageante}.
\end{enumerate}
\end{exoresolu}

\courssub{Repérage des mots-clés}

Soulignons dans le texte les quatre mots-clés de la leçon. Ce sont eux qui portent le sens.

\begin{definition}[Les quatre mots-clés du texte]
\begin{itemize}
\item \esp{el campo} : la campagne, les champs. \esp{Mi abuelo trabajaba en el campo.} « Mon grand-père travaillait aux champs. »
\item \esp{el éxodo rural} : l'exode rural, c'est-à-dire le départ massif des habitants de la campagne vers la ville. \esp{Mis tíos se fueron a la ciudad.} « Mes oncles sont partis à la ville. »
\item \esp{la cooperativa} : la coopérative, une association où les membres travaillent ensemble et partagent les bénéfices. \esp{Los jóvenes crearon una cooperativa agrícola.} « Les jeunes ont créé une coopérative agricole. »
\item \esp{el desarrollo local} : le développement local, c'est-à-dire le progrès économique et social d'un village ou d'une commune, grâce à ses propres habitants. \esp{El desarrollo local devuelve la vida al campo.} « Le développement local redonne la vie à la campagne. »
\end{itemize}
\end{definition}

\begin{center}
\begin{tabularx}{\linewidth}{|X|X|X|}
\hline
\rowcolor{popblueL} \textbf{Français} & \textbf{Español} & \textbf{Remarque} \\
\hline
la campagne, les champs & el campo & mot masculin : \esp{el campo} \\
\hline
l'exode rural & el éxodo rural & accent sur le \esp{é} \\
\hline
la récolte & la cosecha & \esp{cosechar} = récolter \\
\hline
le/la paysan(ne), agriculteur & el agricultor / la agricultora & métier \\
\hline
la migration & la migración & \esp{migrar} = migrer \\
\hline
le transport & el transporte & mot transparent \\
\hline
la route & la carretera & ne pas confondre avec \esp{carrera} (carrière, course) \\
\hline
le développement local & el desarrollo local & \esp{desarrollar} = développer \\
\hline
la coopérative & la cooperativa & mot transparent \\
\hline
le marché & el mercado & \esp{el mercado regional} \\
\hline
\end{tabularx}
\end{center}

\begin{attention}[Faux amis et pièges de vocabulaire]
\begin{itemize}
\item \esp{el campo} signifie « la campagne, les champs », jamais « le camp » militaire (qui se dit \esp{el campamento}).
\item \esp{la carretera} signifie « la route » ; « la carrière » professionnelle se dit \esp{la carrera}. Ne traduis pas « carretera » par « charrette » !
\item \esp{el pueblo} signifie « le village » mais aussi « le peuple » selon le contexte : \esp{el pueblo español} = « le peuple espagnol ».
\item \esp{vacío} prend un accent sur le \esp{í} : \esp{vacío}, « vide ». Sans accent, c'est une faute.
\end{itemize}
\end{attention}

\begin{popfigure}
\begin{tikzpicture}[mindmap, grow cyclic, every node/.style={concept}, concept color=popgreen, text=white,
root concept/.append style={font=\bfseries, minimum size=3.2cm},
level 1/.append style={level distance=3.6cm, sibling angle=90},
level 1 concept/.append style={concept color=popblue, minimum size=2.4cm, font=\small\bfseries}]
\node[root concept]{El mundo rural}
  child { node[align=center]{agricultura\\ \footnotesize cosecha, agricultor} }
  child { node[align=center]{éxodo rural\\ \footnotesize migración, ciudad} }
  child { node[align=center]{movilidad\\ \footnotesize transporte, carretera} }
  child { node[align=center]{desarrollo local\\ \footnotesize cooperativa, mercado} };
\end{tikzpicture}
\legende{Carte mentale du vocabulaire : autour du monde rural (\esp{el mundo rural}), quatre grandes idées : l'agriculture, l'exode rural, la mobilité et le développement local.}
\end{popfigure}

\begin{savaistu}[La España vaciada et les GIC camerounais]
En Espagne, on parle de la \esp{España vaciada} (« l'Espagne vidée ») pour désigner les régions et les villages abandonnés par leurs habitants partis vers Madrid, Barcelone ou les côtes. Certains villages comptent moins de cent habitants, presque tous âgés. Au Cameroun, le même phénomène existe : les jeunes quittent les villages pour Douala et Yaoundé. Mais il existe aussi des solutions locales : les \cle{GIC} (Groupes d'Initiative Commune), où les villageois mettent leur argent et leur travail en commun pour produire et vendre ensemble (cacao, maïs, huile de palme), ressemblent beaucoup aux \esp{cooperativas} espagnoles. Le monde rural n'est pas condamné : il peut renaître grâce à l'organisation collective.
\end{savaistu}

\courssub{Identification des temps du passé dans le texte}

Le texte raconte des souvenirs et des événements passés. En espagnol, deux temps se partagent ce travail : l'\cle{imparfait} (\esp{pretérito imperfecto}) et le \cle{prétérit indéfini} (\esp{pretérito indefinido}). Observons d'abord, la règle détaillée viendra dans la section grammaire.

Reprenons les verbes du texte et classons-les :

\begin{center}
\begin{tabularx}{\linewidth}{|X|X|X|}
\hline
\rowcolor{popblueL} \textbf{Verbe du texte} & \textbf{Temps} & \textbf{Fonction} \\
\hline
\esp{era} (j'étais) & imparfait & situation, contexte \\
\hline
\esp{trabajaba} (travaillait) & imparfait & habitude passée \\
\hline
\esp{se levantaba} (se levait) & imparfait & habitude (\esp{cada mañana}) \\
\hline
\esp{cuidaba} (s'occupait de) & imparfait & habitude passée \\
\hline
\esp{había} (il y avait) & imparfait & description \\
\hline
\esp{dejaron} (ont quitté) & indéfini & action ponctuelle, datée (\esp{en 2010}) \\
\hline
\esp{se fueron} (sont partis) & indéfini & action ponctuelle \\
\hline
\esp{se quedó} (est resté) & indéfini & résultat, action unique \\
\hline
\esp{crearon} (ont créé) & indéfini & action ponctuelle (\esp{el año pasado}) \\
\hline
\esp{construyeron} (ont construit) & indéfini & action achevée \\
\hline
\esp{vendieron} (ont vendu) & indéfini & action achevée \\
\hline
\esp{devuelve} (redonne) & présent & situation actuelle \\
\hline
\end{tabularx}
\end{center}

\begin{aretenir}[Première observation : imparfait ou indéfini ?]
Dans un récit au passé :
\begin{itemize}
\item L'\cle{imparfait} peint le \textbf{décor} et les \textbf{habitudes} : \esp{Cuando yo era niño, mi abuelo trabajaba en el campo.} « Quand j'étais enfant, mon grand-père travaillait aux champs. »
\item Le \cle{prétérit indéfini} raconte les \textbf{actions ponctuelles} et datées qui font avancer l'histoire : \esp{En 2010, mis tíos dejaron el pueblo.} « En 2010, mes oncles ont quitté le village. »
\end{itemize}
Les mots-indices aident à choisir : \esp{cada mañana}, \esp{siempre}, \esp{cuando era niño} appellent l'imparfait ; \esp{en 2010}, \esp{el año pasado}, \esp{ayer} appellent l'indéfini.
\end{aretenir}

\begin{attention}[Erreur fréquente des francophones]
En français, le passé composé et l'imparfait ne coïncident pas toujours avec l'indéfini et l'imparfait espagnols. « Il est parti » se traduit \esp{se fue} (indéfini) et non \esp{se iba}. À l'inverse, « chaque matin, il se levait tôt » exige l'imparfait : \esp{cada mañana se levantaba temprano}, jamais \esp{se levantó} avec \esp{cada mañana}. Repère toujours le mot-indice avant de conjuguer.
\end{attention}

\courssub{Reformulation orale du texte en trois phrases}

Dernière étape de la découverte : chaque élève reformule le texte à l'oral, en espagnol, en seulement trois phrases. C'est un excellent entraînement à l'expression orale et à la synthèse.

\begin{exemplebox}[Modèle de reformulation en trois phrases]
\begin{enumerate}
\item \esp{Antes, el abuelo del narrador trabajaba en el campo y cuidaba sus cosechas.} « Avant, le grand-père du narrateur travaillait aux champs et s'occupait de ses récoltes. »
\item \esp{En 2010, sus tíos se fueron a la ciudad y el pueblo se quedó casi vacío.} « En 2010, ses oncles sont partis à la ville et le village est resté presque vide. »
\item \esp{El año pasado, los jóvenes crearon una cooperativa y el desarrollo local devuelve la vida al campo.} « L'année dernière, les jeunes ont créé une coopérative et le développement local redonne la vie à la campagne. »
\end{enumerate}
Remarque la logique : phrase 1 = le passé heureux (imparfait), phrase 2 = la crise (indéfini), phrase 3 = la solution (indéfini puis présent).
\end{exemplebox}

\begin{experience}[À toi de parler : mi pueblo, mi historia]
Par groupes de trois, imaginez le village de l'un de vos grands-parents (par exemple près de Bafoussam, Garoua ou Kribi). Chaque élève prononce trois phrases au passé, sur le modèle ci-dessus :
\begin{enumerate}
\item une habitude passée à l'imparfait : \esp{Mi abuela cultivaba...}
\item un départ ou un changement à l'indéfini : \esp{Mis primos se fueron a...}
\item une initiative récente : \esp{El año pasado, los jóvenes crearon...}
\end{enumerate}
Le groupe choisit ensuite le meilleur récit et le présente à la classe.
\end{experience}

\begin{exercice}[Vérification de la découverte]
\begin{enumerate}
\item Relis le texte et entoure trois verbes à l'imparfait et trois verbes au prétérit indéfini.
\item Traduis en français : \esp{Con el apoyo del municipio, construyeron una carretera nueva.}
\item Complète avec le mot-clé qui convient : \esp{campo}, \esp{cooperativa}, \esp{carretera}, \esp{cosecha} : a) Los jóvenes crearon una \ldots{} agrícola. b) El abuelo trabajaba en el \ldots{}. c) Construyeron una \ldots{} nueva. d) Cuidaba sus \ldots{} de maíz.
\item Reformule oralement le texte en trois phrases, sans regarder le modèle.
\end{enumerate}
\end{exercice}

\begin{corrige}[Exercice : Vérification de la découverte]
\begin{enumerate}
\item Imparfait : \esp{era}, \esp{trabajaba}, \esp{se levantaba} (aussi \esp{cuidaba}, \esp{había}). Indéfini : \esp{dejaron}, \esp{se fueron}, \esp{crearon} (aussi \esp{se quedó}, \esp{construyeron}, \esp{vendieron}).
\item « Avec le soutien de la commune, ils ont construit une route nouvelle. »
\item a) \esp{cooperativa} ; b) \esp{campo} ; c) \esp{carretera} ; d) \esp{cosechas}.
\item Réponse libre, sur le modèle : passé heureux (imparfait) / départ (indéfini) / initiative récente (indéfini et présent).
\end{enumerate}
\end{corrige}

\begin{aretenir}[Bilan de la découverte]
\begin{itemize}
\item Le texte est un \cle{récit personnel} optimiste sur l'exode rural et le développement local.
\item Les quatre mots-clés : \esp{el campo}, \esp{el éxodo rural}, \esp{la cooperativa}, \esp{el desarrollo local}.
\item Deux temps racontent le passé : l'\cle{imparfait} (décor et habitudes) et le \cle{prétérit indéfini} (actions ponctuelles et datées).
\item Méthode de lecture : titre, lecture globale, mots transparents, glossaire, reformulation.
\end{itemize}
\end{aretenir}

\coursec{Compréhension du texte}

Dans la section précédente, nous avons découvert le texte support \esp{« El pueblo que volvió a vivir »} : nous l'avons lu, écouté, et nous avons repéré les mots nouveaux. Mais lire un texte, ce n'est pas seulement le prononcer : c'est le \cle{comprendre}. Cette section vous apprend à répondre aux questions de compréhension comme au baccalauréat : d'abord le sens global, puis les détails, puis l'interprétation. Pour travailler, relisons le texte intégralement.

\begin{textebox}[El pueblo que volvió a vivir (texte support)]
San Miguel es un pueblo pequeño del oeste de Camerún. Hace treinta años, era un lugar lleno de vida: los agricultores cultivaban cacao y café, los niños jugaban en las calles y el mercado estaba lleno todos los sábados. Pero poco a poco, todo cambió. Las carreteras eran muy malas y los camiones no podían llegar hasta el pueblo para comprar la cosecha. No había ni hospital ni liceo. Entonces muchos jóvenes decidieron irse a la ciudad.

Mis tíos, por ejemplo, dejaron el pueblo en 1995 porque en Douala había más trabajo. Otros vecinos emigraron a Yaoundé o incluso al extranjero. El pueblo se quedó casi vacío: solo vivían allí los abuelos y algunos niños. Las casas se cerraron y los campos se abandonaron.

Pero hace cinco años, todo empezó a cambiar otra vez. Un grupo de jóvenes del pueblo, que habían estudiado en la ciudad, volvieron con una idea: crear una cooperativa para vender el cacao directamente, sin intermediarios. El ayuntamiento les ayudó con dinero y construyó una carretera nueva. Ahora los camiones llegan cada semana, la cooperativa vende el cacao en el mercado regional y los jóvenes ganan buen dinero. Algunos emigrantes incluso han vuelto al pueblo. Hoy San Miguel vive otra vez: hay una escuela nueva, un centro de salud y mucha esperanza.
\end{textebox}

\textbf{Glossaire :}\par
\esp{el pueblo} = le village ; \esp{la cosecha} = la récolte ; \esp{irse} = s'en aller, partir ; \esp{el extranjero} = l'étranger ; \esp{el liceo} = le lycée ; \esp{se quedó vacío} = est devenu vide ; \esp{el ayuntamiento} = la mairie ; \esp{sin intermediarios} = sans intermédiaires ; \esp{el centro de salud} = le centre de santé ; \esp{la esperanza} = l'espoir.

\courssub{Les questions de compréhension globale}

Avant de répondre, apprenons la bonne méthode : elle vaut pour toutes les questions de cette leçon et pour l'examen.

\begin{methode}[Répondre à une question de compréhension]
\begin{enumerate}
\item \textbf{Lire la question deux fois} et souligner le mot interrogatif : \esp{¿quién?} (qui), \esp{¿dónde?} (où), \esp{¿cuándo?} (quand), \esp{¿por qué?} (pourquoi), \esp{¿qué?} (que, quel).
\item \textbf{Repérer dans le texte} la phrase ou le passage qui contient la réponse.
\item \textbf{Répondre par une phrase complète en espagnol}, jamais par un seul mot. Reprendre les mots de la question : à \esp{¿Dónde está San Miguel?} on répond \esp{San Miguel está en el oeste de Camerún.}
\item \textbf{Justifier par une courte citation} du texte entre guillemets, introduite par \esp{como dice el texto} ou \esp{según el texto}.
\item \textbf{Relire sa réponse} : verbe conjugué, accents, ponctuation espagnole (¿...?).
\end{enumerate}
\end{methode}

Voici maintenant les cinq questions globales, avec leurs réponses modèles. Observez comment chaque réponse reprend les mots de la question.

\begin{exemplebox}[Questions globales et réponses modèles]
\textbf{1. ¿Quiénes son los protagonistas del texto?} (Qui sont les protagonistes ?)\par
\esp{Los protagonistas son los habitantes de San Miguel, sobre todo los jóvenes que volvieron al pueblo.} = Les protagonistes sont les habitants de San Miguel, surtout les jeunes qui sont revenus au village.

\textbf{2. ¿Dónde y cuándo ocurre la historia?} (Où et quand se passe l'histoire ?)\par
\esp{La historia ocurre en San Miguel, un pueblo del oeste de Camerún; empieza hace treinta años y termina hoy.} = L'histoire se passe à San Miguel, un village de l'ouest du Cameroun ; elle commence il y a trente ans et finit aujourd'hui.

\textbf{3. ¿Cuál es el problema del pueblo?} (Quel est le problème du village ?)\par
\esp{El problema es el éxodo rural: «muchos jóvenes decidieron irse a la ciudad» y «el pueblo se quedó casi vacío».} = Le problème est l'exode rural : « beaucoup de jeunes ont décidé de partir à la ville » et « le village est devenu presque vide ».

\textbf{4. ¿Cuál es la solución?} (Quelle est la solution ?)\par
\esp{La solución es crear una cooperativa con la ayuda del ayuntamiento: «crear una cooperativa para vender el cacao directamente».} = La solution est de créer une coopérative avec l'aide de la mairie.

\textbf{5. ¿Cuál es el resultado final?} (Quel est le résultat final ?)\par
\esp{El resultado es positivo: «Hoy San Miguel vive otra vez».} = Le résultat est positif : « Aujourd'hui San Miguel revit ».
\end{exemplebox}

\begin{attention}[Ne traduis pas mot à mot !]
Le piège classique du francophone est la traduction littérale. \esp{Había trabajo} signifie « il y avait du travail », et non « avait travail ». De même, \esp{los camiones no podían llegar hasta el pueblo} = « les camions ne pouvaient pas arriver jusqu'au village » (et non « les camions ne pouvaient pas venir jusqu'à »). Comprendre, c'est saisir le \cle{sens global} d'une phrase, pas aligner les mots un par un. Autre exemple : \esp{el pueblo se quedó casi vacío} = « le village est devenu presque vide » (le verbe \esp{quedarse} ne se traduit pas ici par « rester » mais par « devenir », « rester » au sens de « rester vide »).
\end{attention}

\courssub{Les questions de détail}

Les questions de détail portent sur un passage précis. Il faut relire le paragraphe concerné avant de répondre.

\begin{exoresolu}[Questions de détail]
\textbf{a) ¿Por qué se fueron los tíos del narrador?} (Pourquoi les oncles du narrateur sont-ils partis ?)\par
\textbf{b) ¿Qué hicieron los jóvenes que volvieron al pueblo?} (Qu'ont fait les jeunes qui sont revenus ?)
\tcblower
\textbf{a)} On repère le deuxième paragraphe. Réponse complète : \esp{Los tíos del narrador se fueron a Douala en 1995 porque allí había más trabajo.} = Les oncles du narrateur sont partis à Douala en 1995 parce que là-bas il y avait plus de travail. La cause est introduite par \esp{porque}, comme en français « parce que ».\par
\textbf{b)} On repère le troisième paragraphe. Réponse : \esp{Los jóvenes volvieron con una idea: crearon una cooperativa para vender el cacao directamente, sin intermediarios.} = Les jeunes sont revenus avec une idée : ils ont créé une coopérative pour vendre le cacao directement, sans intermédiaires. Notez les deux temps du passé : \esp{volvieron} (action ponctuelle, prétérit indéfini) et la structure \esp{para + infinitivo} pour exprimer le but (« pour vendre »).
\end{exoresolu}

\begin{exemplebox}[Exemple traduit à retenir]
\esp{Mis tíos dejaron el pueblo porque allí había más trabajo.} = « Mes oncles ont quitté le village parce que là-bas il y avait plus de travail. »\par
Observez : \esp{dejaron} = prétérit indéfini (action terminée, datée : 1995) ; \esp{había} = imparfait (situation, décor : le travail existait là-bas de façon durable). Nous étudierons ces deux temps en détail dans la section grammaire.
\end{exemplebox}

\courssub{Les questions d'interprétation}

L'interprétation va plus loin que la lecture littérale : on cherche ce que le texte \emph{suggère}. Il faut alors donner son avis avec \esp{creo que}, \esp{pienso que}, \esp{en mi opinión}, tout en s'appuyant sur le texte.

\begin{exemplebox}[Interpréter : la « carretera nueva »]
\textbf{¿Qué simboliza la carretera nueva en el texto?} (Que symbolise la route neuve ?)\par
Réponse modèle : \esp{En mi opinión, la carretera nueva simboliza la esperanza y el desarrollo: gracias a ella, «los camiones llegan cada semana» y el pueblo vuelve a vivir. La carretera une el pueblo con el mundo exterior.} = À mon avis, la route neuve symbolise l'espoir et le développement : grâce à elle, « les camions arrivent chaque semaine » et le village revit. La route relie le village au monde extérieur.\par
Autres symboles possibles : \esp{las casas cerradas} (les maisons fermées) symbolisent l'abandon et la tristesse ; \esp{la escuela nueva} symbolise l'avenir et la jeunesse.
\end{exemplebox}

\courssub{Vrai ou faux justifié}

Au baccalauréat, un « vrai » ou « faux » sans justification ne rapporte aucun point. La justification est une \cle{citation exacte} du texte.

\begin{exercice}[¿Verdadero o falso? Justifica con una cita del texto]
\begin{enumerate}
\item \esp{San Miguel está en el este de Camerún.}
\item \esp{Los jóvenes se fueron porque las carreteras eran malas y no había trabajo.}
\item \esp{El ayuntamiento no ayudó a los jóvenes.}
\item \esp{Algunos emigrantes volvieron al pueblo.}
\end{enumerate}
\end{exercice}

\begin{corrige}[Exercice : verdadero o falso]
\begin{enumerate}
\item \textbf{Falso.} El texto dice: \esp{«un pueblo pequeño del oeste de Camerún»} (ouest, pas est).
\item \textbf{Verdadero.} \esp{«Las carreteras eran muy malas... No había ni hospital ni liceo. Entonces muchos jóvenes decidieron irse»} et \esp{«en Douala había más trabajo»}.
\item \textbf{Falso.} \esp{«El ayuntamiento les ayudó con dinero y construyó una carretera nueva.»}
\item \textbf{Verdadero.} \esp{«Algunos emigrantes incluso han vuelto al pueblo.»}
\end{enumerate}
\end{corrige}

\courssub{Causes, conséquences et solutions : la logique du texte}

Comprendre un texte de ce type, c'est savoir reconstruire sa \cle{chaîne logique} : les causes de l'exode rural, ses conséquences, puis les solutions. Voici le schéma complet du texte.

\begin{popfigure}
\begin{center}
\begin{tikzpicture}[
  font=\small,
  causa/.style={rectangle, rounded corners, draw=poporange, fill=poporangeL, text width=3.4cm, align=center, inner sep=4pt},
  consec/.style={rectangle, rounded corners, draw=poppink, fill=poppinkL, text width=3.4cm, align=center, inner sep=4pt},
  sol/.style={rectangle, rounded corners, draw=popgreen, fill=popgreenL, text width=3.4cm, align=center, inner sep=4pt},
  centro/.style={rectangle, rounded corners, draw=popblue, fill=popblueL, text width=3.2cm, align=center, inner sep=5pt, font=\small\bfseries},
  flecha/.style={-{Stealth[length=2.5mm]}, thick, popdark}
]
\node[centro] (exodo) at (0,0) {Éxodo rural};
\node[causa] (c1) at (-4.6,1.8) {Falta de trabajo en el pueblo};
\node[causa] (c2) at (-4.6,0) {Malas carreteras};
\node[causa] (c3) at (-4.6,-1.8) {Sin hospital ni liceo};
\node[consec] (k1) at (4.6,1.4) {Pueblo casi vacío};
\node[consec] (k2) at (4.6,0) {Solo abuelos y niños (envejecimiento)};
\node[consec] (k3) at (4.6,-1.4) {Campos abandonados};
\node[sol] (s1) at (0,-3.2) {Cooperativa de cacao};
\node[sol] (s2) at (4.6,-3.2) {Ayuntamiento: dinero y carretera nueva};
\node[sol] (s3) at (-4.6,-3.2) {Mercado regional};
\draw[flecha] (c1) -- (exodo);
\draw[flecha] (c2) -- (exodo);
\draw[flecha] (c3) -- (exodo);
\draw[flecha] (exodo) -- (k1);
\draw[flecha] (exodo) -- (k2);
\draw[flecha] (exodo) -- (k3);
\draw[flecha, popgreen] (s1) -- (0,-1.1);
\draw[flecha, popgreen] (s2) -- (1.6,-1.6);
\draw[flecha, popgreen] (s3) -- (-1.6,-1.6);
\end{tikzpicture}
\end{center}
\legende{Chaîne logique du texte : causes (orange), conséquences (rose) et solutions (vert) de l'exode rural.}
\end{popfigure}

Pour parler des causes et des conséquences en espagnol, il faut maîtriser les \cle{connecteurs logiques}. Les voici, comparés au français.

\begin{center}
\begin{tabularx}{\linewidth}{|X|X|X|}
\hline
\rowcolor{popblueL} Français & Español & Exemple \\
\hline
parce que & porque & \esp{Se fueron porque no había trabajo.} \\
\hline
à cause de & a causa de / por & \esp{El pueblo se vació por las malas carreteras.} \\
\hline
donc, alors & entonces, así que & \esp{No había liceo, así que se fueron.} \\
\hline
c'est pourquoi & por eso & \esp{Las carreteras eran malas; por eso los camiones no llegaban.} \\
\hline
grâce à & gracias a & \esp{Gracias a la cooperativa, el pueblo vive otra vez.} \\
\hline
pour + infinitif (but) & para + infinitivo & \esp{Crearon una cooperativa para vender el cacao.} \\
\hline
\end{tabularx}
\end{center}

\begin{attention}[porque / por eso / para : ne pas confondre]
Les francophones mélangent souvent ces trois mots avec \esp{por}. Retenez : \esp{porque} = parce que (cause, répond à \esp{¿por qué?}) ; \esp{por eso} = c'est pourquoi (conséquence) ; \esp{para + infinitivo} = pour + infinitif (but). Jamais \esp{porque de} ni \esp{para de} : \esp{Vine para ayudar} = « je suis venu pour aider ».
\end{attention}

\begin{exoresolu}[Relier cause et conséquence]
Transformez les deux phrases en une seule avec le connecteur indiqué.\par
\textbf{a)} \esp{No había trabajo en el pueblo. Los jóvenes se fueron a la ciudad.} (\esp{porque})\par
\textbf{b)} \esp{Las carreteras eran muy malas. Los camiones no podían llegar.} (\esp{por eso})
\tcblower
\textbf{a)} \esp{Los jóvenes se fueron a la ciudad porque no había trabajo en el pueblo.} = Les jeunes sont partis à la ville parce qu'il n'y avait pas de travail au village. (La cause suit \esp{porque}.)\par
\textbf{b)} \esp{Las carreteras eran muy malas; por eso los camiones no podían llegar.} = Les routes étaient très mauvaises ; c'est pourquoi les camions ne pouvaient pas arriver. (La conséquence suit \esp{por eso}.) Attention à l'ordre : on ne dit pas \esp{por eso las carreteras eran malas}, ce qui inverserait la logique.
\end{exoresolu}

\courssub{Débat bref : ¿campo o ciudad?}

Pour terminer la compréhension, passons à l'oral : un mini-débat permet de réutiliser le lexique et les connecteurs dans une situation de communication réelle.

\begin{experience}[Débat : ¿Es mejor vivir en el campo o en la ciudad?]
La classe se divise en deux groupes. Le groupe A défend la vie à la campagne, le groupe B la vie en ville. Chaque élève prépare deux arguments avec les expressions :\par
\esp{Pienso que vivir en el campo es mejor porque...} ; \esp{En mi opinión, la ciudad es mejor porque...} ; \esp{Estoy de acuerdo / No estoy de acuerdo porque...}\par
Arguments possibles pour le campo : \esp{el aire es puro, la vida es tranquila, hay comunidad}. Pour la ciudad : \esp{hay más trabajo, hay hospitales y liceos, hay transporte}. Le professeur conclut en rappelant la leçon du texte : avec le développement local (coopérative, routes, services), la campagne peut offrir un bel avenir.
\end{experience}

\begin{savaistu}[Le monde rural au Cameroun et en Guinée équatoriale]
Le Cameroun est souvent appelé « l'Afrique en miniature » pour sa diversité géographique. La région de l'Ouest (Bamileke, Bamoun) est une zone de \emph{montagne} à forte densité rurale, productrice de café, cacao, maïs, haricots. L'exode rural y est massif vers Douala et Yaoundé, mais des initiatives de \esp{desarrollo local} émergent : coopératives agricoles (ex. \esp{cooperativas de caficultores} dans le Noun), pistes rurales réhabilitées, \esp{telecentros} communautaires. En Guinée équatoriale (pays hispanophone voisin), l'île de Bioko et la région continentale de Río Muni connaissent aussi une transformation : l'argent du pétrole a financé des routes (\esp{carreteras}) et des infrastructures, mais l'agriculture traditionnelle (cacao, café, bois) peine à retenir la jeunesse. Le défi commun : fixer les populations par l'emploi local et les services (santé, éducation), comme le montre le texte de San Miguel.
\end{savaistu}

\begin{aretenir}[Ce qu'il faut retenir de cette section]
\begin{itemize}
\item Répondre toujours par une \textbf{phrase complète en espagnol}, en reprenant les mots de la question.
\item Justifier par une \textbf{courte citation} entre guillemets : c'est elle qui rapporte des points.
\item Distinguer les trois niveaux : compréhension globale (qui, où, quand), détail (relire le paragraphe), interprétation (\esp{en mi opinión} + appui sur le texte).
\item Connaître la chaîne logique du texte : causes (\esp{porque}) $\rightarrow$ conséquences (\esp{por eso}) $\rightarrow$ solutions (\esp{gracias a}, \esp{para}).
\item Ne jamais traduire mot à mot : \esp{había trabajo} = « il y avait du travail » ; \esp{se quedó vacío} = « est devenu vide ».
\end{itemize}
\end{aretenir}

La compréhension est maintenant assurée. Dans la section suivante, nous étudierons en détail le \textbf{lexique thématique} du monde rural et du développement local, pour enrichir votre vocabulaire avant la grammaire et la production écrite.

\coursec{Lexique thématique}

Dans la section précédente, nous avons lu et compris le texte \esp{Un joven entre el campo y la ciudad} : nous avons identifié les causes de l'\esp{éxodo rural} et les initiatives de \esp{desarrollo local} qui ont permis au village de renaître. Mais pour parler librement du monde rural, il ne suffit pas de comprendre un texte : il faut disposer d'un \cle{vocabulaire} solide et bien organisé. C'est l'objectif de cette section : construire, champ par champ, tout le lexique nécessaire pour décrire la campagne, expliquer la mobilité des populations et raconter les projets de développement local au Cameroun et dans le monde hispanophone.

\courssub{Pourquoi apprendre le vocabulaire par champs lexicaux ?}

\begin{definition}[Champ lexical]
Un \cle{champ lexical} est l'ensemble des mots qui se rapportent à une même idée, à un même thème. Par exemple, autour du thème de la campagne, on trouve \esp{el campo}, \esp{el agricultor}, \esp{la cosecha}, \esp{la tierra}, \esp{el ganado}. Apprendre les mots par familles thématiques permet de les mémoriser plus facilement et de les réactiver rapidement à l'oral comme à l'écrit.
\end{definition}

Avant d'énoncer les listes, observons d'abord les mots en situation. Lis ce court texte original et souligne tous les mots qui appartiennent au thème du monde rural et de la mobilité.

\begin{textebox}[Un joven entre el campo y la ciudad]
Me llamo Samuel y nací en un pueblo pequeño cerca de Bafoussam, en el oeste de Camerún. Mi padre es agricultor: cultiva maíz y café, y cuida también un poco de ganado. Cuando era niño, yo ayudaba a mi familia en los campos durante la cosecha. La vida era tranquila, pero el pueblo no tenía ni hospital ni carretera asfaltada, y el camino al mercado era muy largo.

A los dieciocho años, como muchos jóvenes de mi región, dejé el pueblo y me fui a Duala para buscar trabajo. Este movimiento se llama éxodo rural: miles de personas abandonan el campo cada año para vivir en la ciudad. En Duala encontré un empleo, pero la vida era cara y difícil.

Sin embargo, hace dos años, todo cambió. El municipio creó una cooperativa agrícola con el apoyo de una ONG. El proyecto ayuda a los jóvenes a cultivar, a vender sus productos en el mercado y a construir una nueva carretera. Gracias a esta iniciativa, volví al campo. Hoy, con otros jóvenes, doy vida a mi pueblo. El desarrollo local no es un sueño: es una realidad.
\end{textebox}

\textbf{Glossaire}\par
\esp{nací} = je suis né ; \esp{cuida} = il s'occupe de ; \esp{asfaltada} = goudronnée ; \esp{buscar trabajo} = chercher du travail ; \esp{un empleo} = un emploi ; \esp{hace dos años} = il y a deux ans ; \esp{una ONG} = une ONG ; \esp{un sueño} = un rêve.

\textbf{Questions de repérage lexical}\par
\begin{enumerate}
\item Trouve dans le texte trois mots du champ de l'agriculture.
\item Trouve deux expressions qui expriment le départ du village.
\item Quel mot désigne le mouvement de population du campo vers la ciudad ?
\item Releve deux mots liés au développement local.
\end{enumerate}

\courssub{Champ 1 : El mundo rural y la agricultura}

\begin{definition}[Vocabulaire de la campagne et de l'agriculture]
\begin{itemize}
\item \esp{el campo} = la campagne (par opposition à la ville) ; aussi : le champ (terrain cultivé)
\item \esp{el agricultor / la agricultora} = l'agriculteur / l'agricultrice
\item \esp{la tierra} = la terre (le sol que l'on cultive)
\item \esp{cultivar} = cultiver
\item \esp{la cosecha} = la récolte
\item \esp{el pueblo} = le village
\item \esp{el ganado} = le bétail
\item \esp{el campesino / la campesina} = le paysan / la paysanne
\end{itemize}
\end{definition}

\begin{exemplebox}[Les mots en contexte]
\esp{El agricultor cultiva la tierra con sus hijos.} = L'agriculteur cultive la terre avec ses enfants.

\esp{El agricultor vende su cosecha en el mercado.} = L'agriculteur vend sa récolte au marché.

\esp{En mi pueblo, casi todas las familias tienen ganado.} = Dans mon village, presque toutes les familles ont du bétail.

\esp{La cosecha de cacao fue muy buena este año.} = La récolte de cacao a été très bonne cette année.
\end{exemplebox}

\begin{attention}[Faux ami : « el campo » n'est pas « le camp »]
\esp{El campo} signifie \cle{la campagne} ou \cle{le champ}, mais jamais « le camp » (militaire, de vacances ou de réfugiés). Pour « le camp », l'espagnol dit \esp{el campamento}.

\esp{Vivo en el campo.} = J'habite à la campagne. (et non « dans un camp »)

\esp{Los refugiados viven en un campamento.} = Les réfugiés vivent dans un camp.

Autre piège : \esp{el pueblo} signifie « le village » mais aussi « le peuple ». Le contexte décide : \esp{mi pueblo está cerca de Bafoussam} = mon village est près de Bafoussam ; \esp{el pueblo camerunés} = le peuple camerounais.
\end{attention}

\courssub{Champ 2 : La movilidad}

\begin{definition}[Éxodo rural]
L'\esp{éxodo rural} (l'exode rural) est le \esp{movimiento de población que abandona el campo para vivir en la ciudad}, c'est-à-dire le départ massif et durable des habitants de la campagne vers la ville, généralement pour chercher du travail, des études ou de meilleurs services.
\end{definition}

\begin{definition}[Vocabulaire de la mobilité]
\begin{itemize}
\item \esp{el éxodo rural} = l'exode rural
\item \esp{la migración} = la migration
\item \esp{emigrar} = émigrer
\item \esp{mudarse} = déménager, changer de domicile
\item \esp{dejar el pueblo} = quitter le village
\item \esp{irse a la ciudad} = partir à la ville (verbe pronominal irrégulier : \esp{me voy, te vas, se va...})
\item \esp{volver al campo} = retourner à la campagne
\item \esp{buscar trabajo} = chercher du travail
\end{itemize}
\end{definition}

\begin{exemplebox}[Exemples traduits]
\esp{Muchos jóvenes dejan el pueblo porque no hay trabajo.} = Beaucoup de jeunes quittent le village parce qu'il n'y a pas de travail.

\esp{Mi hermano se mudó a Yaundé el año pasado.} = Mon frère a déménagé à Yaoundé l'année dernière.

\esp{Después de diez años en la ciudad, volvió al campo.} = Après dix ans en ville, il est retourné à la campagne.
\end{exemplebox}

\begin{attention}[« Mudarse » ou « moverse » ?]
Les francophones confondent souvent ces deux verbes pronominaux :

\esp{mudarse} = \cle{déménager}, changer de maison ou de ville : \esp{Nos mudamos a Duala.} = Nous avons déménagé à Douala.

\esp{moverse} = \cle{bouger} (le corps, se déplacer dans l'espace) : \esp{No te muevas.} = Ne bouge pas.

On ne dit donc jamais \esp{me moví a la ciudad} pour « j'ai déménagé en ville » : il faut \esp{me mudé a la ciudad}.
\end{attention}

\courssub{Champ 3 : Las infraestructuras}

\begin{definition}[Vocabulaire des infrastructures et des services]
\begin{itemize}
\item \esp{el transporte} = le transport
\item \esp{la carretera} = la route (goudronnée, entre les villes)
\item \esp{el camino} = le chemin, la piste
\item \esp{el mercado} = le marché
\item \esp{los servicios} = les services (publics)
\item \esp{la escuela} = l'école
\item \esp{el hospital} = l'hôpital (attention : le \esp{h} est muet : « os-pi-tal »)
\item \esp{el agua potable} = l'eau potable
\item \esp{la electricidad} = l'électricité
\end{itemize}
\end{definition}

L'absence d'infrastructures est l'une des causes principales de l'exode rural ; leur construction est au cœur du développement local. Compare :

\esp{En el pueblo no hay hospital ni carretera asfaltada.} = Au village, il n'y a ni hôpital ni route goudronnée.

\esp{El camino al mercado es muy largo en la estación de las lluvias.} = La piste vers le marché est très longue à la saison des pluies.

\courssub{Champ 4 : El desarrollo local}

\begin{definition}[Vocabulaire du développement local]
\begin{itemize}
\item \esp{el desarrollo local} = le développement local
\item \esp{la cooperativa} = la coopérative
\item \esp{el apoyo} = le soutien, l'appui
\item \esp{apoyar} = soutenir
\item \esp{el municipio} = la commune, la municipalité
\item \esp{crear} = créer
\item \esp{vender} = vendre
\item \esp{el proyecto} = le projet
\item \esp{dar vida al pueblo} = redonner vie au village
\end{itemize}
\end{definition}

\begin{exemplebox}[Exemples traduits]
\esp{La cooperativa apoya a los jóvenes.} = La coopérative soutient les jeunes.

\esp{El municipio creó un proyecto para construir una carretera.} = La commune a créé un projet pour construire une route.

\esp{Gracias a la cooperativa, los agricultores venden su café a mejor precio.} = Grâce à la coopérative, les agriculteurs vendent leur café à meilleur prix.

\esp{Los jóvenes dan vida al pueblo.} = Les jeunes redonnent vie au village.
\end{exemplebox}

\begin{propriete}[Deux connecteurs indispensables : « gracias a » et « sin embargo »]
\esp{Gracias a} + nom = \cle{grâce à} (cause positive) : \esp{Gracias al proyecto, el pueblo tiene electricidad.} = Grâce au projet, le village a l'électricité.

\esp{Sin embargo} = \cle{cependant, pourtant} (opposition) ; il se place en début de phrase, suivi d'une virgule : \esp{La vida en la ciudad es cara. Sin embargo, muchos jóvenes emigran.} = La vie en ville est chère. Cependant, beaucoup de jeunes émigrent.
\end{propriete}

\begin{savaistu}[« Cooperativa », un mot transparent et une réalité camerounaise]
\esp{Cooperativa} est un \cle{mot transparent} : il se comprend immédiatement grâce au français « coopérative », car les deux mots viennent du latin \emph{cooperari} (travailler ensemble). Au Cameroun, les coopératives de café-cacao du Littoral et de l'Ouest fonctionnent exactement selon le même modèle qu'en Amérique latine (Colombie, Mexique, Pérou) : les petits agriculteurs se regroupent pour vendre leur récolte ensemble, obtenir de meilleurs prix et financer des services communs. Le vocabulaire du développement local est donc presque identique de Yaoundé à Bogotá : \esp{cooperativa}, \esp{apoyo}, \esp{proyecto}, \esp{mercado}.
\end{savaistu}

\courssub{Carte mentale et tableau récapitulatif}

\begin{popfigure}
\begin{tikzpicture}[mindmap, grow cyclic, every node/.style={concept, align=center, font=\small}, root concept/.append style={concept color=popblue, font=\small\bfseries}, level 1/.append style={level distance=3.4cm, sibling angle=90}, level 2/.append style={level distance=2.2cm, sibling angle=45, font=\scriptsize}]
\node[root concept]{Mundo rural y movilidad}
  child[concept color=popgreen] { node {Agricultura}
    child { node {el agricultor} }
    child { node {la cosecha} }
    child { node {el ganado} }
    child { node {cultivar} }
  }
  child[concept color=poporange] { node {Movilidad}
    child { node {éxodo rural} }
    child { node {mudarse} }
    child { node {emigrar} }
    child { node {buscar trabajo} }
  }
  child[concept color=popgold] { node {Infraestructuras}
    child { node {la carretera} }
    child { node {el mercado} }
    child { node {la escuela} }
    child { node {el hospital} }
  }
  child[concept color=poppurple] { node {Desarrollo local}
    child { node {la cooperativa} }
    child { node {el apoyo} }
    child { node {el proyecto} }
    child { node {el municipio} }
  };
\end{tikzpicture}
\legende{Carte mentale des quatre champs lexicaux de la leçon.}
\end{popfigure}

\begin{center}
\begin{tabularx}{\linewidth}{|l|X|X|X|}
\hline
\rowcolor{popblueL} Campo léxico & Palabra & Traducción & Ejemplo \\
\hline
Agricultura & el agricultor & l'agriculteur & El agricultor cultiva maíz. \\
\hline
Agricultura & la cosecha & la récolte & La cosecha fue buena. \\
\hline
Agricultura & el ganado & le bétail & Mi tío tiene ganado. \\
\hline
Movilidad & el éxodo rural & l'exode rural & El éxodo rural vacía los pueblos. \\
\hline
Movilidad & mudarse & déménager & Se mudó a Duala. \\
\hline
Movilidad & buscar trabajo & chercher du travail & Emigra para buscar trabajo. \\
\hline
Infraestructuras & la carretera & la route & La carretera llega al pueblo. \\
\hline
Infraestructuras & el mercado & le marché & Vende su cosecha en el mercado. \\
\hline
Desarrollo local & la cooperativa & la coopérative & La cooperativa apoya a los jóvenes. \\
\hline
Desarrollo local & el proyecto & le projet & El municipio creó un proyecto. \\
\hline
\end{tabularx}
\end{center}

\begin{attention}[Genre et orthographe : les pièges du lexique]
\begin{itemize}
\item \esp{el agricultor / la agricultora} : le féminin se forme en \esp{-a}, comme \esp{el profesor / la profesora}.
\item \esp{la migración}, \esp{la estación} portent un accent écrit sur la dernière syllabe (mots terminés en \esp{-ción}, \esp{-sión} : oxytones en \esp{-n} → accent obligatoire).
\item \esp{la electricidad} (comme \esp{la verdad}, \esp{la ciudad}, \esp{la libertad}) est oxytone mais se termine par \esp{-d} : \textbf{pas d'accent écrit}.
\item \esp{éxodo} porte un accent sur le \esp{é} : c'est un mot proparoxyton (accentué sur l'antépénultième syllabe), comme \esp{médico}, \esp{máquina}. Prononce « é-kso-do ».
\item Ne traduis pas « le développement » par \esp{*desarollamiento} : le bon mot est \esp{el desarrollo} (avec deux \esp{ll} prononcées « y » : « dé-sa-rro-yo »).
\end{itemize}
\end{attention}

\courssub{Activités lexicales}

\begin{exoresolu}[L'intrus]
Énoncé — Dans chaque série, barre l'intrus et justifie ton choix en français.

a) \esp{el agricultor — la cosecha — el ganado — el hospital}

b) \esp{mudarse — emigrar — dejar el pueblo — cultivar}

c) \esp{la carretera — el camino — el mercado — la cooperativa}
\tcblower
Solution détaillée :

a) L'intrus est \esp{el hospital} : les trois autres mots appartiennent au champ de l'agriculture (l'agriculteur, la récolte, le bétail), tandis que l'hôpital appartient au champ des services et des infrastructures.

b) L'intrus est \esp{cultivar} : \esp{mudarse}, \esp{emigrar} et \esp{dejar el pueblo} expriment tous un départ, une mobilité ; \esp{cultivar} (cultiver) appartient au champ de l'agriculture.

c) L'intrus est \esp{la cooperativa} : \esp{la carretera}, \esp{el camino} et \esp{el mercado} désignent des lieux ou des voies concrètes (infrastructures), alors que la coopérative est une organisation humaine liée au développement local.
\end{exoresolu}

\begin{exercice}[Phrases à trous]
Complète avec le mot qui convient : \esp{éxodo rural — cooperativa — cosecha — carretera — mercado — apoyo}.

\begin{enumerate}
\item Los jóvenes dejan el pueblo : es el \ldots\ldots\ldots\ldots
\item El agricultor vende su \ldots\ldots\ldots\ldots en el \ldots\ldots\ldots\ldots
\item Gracias al \ldots\ldots\ldots\ldots del municipio, el pueblo tiene una nueva \ldots\ldots\ldots\ldots
\item La \ldots\ldots\ldots\ldots ayuda a los campesinos a vender su café.
\end{enumerate}
\end{exercice}

\begin{corrige}[Exercice : phrases à trous]
1. \esp{éxodo rural} — 2. \esp{cosecha} ; \esp{mercado} — 3. \esp{apoyo} ; \esp{carretera} — 4. \esp{cooperativa}.

Traduction des phrases complètes : 1. Les jeunes quittent le village : c'est l'exode rural. 2. L'agriculteur vend sa récolte au marché. 3. Grâce au soutien de la commune, le village a une nouvelle route. 4. La coopérative aide les paysans à vendre leur café.
\end{corrige}

\begin{exercice}[Champ lexical à compléter et mini-mots croisés]
A. Classe les mots suivants dans les quatre champs de la carte mentale : \esp{el ganado — mudarse — el hospital — el proyecto — la tierra — emigrar — la escuela — el apoyo — el camino — volver al campo}.

B. Mots croisés (définitions en français, réponses en espagnol) :
\begin{enumerate}
\item Horizontal : la récolte (7 lettres).
\item Vertical : déménager (7 lettres).
\item Horizontal : le bétail (6 lettres).
\item Vertical : la commune (9 lettres).
\end{enumerate}
\end{exercice}

\begin{corrige}[Exercice : champ lexical et mots croisés]
A. Agricultura : \esp{el ganado, la tierra}. Movilidad : \esp{mudarse, emigrar, volver al campo}. Infraestructuras : \esp{el hospital, la escuela, el camino}. Desarrollo local : \esp{el proyecto, el apoyo}.

B. 1. \esp{cosecha} — 2. \esp{mudarse} — 3. \esp{ganado} — 4. \esp{municipio}.
\end{corrige}

\begin{aretenir}[L'essentiel du lexique]
\begin{itemize}
\item Quatre champs à maîtriser : l'agriculture (\esp{el agricultor, la cosecha, el ganado}), la mobilité (\esp{el éxodo rural, mudarse, emigrar}), les infrastructures (\esp{la carretera, el mercado, los servicios}) et le développement local (\esp{la cooperativa, el apoyo, el proyecto, el municipio}).
\item Deux expressions de retour et de vie : \esp{volver al campo} (retourner à la campagne), \esp{dar vida al pueblo} (redonner vie au village).
\item Deux connecteurs pour argumenter : \esp{gracias a} (cause positive) et \esp{sin embargo} (opposition).
\item Pièges : \esp{el campo} n'est pas « le camp » ; \esp{mudarse} (déménager) n'est pas \esp{moverse} (bouger) ; attention aux accents de \esp{éxodo}, \esp{migración}, \esp{estación} (oxytones en \esp{-n}/\esp{-s}) et à l'absence d'accent sur \esp{electricidad} (oxytones en consonne autre que \esp{n}/\esp{s}).
\end{itemize}
\end{aretenir}

Ce vocabulaire est maintenant ton outil de travail. Dans la prochaine section, nous verrons comment le mettre en mouvement dans le temps : le \esp{pretérito indefinido} et l'\esp{imperfecto} te permettront de raconter au passé une expérience rurale, comme Samuel dans notre texte.

\coursec{Grammaire : prétérit indéfini et imparfait}

Dans la section précédente, nous avons construit le lexique du monde rural : \esp{el campo}, \esp{la cosecha}, \esp{el éxodo rural}, \esp{la cooperativa}, \esp{el agricultor}, \esp{la migración}, \esp{el desarrollo local}. Or, pour raconter la vie d'un village, son passé, ses transformations, il faut maîtriser les deux grands temps du passé en espagnol : l'\cle{imparfait} (\esp{pretérito imperfecto}) et le \cle{prétérit indéfini} (\esp{pretérito indefinido}). C'est l'outil grammatical indispensable pour raconter une expérience rurale, objectif de notre leçon.

\courssub{Observation : deux façons de voir le passé}

\begin{textebox}[Deux phrases du texte support]
Mis tíos \textbf{dejaron} el pueblo en 2010 y \textbf{se instalaron} en Douala. Antes, mi abuelo \textbf{se levantaba} temprano cada día, \textbf{trabajaba} en el campo y \textbf{vendía} la cosecha en el mercado. \textbf{Era} una vida dura, pero tranquila.
\end{textebox}

Traduction : « Mes oncles ont quitté le village en 2010 et se sont installés à Douala. Avant, mon grand-père se levait tôt chaque jour, travaillait aux champs et vendait la récolte au marché. C'était une vie dure, mais tranquille. »

Observons les verbes en gras :

\begin{itemize}
\item \esp{dejaron}, \esp{se instalaron} : actions \textbf{uniques, datées, achevées} — elles font avancer le récit. C'est le \cle{prétérit indéfini}.
\item \esp{se levantaba}, \esp{trabajaba}, \esp{vendía}, \esp{era} : \textbf{habitudes répétées} et \textbf{description} d'une situation durable — c'est le décor. C'est l'\cle{imparfait}.
\end{itemize}

\begin{propriete}[Règle de choix fondamentale]
\begin{itemize}
\item \textbf{IMPARFAIT} (\esp{imperfecto}) = habitude, répétition, description, décor, action en cours ou simultanée, âge, heure. Le passé est vu « de l'intérieur », sans début ni fin précis.
\item \textbf{INDÉFINI} (\esp{indefinido}) = action unique, ponctuelle, datée, achevée, succession d'événements qui font avancer le récit. Le passé est vu « de l'extérieur », comme un point fermé.
\end{itemize}
\end{propriete}

\begin{exemplebox}[Les deux temps face à face]
\esp{Cuando era niño, jugaba en el campo.} « Quand j'étais enfant, je jouais à la campagne. » (imparfait : habitude et description)\par
\esp{El año pasado visité el pueblo de mis abuelos.} « L'année dernière, j'ai visité le village de mes grands-parents. » (indéfini : événement unique et daté)
\end{exemplebox}

\courssub{Formation de l'imparfait}

L'imparfait est le temps le plus régulier de l'espagnol : il n'a que \textbf{trois verbes irréguliers}.

\begin{center}
\begin{tabular}{|l|l|l|}
\hline
\rowcolor{popblueL} \textbf{-ar : hablar} & \textbf{-er : comer} & \textbf{-ir : vivir} \\
\hline
habl\textbf{aba} & com\textbf{ía} & viv\textbf{ía} \\
habl\textbf{abas} & com\textbf{ías} & viv\textbf{ías} \\
habl\textbf{aba} & com\textbf{ía} & viv\textbf{ía} \\
habl\textbf{ábamos} & com\textbf{íamos} & viv\textbf{íamos} \\
habl\textbf{abais} & com\textbf{íais} & viv\textbf{íais} \\
habl\textbf{aban} & com\textbf{ían} & viv\textbf{ían} \\
\hline
\end{tabular}
\end{center}

\begin{propriete}[Terminaisons de l'imparfait]
\begin{itemize}
\item Verbes en \textbf{-ar} : \textbf{-aba, -abas, -aba, -ábamos, -abais, -aban}.
\item Verbes en \textbf{-er} et \textbf{-ir} (mêmes terminaisons) : \textbf{-ía, -ías, -ía, -íamos, -íais, -ían}.
\item Attention aux accents : \esp{ábamos} (1\iere{} personne du pluriel en -ar) et toutes les formes en \esp{-ía} portent un accent sur le \esp{í}.
\end{itemize}
\end{propriete}

\begin{propriete}[Les trois seuls irréguliers de l'imparfait]
\begin{itemize}
\item \esp{ir} (aller) : iba, ibas, iba, íbamos, ibais, iban.
\item \esp{ser} (être) : era, eras, era, éramos, erais, eran.
\item \esp{ver} (voir) : veía, veías, veía, veíamos, veíais, veían.
\end{itemize}
\end{propriete}

\begin{exemplebox}[Imparfait en contexte rural]
\esp{Mi abuela cocinaba ñame todos los días.} « Ma grand-mère cuisinait de l'igname tous les jours. »\par
\esp{Los campesinos iban al mercado los sábados.} « Les paysans allaient au marché le samedi. »\par
\esp{El pueblo era pequeño y las casas eran de barro.} « Le village était petit et les maisons étaient en terre. »
\end{exemplebox}

\courssub{Formation du prétérit indéfini}

\begin{center}
\begin{tabular}{|l|l|l|}
\hline
\rowcolor{popgreenL} \textbf{-ar : hablar} & \textbf{-er : comer} & \textbf{-ir : vivir} \\
\hline
habl\textbf{é} & com\textbf{í} & viv\textbf{í} \\
habl\textbf{aste} & com\textbf{iste} & viv\textbf{iste} \\
habl\textbf{ó} & com\textbf{ió} & viv\textbf{ió} \\
habl\textbf{amos} & com\textbf{imos} & viv\textbf{imos} \\
habl\textbf{asteis} & com\textbf{isteis} & viv\textbf{isteis} \\
habl\textbf{aron} & com\textbf{ieron} & viv\textbf{ieron} \\
\hline
\end{tabular}
\end{center}

\begin{propriete}[Terminaisons de l'indéfini]
\begin{itemize}
\item Verbes en \textbf{-ar} : \textbf{-é, -aste, -ó, -amos, -asteis, -aron}.
\item Verbes en \textbf{-er} et \textbf{-ir} : \textbf{-í, -iste, -ió, -imos, -isteis, -ieron}.
\item Les accents sont obligatoires à la 1\iere{} et à la 3\ieme{} personne du singulier : \esp{hablé}, \esp{habló}, \esp{comí}, \esp{comió}, \esp{viví}, \esp{vivió}. Sans accent, le sens change : \esp{hablo} = « je parle » (présent), \esp{habló} = « il parla ».
\end{itemize}
\end{propriete}

\begin{propriete}[Indéfini : les irréguliers fréquents (radical spécial, terminaisons sans accent)]
\begin{itemize}
\item \esp{ir} / \esp{ser} : fui, fuiste, fue, fuimos, fuisteis, fueron.
\item \esp{hacer} : hice, hiciste, hizo, hicimos, hicisteis, hicieron.
\item \esp{tener} : tuve, tuviste, tuvo, tuvimos, tuvisteis, tuvieron.
\item \esp{estar} : estuve, estuviste, estuvo, estuvimos, estuvisteis, estuvieron.
\item \esp{poder} : pude, pudiste, pudo, pudimos, pudisteis, pudieron.
\item \esp{poner} : puse, pusiste, puso, pusimos, pusisteis, pusieron.
\item \esp{decir} : dije, dijiste, dijo, dijimos, dijisteis, dijeron.
\item \esp{venir} : vine, viniste, vino, vinimos, vinisteis, vinieron.
\item \esp{querer} : quise, quisiste, quiso, quisimos, quisisteis, quisieron.
\item \esp{traer} : traje, trajiste, trajo, trajimos, trajisteis, trajeron.
\item \esp{dar} : di, diste, dio, dimos, disteis, dieron (pas d'accents).
\item \esp{ver} : vi, viste, vio, vimos, visteis, vieron (pas d'accents).
\end{itemize}
\end{propriete}

\begin{attention}[Changements orthographiques à la 3\ieme{} personne de l'indéfini]
Certains verbes réguliers changent d'orthographe à la 3\ieme{} personne (singulier \esp{-ió/-yó} et pluriel \esp{-ieron/-yeron}) pour conserver le son :
\begin{itemize}
\item \textbf{-uir} $\rightarrow$ \textbf{y} : \esp{construir} $\rightarrow$ \esp{construyó, construyeron} ; \esp{huir} $\rightarrow$ \esp{huyó, huyeron}.
\item \textbf{-eer / -oír} $\rightarrow$ \textbf{y} : \esp{leer} $\rightarrow$ \esp{leyó, leyeron} ; \esp{oír} $\rightarrow$ \esp{oyó, oyeron}.
\item \textbf{-uir} sans accent (seguir, conseguir) : \esp{í} $\rightarrow$ \textbf{y} : \esp{siguió, siguieron}.
\item \textbf{Changement de radical e $\rightarrow$ i / o $\rightarrow$ u} (verbes en -ir seulement) : \esp{pedir} $\rightarrow$ \esp{pidió, pidieron} ; \esp{dormir} $\rightarrow$ \esp{durmió, durmieron} ; \esp{servir} $\rightarrow$ \esp{sirvió, sirvieron}.
\end{itemize}
\end{attention}

\begin{exemplebox}[Indéfini en contexte rural]
\esp{Mis tíos dejaron el pueblo en 2010.} « Mes oncles quittèrent le village en 2010. »\par
\esp{Ayer coseché el cacao con mi padre.} « Hier, j'ai récolté le cacao avec mon père. »\par
\esp{La cooperativa puso un almacén nuevo el mes pasado.} « La coopérative a installé un nouveau magasin le mois dernier. »\par
\esp{El joven vino a la ciudad y encontró trabajo.} « Le jeune homme vint en ville et trouva du travail. »
\end{exemplebox}

\begin{attention}[ir et ser : même forme à l'indéfini]
\esp{ir} (aller) et \esp{ser} (être) partagent exactement les mêmes formes à l'indéfini : \esp{fui, fuiste, fue, fuimos, fuisteis, fueron}. Seul le contexte permet de trancher :
\begin{itemize}
\item \esp{Fui al campo el sábado.} « Je suis allé à la campagne samedi. » (verbe \esp{ir})
\item \esp{Fue un buen día.} « Ce fut une bonne journée. » (verbe \esp{ser})
\end{itemize}
\end{attention}

\courssub{Emplois détaillés et marqueurs temporels}

\begin{propriete}[Emplois de l'imparfait]
\begin{enumerate}
\item \textbf{Habitude, répétition} : \esp{Cada día se levantaba a las cinco.} « Chaque jour, il se levait à cinq heures. »
\item \textbf{Description, décor} : \esp{El cielo estaba azul y hacía calor.} « Le ciel était bleu et il faisait chaud. »
\item \textbf{Action en cours} : \esp{Trabajaba en el campo.} « Il travaillait aux champs (à ce moment-là). »
\item \textbf{Âge et heure} : \esp{Tenía diez años. Eran las seis.} « Il avait dix ans. Il était six heures. »
\end{enumerate}
Marqueurs typiques : \esp{cada día} (chaque jour), \esp{siempre} (toujours), \esp{a menudo} (souvent), \esp{mientras} (pendant que), \esp{antes} (autrefois), \esp{todos los años} (tous les ans), \esp{en aquella época} (à cette époque).
\end{propriete}

\begin{propriete}[Emplois de l'indéfini]
\begin{enumerate}
\item \textbf{Action ponctuelle achevée} : \esp{Compró una moto.} « Il a acheté une moto. »
\item \textbf{Événement daté} : \esp{En 2010, mi familia se mudó a Yaoundé.} « En 2010, ma famille a déménagé à Yaoundé. »
\item \textbf{Succession d'actions} : \esp{Se levantó, desayunó y salió.} « Il se leva, déjeuna et sortit. »
\item \textbf{Début ou fin d'une action} : \esp{Empezó a llover.} « Il se mit à pleuvoir. »
\end{enumerate}
Marqueurs typiques : \esp{ayer} (hier), \esp{anoche} (hier soir), \esp{en 2010}, \esp{el año pasado} (l'année dernière), \esp{de repente} (soudain), \esp{una vez} (une fois), \esp{entonces} (alors), \esp{el otro día} (l'autre jour).
\end{propriete}

\begin{aretenir}[La combinaison des deux temps : la structure narrative]
Dans un récit, l'imparfait peint le \textbf{décor} et l'indéfini fait surgir l'\textbf{événement} :\par
\esp{Mientras trabajaba en el campo, llegó la noticia.} « Pendant qu'il travaillait aux champs, la nouvelle arriva. »\par
\esp{Cuando vivíamos en el pueblo, construyeron una carretera nueva.} « Quand nous vivions au village, on construisit une nouvelle route. »
\end{aretenir}

\begin{popfigure}
\begin{tikzpicture}
\draw[-{Stealth}, thick, popdark] (0,0) -- (11,0) node[right, font=\small] {pasado};
\draw[popblue, very thick, decorate, decoration={snake, amplitude=1.2mm, segment length=6mm}] (0.6,0.9) -- (10.2,0.9);
\node[popblue, font=\small, align=center] at (5.4,1.7) {IMPERFECTO : decor, hábitos\\ \esp{trabajaba, se levantaba, era}};
\fill[poporange] (2.2,0) circle (2.5pt);
\fill[poporange] (5.6,0) circle (2.5pt);
\fill[poporange] (8.6,0) circle (2.5pt);
\node[poporange, font=\small, align=center] at (2.2,-0.7) {\esp{dejaron}\\ (2010)};
\node[poporange, font=\small, align=center] at (5.6,-0.7) {\esp{llegó}\\ la noticia};
\node[poporange, font=\small, align=center] at (8.6,-0.7) {\esp{volvió}\\ ayer};
\node[poporange, font=\small, align=center] at (5.6,-1.6) {INDEFINIDO : acciones puntuales};
\end{tikzpicture}
\legende{Frise des temps : la vague bleue (imparfait) forme le décor continu ; les points orange (indéfini) sont les événements ponctuels.}
\end{popfigure}

\begin{popfigure}
\begin{tikzpicture}[font=\small]
\node[draw, fill=popblueL, rounded corners, align=center, minimum width=3.2cm] (imp) at (0,0) {IMPERFECTO\\ -ar : -aba\\ -er/-ir : -ía};
\node[draw, fill=popgreenL, rounded corners, align=center, minimum width=3.2cm] (ind) at (6.5,0) {INDEFINIDO\\ -ar : -é, -ó\\ -er/-ir : -í, -ió};
\node[draw, fill=poporangeL, rounded corners, align=center, minimum width=3.2cm] (irr) at (3.25,-2.2) {Irregulares\\ imp. : ir, ser, ver\\ ind. : fui, hice, tuve...};
\draw[-{Stealth}, thick, popdark] (imp) -- (irr);
\draw[-{Stealth}, thick, popdark] (ind) -- (irr);
\end{tikzpicture}
\legende{Carte mentale des terminaisons : deux temps réguliers, quelques irréguliers à mémoriser.}
\end{popfigure}

\courssub{Comparaison avec le français}

\begin{center}
\begin{tabularx}{\linewidth}{|X|X|X|}
\hline
\rowcolor{popblueL} \textbf{Français} & \textbf{Español} & \textbf{Remarque} \\
\hline
je jouais (imparfait) & jugaba & imparfait $\approx$ imparfait français \\
\hline
j'allais & iba & irrégulier dans les deux langues \\
\hline
j'ai visité / je visitai & visité & indéfini $\approx$ passé composé et passé simple \\
\hline
je suis allé & fui & une seule forme en espagnol pour \esp{ir} et \esp{ser} \\
\hline
pendant que je travaillais & mientras trabajaba & \esp{mientras} appelle l'imparfait \\
\hline
soudain, il arriva & de repente, llegó & \esp{de repente} appelle l'indéfini \\
\hline
\end{tabularx}
\end{center}

\begin{attention}[Erreur fréquente des francophones]
Ne traduis pas le passé composé français automatiquement par l'indéfini, ni l'imparfait français par... n'importe quoi ! Demande-toi d'abord : habitude ou événement unique ?
\begin{itemize}
\item « J'allais souvent au village. » = \esp{Iba a menudo al pueblo.} (imparfait : habitude)
\item « Je suis allé au village hier. » = \esp{Fui al pueblo ayer.} (indéfini : événement daté)
\end{itemize}
Autre piège : en français « il parlait » et « il a parlé » se distinguent à l'oreille ; en espagnol, l'accent écrit fait tout : \esp{hablaba} (il parlait) contre \esp{habló} (il parla).
Attention aux \textbf{faux amis} : \esp{asistir} = assister (être présent), pas « aider » (\esp{ayudar}) ; \esp{realizar} = réaliser (accomplir), pas « réaliser » au sens de « comprendre » (\esp{darse cuenta}).
\end{attention}

\begin{methode}[Conjuguer au passé en quatre étapes]
\begin{enumerate}
\item \textbf{Repérer le marqueur temporel} : \esp{ayer}, \esp{en 2010}, \esp{de repente} orientent vers l'indéfini ; \esp{cada día}, \esp{siempre}, \esp{mientras}, \esp{antes} vers l'imparfait.
\item \textbf{Se poser la question clé} : habitude, description, décor, action en cours ? $\rightarrow$ imparfait. Événement unique, daté, achevé, succession ? $\rightarrow$ indéfini.
\item \textbf{Choisir le temps} et appliquer les terminaisons du groupe (-ar ou -er/-ir).
\item \textbf{Vérifier l'irrégularité} :
\begin{itemize}
\item Imparfait : seulement \esp{ir} (\esp{iba}), \esp{ser} (\esp{era}), \esp{ver} (\esp{veía}).
\item Indéfini : radicaux spéciaux (\esp{fui}, \esp{hice}, \esp{tuve}, \esp{estuve}, \esp{pude}, \esp{puse}, \esp{dije}, \esp{vine}, \esp{quise}, \esp{traje}, \esp{di}, \esp{vi}...) et changements orthographiques 3\ieme{} pers. (\esp{construyó}, \esp{leyó}, \esp{durmió}, \esp{pidió}).
\end{itemize}
\end{enumerate}
\end{methode}

\begin{exoresolu}[Choisis et conjugue le verbe entre parenthèses]
Énoncé. Conjuguez au passé (imparfait ou indéfini) en justifiant :\par
1. Cuando (ser) \dots\ niño, (pasar) \dots\ las vacaciones en el pueblo.\par
2. Ayer mis primos (llegar) \dots\ de Douala y (traer) \dots\ regalos.\par
3. Mientras mi abuelo (dormir) \dots, (empezar) \dots\ a llover.\par
4. El año pasado la cooperativa (construir) \dots\ un almacén.
\tcblower
Solution détaillée :\par
1. \esp{Cuando era niño, pasaba las vacaciones en el pueblo.} « Quand j'étais enfant, je passais les vacances au village. » Marqueur \esp{cuando} + situation durable et habitude répétée : imparfait. \esp{ser} est irrégulier : \esp{era}.\par
2. \esp{Ayer mis primos llegaron de Douala y trajeron regalos.} « Hier, mes cousins sont arrivés de Douala et ont apporté des cadeaux. » Marqueur \esp{ayer} + succession d'événements : indéfini. \esp{llegar} régulier (\esp{llegaron}) ; \esp{traer} irrégulier : \esp{trajeron}.\par
3. \esp{Mientras mi abuelo dormía, empezó a llover.} « Pendant que mon grand-père dormait, il se mit à pleuvoir. » \esp{mientras} + action en cours : imparfait (\esp{dormía}) ; début d'action : indéfini (\esp{empezó} avec accent).\par
4. \esp{El año pasado la cooperativa construyó un almacén.} « L'année dernière, la coopérative a construit un magasin. » Marqueur \esp{el año pasado} : indéfini ; \esp{construir} : \esp{construyó} (changement i $\rightarrow$ y à la 3\ieme{} personne).
\end{exoresolu}

\begin{experience}[Raconter son village : jeu de rôle guidé]
\textbf{Consigne :} En binôme. L'élève A est un jeune revenu au village pour les vacances ; l'élève B est un grand-parent qui décrit le village d'avant.
\begin{enumerate}
\item \textbf{Préparation (5 min) :} A note 3 questions à l'indéfini (\esp{¿Qué pasó...?}, \esp{¿Cuándo llegaron...?}). B note 3 descriptions à l'imparfait (\esp{Antes había...}, \esp{La gente vivía...}, \esp{Era...}).
\item \textbf{Dialogue (3 min) :} A interroge, B répond. Exemple :
\esp{A: ¿Cuándo llegó la carretera?}
\esp{B: Antes no había carretera, el camino era de tierra. Llegó en 2015.}
\item \textbf{Restitution :} Chaque binôme présente une phrase « décor + événement » type \esp{Mientras..., pasó...}.
\end{enumerate}
\end{experience}

\begin{savaistu}[Le monde rural au Cameroun et en Guinée équatoriale]
Le Cameroun compte environ 45\% de population rurale. L'\esp{exodo rural} vers Douala, Yaoundé ou les plantations du Sud-Ouest est massif. Des initiatives de \esp{desarrollo local} émergent : coopératives cacaoyères (\esp{cooperativas cacaoteras}) dans le Centre et le Sud, groupements de femmes pour la transformation du manioc (\esp{ñame}, \esp{mandioca}), projets d'\esp{agroforestería} appuyés par la coopération espagnole. En Guinée équatoriale (pays hispanophone), l'île de Bioko et la région continentale de Río Muni connaissent des dynamiques similaires : modernisation de l'agriculture, routes (\esp{carreteras}) et éducation rurale. Raconter ces changements en espagnol, c'est utiliser l'imparfait pour l'ancien temps (\esp{antes no había luz, se caminaba horas}) et l'indéfini pour les réalisations récentes (\esp{el año pasado instalaron paneles solares, construyeron un centro de salud}).
\end{savaistu}

\begin{exercice}[À ton tour]
Mets les verbes entre parenthèses à l'imparfait ou à l'indéfini, puis traduis les phrases 2 et 4 en français :\par
1. Antes, nosotros (ir) \dots\ al río todos los domingos.\par
2. De repente, (sonar) \dots\ el teléfono mientras (comer) \dots.\par
3. En 2015, mi tío (vender) \dots\ su cosecha de cacao y (comprar) \dots\ una moto.\par
4. El pueblo (ser) \dots\ muy tranquilo, pero un día (llegar) \dots\ la carretera nueva.\par
5. ¿Qué (hacer) \dots\ tú ayer por la tarde?
\end{exercice}

\begin{corrige}[Exercice]
1. \esp{Antes, nosotros íbamos al río todos los domingos.} (habitude : imparfait de \esp{ir}, irrégulier).\par
2. \esp{De repente, sonó el teléfono mientras comíamos.} « Soudain, le téléphone a sonné pendant que nous mangions. » (événement : indéfini \esp{sonó} ; action en cours : imparfait \esp{comíamos}).\par
3. \esp{En 2015, mi tío vendió su cosecha de cacao y compró una moto.} (date précise + succession : indéfini).\par
4. \esp{El pueblo era muy tranquilo, pero un día llegó la carretera nueva.} « Le village était très tranquille, mais un jour, la nouvelle route arriva. » (description : imparfait \esp{era} ; événement : indéfini \esp{llegó}).\par
5. \esp{¿Qué hiciste ayer por la tarde?} (événement daté par \esp{ayer} : indéfini irrégulier de \esp{hacer}).
\end{corrige}

\begin{aretenir}[L'essentiel]
\begin{itemize}
\item Imparfait = décor, habitude, description : terminaisons \esp{-aba} / \esp{-ía} ; seuls irréguliers : \esp{iba}, \esp{era}, \esp{veía}.
\item Indéfini = événement ponctuel, daté, achevé : terminaisons \esp{-é/-ó} et \esp{-í/-ió} ; nombreux irréguliers (\esp{fui}, \esp{hice}, \esp{tuve}, \esp{traje}, \esp{di}, \esp{vi}...) et changements orthographiques 3\ieme{} pers. (\esp{-ió/-yó}, \esp{-ieron/-yeron}).
\item La combinaison \esp{mientras} + imparfait, puis indéfini, est la structure reine du récit au passé : \esp{Mientras trabajaba en el campo, llegó la noticia.}
\end{itemize}
\end{aretenir}

\coursec{Méthode du commentaire et commentaire modèle}

Dans la section précédente, nous avons appris à raconter le passé avec le \cle{prétérit indéfini} et l'\cle{imparfait} : \esp{El joven llegó a la ciudad} (action ponctuelle, achevée) et \esp{El campo era tranquilo y la gente trabajaba unida} (description, habitude). Ces deux temps seront vos meilleurs outils pour commenter un texte narratif ou journalistique. Nous allons maintenant utiliser cette compétence grammaticale dans l'exercice roi de la classe de Première : le \cle{commentaire de texte} (\esp{comentario de texto}).

\courssub{Qu'est-ce qu'un commentaire de texte ?}

\begin{definition}[Le comentario de texto]
Le \esp{comentario de texto} est un exercice d'expression écrite dans lequel l'élève présente un texte, en dégage les idées principales, les analyse de manière organisée et propose une conclusion personnelle mais objective. En espagnol, il se rédige entièrement dans la langue étudiée, avec un plan rigoureux en trois parties : \esp{introducción}, \esp{desarrollo}, \esp{conclusión}.
\end{definition}

Contrairement au résumé (\esp{resumen}), qui se contente de reformuler le texte, le commentaire \emph{analyse} : il répond à une question (la \cle{problématique}) et montre comment le texte y répond. Au baccalauréat et dans les évaluations MINESEC, le commentaire est noté sur le respect du plan, l'exactitude linguistique et la richesse du vocabulaire.

\courssub{Le plan en trois parties}

\begin{propriete}[Structure obligatoire du commentaire]
\begin{itemize}
\item \textbf{Introducción} : une phrase d'accroche sur le thème général, la présentation du texte (auteur ou source, thème, type de texte), la problématique et l'annonce du plan.
\item \textbf{Desarrollo} : deux axes (ou idées principales), chacun développé dans un paragraphe avec des exemples tirés du texte et des connecteurs logiques.
\item \textbf{Conclusión} : le bilan de l'analyse (réponse à la problématique) et une ouverture vers un thème voisin.
\end{itemize}
\end{propriete}

\begin{popfigure}
\begin{tikzpicture}[node distance=0.6cm and 1.1cm, every node/.style={font=\small}]
\node[draw=popblue, fill=popblueL, rounded corners, thick, align=center, text width=3.6cm] (intro) {\textbf{Introducción}\\ Accroche + texto\\ + problemática + plan};
\node[draw=poporange, fill=poporangeL, rounded corners, thick, align=center, text width=3.2cm, right=of intro] (idea1) {\textbf{Idea 1}\\ Causas y consecuencias\\ del éxodo rural};
\node[draw=popgreen, fill=popgreenL, rounded corners, thick, align=center, text width=3.2cm, below=of idea1] (idea2) {\textbf{Idea 2}\\ Soluciones de\\ desarrollo local};
\node[draw=poppurple, fill=poppurpleL, rounded corners, thick, align=center, text width=3.6cm, right=of idea1, yshift=-1.1cm] (conclu) {\textbf{Conclusión}\\ Bilan + apertura};
\node[draw=popdark, dashed, rounded corners, fit=(idea1)(idea2), inner sep=0.25cm, label={[popdark, font=\small\bfseries]above:Desarrollo}] (des) {};
\draw[-{Stealth}, very thick, popblue] (intro) -- (des.west);
\draw[-{Stealth}, very thick, poppurple] (des.east) -- (conclu);
\end{tikzpicture}
\legende{Schéma du plan de commentaire : Introducción $\rightarrow$ Desarrollo (Idea 1 + Idea 2) $\rightarrow$ Conclusión.}
\end{popfigure}

\courssub{Rédiger l'introduction : la méthode pas à pas}

\begin{methode}[Comment rédiger la introducción en quatre étapes]
\begin{enumerate}
\item \textbf{Phrase d'accroche} : une phrase générale sur le thème. \esp{En muchos países de África, el campo se vacía y las ciudades crecen muy rápido.} « Dans de nombreux pays d'Afrique, la campagne se vide et les villes grandissent très vite. »
\item \textbf{Présentation du texte} : source, thème, type. \esp{El texto que vamos a comentar es un artículo periodístico sobre el éxodo rural y el desarrollo local en Camerún.} « Le texte que nous allons commenter est un article de presse sur l'exode rural et le développement local au Cameroun. »
\item \textbf{Problématique} : la question qui guide l'analyse. \esp{¿Por qué los jóvenes abandonan el campo y qué soluciones existen?} « Pourquoi les jeunes abandonnent-ils la campagne et quelles solutions existent ? »
\item \textbf{Annonce du plan} : \esp{En primer lugar, analizaremos las causas y consecuencias del éxodo rural; en segundo lugar, estudiaremos las iniciativas de desarrollo local.} « D'abord, nous analyserons les causes et les conséquences de l'exode rural ; ensuite, nous étudierons les initiatives de développement local. »
\end{enumerate}
\end{methode}

\begin{attention}[Le piège du « yo pienso que »]
Les francophones transposent souvent « je pense que » par \esp{yo pienso que} répété à chaque phrase. C'est lourd et trop subjectif pour un commentaire. Préférez des tournures objectives :
\begin{itemize}
\item \esp{El texto muestra que...} « Le texte montre que... »
\item \esp{El autor explica que...} « L'auteur explique que... »
\item \esp{Se puede decir que...} « On peut dire que... »
\item \esp{Según el texto...} « Selon le texte... »
\end{itemize}
Autre piège : ne traduisez jamais « d'abord » par \esp{primero lugar} : la forme correcte est \esp{en primer lugar}.
\end{attention}

\courssub{Les connecteurs logiques : la colonne vertébrale du commentaire}

\begin{center}
\begin{tabularx}{\linewidth}{|X|X|X|}
\hline
\rowcolor{popblueL} \textbf{Español} & \textbf{Français} & \textbf{Emploi} \\
\hline
\esp{en primer lugar} & d'abord, en premier lieu & ouvrir l'axe 1 \\
\hline
\esp{además} & de plus, en outre & ajouter un argument \\
\hline
\esp{sin embargo} & cependant, pourtant & opposer, nuancer \\
\hline
\esp{por ejemplo} & par exemple & illustrer avec le texte \\
\hline
\esp{en segundo lugar} & ensuite, en second lieu & ouvrir l'axe 2 \\
\hline
\esp{por eso / por lo tanto} & c'est pourquoi, donc & conséquence \\
\hline
\esp{en conclusión} & en conclusion & ouvrir le bilan \\
\hline
\end{tabularx}
\end{center}

\begin{exemplebox}[Connecteurs en contexte]
\esp{En primer lugar, el texto muestra las causas del éxodo rural.} « D'abord, le texte montre les causes de l'exode rural. »

\esp{Además, explica las consecuencias para los pueblos.} « De plus, il explique les conséquences pour les villages. »

\esp{Sin embargo, existen soluciones como las cooperativas.} « Cependant, il existe des solutions comme les coopératives. »

\esp{En conclusión, el desarrollo local es una esperanza para el campo.} « En conclusion, le développement local est un espoir pour la campagne. »
\end{exemplebox}

\courssub{Les deux axes de développement appliqués à notre texte}

\textbf{Axe 1 : les causes et les conséquences de l'exode rural.} Dans le texte support, les causes sont économiques (pauvreté du \esp{campo}, absence de \esp{carreteras} et de transport), sociales (manque d'écoles et d'hôpitaux) et culturelles (l'attrait de la ville). Les conséquences : le vieillissement des villages, la baisse des \esp{cosechas}, la surpopulation des villes comme Douala et Yaoundé. Pour cet axe, utilisez l'imparfait pour décrire la vie rurale d'autrefois (\esp{la gente vivía de la agricultura}) et le prétérit indéfini pour les événements (\esp{muchos jóvenes se fueron a la ciudad}).

\textbf{Axe 2 : les solutions de développement local et leur portée.} Le texte présente les \esp{cooperativas} de cacao, les projets d'électrification, les routes rurales et le retour de certains jeunes diplômés. La portée : ces initiatives redonnent espoir, créent des emplois et freinent la \esp{migración}. Montrez que ces solutions sont positives (\esp{además, dan trabajo a las familias}) mais limitées (\esp{sin embargo, todavía falta dinero}).

\courssub{Commentaire modèle rédigé}

\begin{textebox}[Comentario de texto : modelo completo]
En muchos países de África, los jóvenes abandonan el campo para buscar una vida mejor en la ciudad. El texto que vamos a comentar es un artículo periodístico sobre el éxodo rural y el desarrollo local en Camerún. ¿Por qué se vacían los pueblos y qué soluciones existen? En primer lugar, analizaremos las causas y consecuencias de este fenómeno; en segundo lugar, estudiaremos las iniciativas de desarrollo local.

En primer lugar, el texto muestra las causas del éxodo rural: la pobreza, la falta de carreteras y la ausencia de escuelas y hospitales. Antes, la gente vivía de la agricultura y trabajaba unida, pero muchos jóvenes se fueron a Douala o a Yaoundé. Además, el éxodo tiene consecuencias graves: los pueblos envejecen y las cosechas bajan.

En segundo lugar, el artículo presenta soluciones de desarrollo local, por ejemplo las cooperativas de cacao y las nuevas carreteras rurales. Estas iniciativas crean empleo y devuelven la esperanza a las familias. Sin embargo, todavía falta dinero y apoyo del Estado.

En conclusión, el desarrollo local es una esperanza para el campo camerunés. Se puede decir que el futuro de los pueblos depende de estos proyectos. Este tema nos invita a pensar también en la educación rural.
\end{textebox}

\textbf{Glossaire :} \esp{se vacían} = se vident ; \esp{fenómeno} = phénomène ; \esp{envejecen} = vieillissent ; \esp{devuelven} = redonnent, rendent ; \esp{apoyo} = soutien ; \esp{falta} = il manque.

Observez le modèle : chaque partie est visible, les connecteurs \esp{en primer lugar}, \esp{además}, \esp{sin embargo}, \esp{en conclusión} structurent le propos, et les temps du passé alternent correctement (\esp{vivía}, imparfait de description ; \esp{se fueron}, indéfini d'action achevée).

\begin{exoresolu}[Transformer des notes en phrase de commentaire]
Énoncé : à partir de ces notes — « texto / mostrar / causas / éxodo rural » — rédigez une phrase complète de développement avec un connecteur, puis traduisez-la.
\tcblower
Solution : on choisit le connecteur d'ouverture de l'axe 1, on conjugue le verbe à la troisième personne du présent (le texte « montre », vérité actuelle) et on complète avec l'article défini : \esp{En primer lugar, el texto muestra las causas del éxodo rural.} Traduction : « D'abord, le texte montre les causes de l'exode rural. » Erreurs à éviter : ne pas écrire \esp{el texto mostrar} (infinitif non conjugué) ni oublier l'article \esp{las} devant \esp{causas}.
\end{exoresolu}

\courssub{Entraînement guidé et critères d'évaluation}

\begin{exercice}[Rédiger une introducción]
À partir du texte support de cette leçon (l'article sur l'exode rural et le développement local au Cameroun), rédigez une introduction complète de 4 à 6 phrases en suivant la méthode : accroche, présentation du texte, problématique, annonce du plan. Utilisez au moins deux expressions de la boîte méthode.
\end{exercice}

\begin{corrige}[Exercice : introducción]
Modèle de réponse attendue : \esp{Hoy en día, muchos pueblos africanos pierden a sus jóvenes. El texto que vamos a comentar es un artículo sobre el éxodo rural y el desarrollo local en Camerún. El autor explica por qué la gente deja el campo y presenta soluciones concretas. ¿Cómo frenar el éxodo rural? En primer lugar, estudiaremos las causas y consecuencias del fenómeno; en segundo lugar, analizaremos las iniciativas de desarrollo local.} Critères de réussite : quatre étapes présentes, verbes correctement conjugués au présent, annonce du plan avec \esp{en primer lugar / en segundo lugar}.
\end{corrige}

\begin{aretenir}[Critères d'évaluation type Bac]
\begin{itemize}
\item \textbf{Respect du plan} : introduction, développement en deux axes, conclusion visibles et équilibrées.
\item \textbf{Exactitude linguistique} : accords, conjugaisons, accents (\esp{éxodo}, \esp{migración}, \esp{conclusión}).
\item \textbf{Emploi des temps du passé} : imparfait pour les descriptions et habitudes, indéfini pour les actions achevées.
\item \textbf{Richesse du lexique} : vocabulaire du thème (\esp{campo, cosecha, cooperativa, desarrollo local}) et connecteurs variés.
\end{itemize}
\end{aretenir}

\begin{savaistu}[Le commentaire, une tradition hispanique]
En Espagne, le \esp{comentario de texto} est l'épreuve centrale de l'examen d'entrée à l'université (la \esp{Selectividad}, aujourd'hui \esp{EBAU}). Les lycéens espagnols y commentent des textes littéraires ou journalistiques, exactement comme vous le faites en espagnol LV2 au Cameroun : maîtriser cet exercice, c'est acquérir une compétence reconnue dans tout le monde hispanophone.
\end{savaistu}

\coursec{Production guidée et entraînement à la traduction}

Après avoir appris à lire et à commenter un texte sur le monde rural, il est temps de passer à la \cle{production} : à votre tour de raconter, de traduire et de réutiliser le lexique et les temps du passé étudiés dans cette leçon. Cette dernière étape est essentielle : c'est elle qui transforme des connaissances passives (comprendre) en compétences actives (écrire et traduire). Nous allons procéder en trois temps : une production écrite guidée, un entraînement au thème (français $\rightarrow$ espagnol), puis un entraînement à la version (espagnol $\rightarrow$ français), avec une correction collective et une réécriture individuelle.

\courssub{Production écrite guidée : raconter une visite au village}

\textbf{Consigne.} \esp{Cuenta una visita al pueblo.} Raconte au passé une visite au village (le tien, celui de tes grands-parents, ou un village imaginaire du Cameroun). Tu répondras obligatoirement aux cinq questions de la grille suivante :

\begin{center}
\begin{tabularx}{\linewidth}{|l|X|X|}
\hline
\rowcolor{popblueL} \textbf{Question} & \textbf{Rôle dans le récit} & \textbf{Temps verbal} \\
\hline
\esp{¿Cuándo?} & Situer dans le temps & marqueur temporel \\
\hline
\esp{¿Con quién?} & Présenter les personnages & imparfait (\esp{era}, \esp{tenía}) \\
\hline
\esp{¿Qué hacían las personas?} & Décrire les habitudes & imparfait \\
\hline
\esp{¿Qué pasó?} & Raconter les événements & prétérit indéfini \\
\hline
\esp{¿Qué cambió?} & Bilan, transformation & indéfini \\
\hline
\end{tabularx}
\end{center}

\begin{methode}[Le plan imposé en quatre étapes]
\begin{enumerate}
\item \textbf{Le décor à l'imparfait} (2-3 lignes) : quand, où, avec qui, comment était le village. Exemple : \esp{Cuando era pequeño, iba todos los veranos al pueblo de mis abuelos, cerca de Bafoussam.} « Quand j'étais petit, j'allais tous les étés au village de mes grands-parents, près de Bafoussam. »
\item \textbf{La vie quotidienne à l'imparfait} (2-3 lignes) : ce que faisaient les gens habituellement. Exemple : \esp{Los agricultores trabajaban en el campo y las mujeres vendían la cosecha en el mercado.} « Les agriculteurs travaillaient dans les champs et les femmes vendaient la récolte au marché. »
\item \textbf{L'événement au prétérit indéfini} (2-3 lignes) : ce qui s'est passé une fois, un jour précis. Exemple : \esp{Un día, llegó una moto nueva al pueblo y todos salimos a verla.} « Un jour, une moto neuve est arrivée au village et nous sommes tous sortis la voir. »
\item \textbf{Le changement à l'indéfini} (1-2 lignes) : ce qui a changé. Exemple : \esp{Después, el pueblo cambió mucho: construyeron una carretera y los jóvenes crearon una cooperativa.} « Ensuite, le village a beaucoup changé : ils ont construit une route et les jeunes ont créé une coopérative. »
\end{enumerate}
\end{methode}

\begin{aretenir}[Contraintes obligatoires]
\begin{itemize}
\item Longueur : \cle{8 à 10 lignes}.
\item Au moins \cle{5 mots du lexique} de la leçon : \esp{campo}, \esp{agricultor}, \esp{cosecha}, \esp{éxodo rural}, \esp{migración}, \esp{transporte}, \esp{carretera}, \esp{desarrollo local}, \esp{cooperativa}.
\item Respecter l'équilibre des temps : imparfait pour le décor et les habitudes, indéfini pour les événements.
\end{itemize}
\end{aretenir}

\begin{popfigure}
\begin{center}
\begin{tikzpicture}[
  node distance=0.55cm and 1.1cm,
  etape/.style={rectangle, rounded corners, draw=popblue, fill=popblueL, thick, align=center, minimum width=3.4cm, minimum height=0.9cm, font=\small},
  fleche/.style={-{Stealth[length=3mm]}, thick, popdark}
]
\node[etape, align=center] (n1){1. Marcador temporal\\ \esp{el año pasado, cuando era niño...}};
\node[etape, below=of n1, align=center] (n2){2. Elección del tiempo\\ ¿descripción? $\rightarrow$ imperfecto\\ ¿acontecimiento? $\rightarrow$ indefinido};
\node[etape, below=of n2, align=center] (n3){3. Conjugación\\ persona + terminación correcta};
\node[etape, below=of n3, fill=popgreenL, draw=popgreen, align=center] (n4){4. Frase completa\\ sujeto + verbo + complementos};
\draw[fleche] (n1) -- (n2);
\draw[fleche] (n2) -- (n3);
\draw[fleche] (n3) -- (n4);
\end{tikzpicture}
\end{center}
\legende{Organigramme de la production guidée : du marqueur temporel à la phrase complète.}
\end{popfigure}

\begin{exemplebox}[Deux exemples fondamentaux à imiter]
\esp{Cuando era pequeño, mi abuela cultivaba maíz.} « Quand j'étais petit, ma grand-mère cultivait du maïs. » Ici, deux imparfaits : \esp{era} (décor, âge) et \esp{cultivaba} (habitude régulière).\par\medskip
\esp{El año pasado, mi hermano creó una cooperativa.} « L'année dernière, mon frère a créé une coopérative. » Ici, un indéfini : \esp{creó}, événement unique et daté.
\end{exemplebox}

\begin{attention}[« Il y avait » ou « il y a eu » ?]
C'est le piège classique des francophones. En espagnol, le choix du temps change tout :
\begin{itemize}
\item « Il y \textbf{avait} » (description, état) = \esp{había}, imparfait : \esp{En el pueblo había una escuela pequeña.} « Il y avait une petite école au village. »
\item « Il y \textbf{a eu} » (événement) = \esp{hubo}, indéfini : \esp{Ayer hubo una fiesta en la cooperativa.} « Hier, il y a eu une fête à la coopérative. »
\end{itemize}
Ne traduis jamais « il y avait » par \esp{hubo} ni « il y a eu » par \esp{había}.
\end{attention}

\begin{textebox}[Modèle de production : Mi visita al pueblo]
Cuando era niño, pasaba las vacaciones en el pueblo de mis abuelos, en la región del Oeste. El pueblo era pequeño y tranquilo: había casas de barro, una escuela y un mercado muy animado. Todos los días, los agricultores salían temprano hacia el campo y mi abuela cultivaba maíz y frijoles detrás de la casa.\par
El año pasado volví al pueblo con mi hermano mayor y todo había cambiado. Construyeron una carretera nueva y ahora el transporte es mucho más fácil. Mi hermano, que antes pensaba emigrar a la ciudad, creó una cooperativa con otros jóvenes del pueblo. Juntos vendían la cosecha directamente en Douala y ganaban más dinero.\par
Un día, durante mi visita, hubo una gran reunión en la plaza: el alcalde habló del desarrollo local y de la lucha contra el éxodo rural y la migración. Al final de la fiesta, mi abuela me dijo: « Ahora los jóvenes se quedan aquí ». Yo sonreí y pensé que el pueblo tenía un futuro.
\end{textebox}

\textbf{Glossaire.} \esp{barro} : boue, torchis ; \esp{frijoles} : haricots ; \esp{volver} : revenir ; \esp{ganar dinero} : gagner de l'argent ; \esp{el alcalde} : le maire ; \esp{quedarse} : rester ; \esp{el futuro} : l'avenir.

\textbf{Questions sur le modèle.} 1) Relevez trois verbes à l'imparfait et trois au prétérit indéfini, et justifiez chaque choix. 2) Quels mots du lexique de la leçon l'auteur a-t-il utilisés ? 3) Comment le modèle respecte-t-il le plan en quatre étapes ?

\courssub{Thème : du français vers l'espagnol}

Le thème entraîne la précision : il faut choisir le bon temps, conjuguer sans erreur et employer le lexique exact de la leçon.

\begin{exoresolu}[Thème guidé : deux phrases modèles]
Traduisez : « Quand j'étais petit, ma grand-mère cultivait le maïs. L'année dernière, mon frère a créé une coopérative. »
\tcblower
\textbf{Étape 1 — Analyse des temps français.} « J'étais » et « cultivait » sont à l'imparfait (décor + habitude) $\rightarrow$ imparfait espagnol. « A créé » est un passé composé d'événement unique et daté (« l'année dernière ») $\rightarrow$ prétérit indéfini.\par
\textbf{Étape 2 — Traduction.} \esp{Cuando era pequeño, mi abuela cultivaba maíz. El año pasado, mi hermano creó una cooperativa.}\par
\textbf{Étape 3 — Vérification.} \esp{era} : 1\up{re} personne du singulier de \esp{ser} à l'imparfait ; \esp{cultivaba} : imparfait en \esp{-aba} ; \esp{creó} : indéfini de \esp{crear}, 3\up{e} personne, avec accent obligatoire sur le \esp{-ó}.
\end{exoresolu}

\begin{exercice}[Thème : cinq phrases]
Traduisez en espagnol :
\begin{enumerate}
\item Avant, les jeunes quittaient le village parce qu'il n'y avait pas de travail.
\item Hier, nous avons visité la ferme de mon oncle et il nous a montré la récolte.
\item Quand j'habitais au village, je prenais le bendskin tous les matins pour aller à l'école.
\item L'année dernière, les femmes ont créé une coopérative pour vendre le cacao.
\item Il y avait beaucoup de monde au marché, et soudain il y a eu un grand bruit.
\end{enumerate}
\end{exercice}

\begin{corrige}[Exercice de thème]
\begin{enumerate}
\item \esp{Antes, los jóvenes se iban del pueblo porque no había trabajo.} — Habitude (\esp{se iban}) et description (\esp{había}) : imparfait partout.
\item \esp{Ayer visitamos la finca de mi tío y él nos mostró la cosecha.} — Deux événements datés : indéfini (\esp{visitamos}, \esp{mostró}).
\item \esp{Cuando vivía en el pueblo, tomaba la moto-taxi todas las mañanas para ir a la escuela.} — Habitudes répétées : imparfait (\esp{vivía}, \esp{tomaba}).
\item \esp{El año pasado, las mujeres crearon una cooperativa para vender el cacao.} — Événement daté : indéfini (\esp{crearon}).
\item \esp{Había mucha gente en el mercado, y de repente hubo un gran ruido.} — Description puis événement : \esp{había} puis \esp{hubo}.
\end{enumerate}
\end{corrige}

\courssub{Version : de l'espagnol vers le français}

\begin{methode}[Traduire un verbe au passé : la démarche en trois questions]
\begin{enumerate}
\item \textbf{Identifier le temps espagnol.} Est-ce un imparfait (\esp{-aba}, \esp{-ía}) ou un indéfini (\esp{-é}, \esp{-ó}, \esp{-ieron}) ?
\item \textbf{Choisir l'équivalent français.} Imparfait espagnol $\rightarrow$ imparfait français. Indéfini espagnol $\rightarrow$ passé composé (langue courante) ou passé simple (récit littéraire).
\item \textbf{Vérifier le sens global.} Le contexte confirme-t-il une habitude/description (imparfait) ou un événement (passé composé) ?
\end{enumerate}
\end{methode}

\begin{center}
\begin{tabularx}{\linewidth}{|X|X|X|}
\hline
\rowcolor{popblueL} \textbf{Espagnol} & \textbf{Français} & \textbf{Remarque} \\
\hline
\esp{mi abuela cultivaba maíz} & ma grand-mère cultivait du maïs & imparfait $\rightarrow$ imparfait \\
\hline
\esp{mi hermano creó una cooperativa} & mon frère a créé une coopérative & indéfini $\rightarrow$ passé composé \\
\hline
\esp{había una escuela} & il y avait une école & description \\
\hline
\esp{hubo una reunión} & il y a eu une réunion & événement \\
\hline
\esp{los jóvenes se fueron a la ciudad} & les jeunes sont partis en ville & indéfini de \esp{irse}, irrégulier \\
\hline
\end{tabularx}
\end{center}

\begin{exercice}[Version : cinq phrases extraites du texte support]
Traduisez en français :
\begin{enumerate}
\item \esp{Antes, muchos campesinos trabajaban de sol a sol en los campos de cacao.}
\item \esp{En 2010, el gobierno construyó una carretera que unió el pueblo con la ciudad.}
\item \esp{Cada semana, mi tío llevaba la cosecha al mercado en una moto vieja.}
\item \esp{El éxodo rural dejó los campos casi vacíos y las escuelas perdieron alumnos.}
\item \esp{Gracias a la cooperativa, los agricultores vendían sus productos a mejor precio.}
\end{enumerate}
\end{exercice}

\begin{corrige}[Exercice de version]
\begin{enumerate}
\item « Autrefois, beaucoup de paysans travaillaient du lever au coucher du soleil dans les champs de cacao. » — \esp{trabajaban} : imparfait d'habitude.
\item « En 2010, le gouvernement a construit une route qui a relié le village à la ville. » — \esp{construyó}, \esp{unió} : indéfinis, événements datés.
\item « Chaque semaine, mon oncle portait la récolte au marché sur une vieille moto. » — \esp{llevaba} : imparfait (\esp{cada semana} = répétition).
\item « L'exode rural a laissé les champs presque vides et les écoles ont perdu des élèves. » — \esp{dejó}, \esp{perdieron} : indéfinis, bilan d'événements.
\item « Grâce à la coopérative, les agriculteurs vendaient leurs produits à meilleur prix. » — \esp{vendían} : imparfait, situation durable.
\end{enumerate}
\end{corrige}

\courssub{Correction collective et réécriture individuelle}

La correction collective est un moment d'apprentissage à part entière. Le professeur écrit au tableau deux ou trois productions anonymes d'élèves, et la classe repère ensemble les erreurs selon la grille suivante :

\begin{center}
\begin{tabularx}{\linewidth}{|l|X|X|}
\hline
\rowcolor{poporangeL} \textbf{Type d'erreur} & \textbf{Exemple fautif} & \textbf{Correction} \\
\hline
Aspect verbal & \esp{*Cuando fui pequeño...} & \esp{Cuando era pequeño...} (décor = imparfait) \\
\hline
Aspect verbal & \esp{*Mi abuela cultivó maíz todos los días} & \esp{cultivaba} (habitude = imparfait) \\
\hline
Lexique & \esp{*la cosecha del campo de agricultura} & préférer \esp{la cosecha de maíz} \\
\hline
Conjugaison & \esp{*mi hermano creo una cooperativa} & \esp{creó} (accent de l'indéfini) \\
\hline
Faux ami & \esp{*la carretera estaba larga} & « longue » = \esp{larga}, attention au genre \\
\hline
\end{tabularx}
\end{center}

\begin{attention}[Les trois erreurs reines des francophones]
\begin{enumerate}
\item Mettre l'indéfini partout parce que le français utilise le passé composé : relire chaque phrase et se demander « habitude ou événement ? ».
\item Oublier l'accent de la 3\up{e} personne de l'indéfini : \esp{creó}, \esp{construyó}, \esp{llegó}. Sans accent, \esp{creo} signifie « je crois » !
\item Traduire « il y avait » par \esp{hubo} : c'est \esp{había}.
\end{enumerate}
\end{attention}

\begin{experience}[Réécriture individuelle]
Après la correction collective, chaque élève reprend sa production initiale et la \textbf{réécrit entièrement} au propre, en corrigeant les temps verbaux et le lexique, et en soulignant les 5 mots du lexique de la leçon. La réécriture est ramassée et notée selon trois critères : exactitude des temps (imparfait/indéfini), richesse du lexique rural, respect du plan en quatre étapes. Pour finir, deux volontaires lisent leur texte à voix haute : la classe doit identifier, à l'oreille, les verbes à l'imparfait et ceux à l'indéfini.
\end{experience}

\begin{aretenir}[Bilan de la section]
Produire et traduire au passé exige une seule question-clé : \textbf{description/habitude ou événement ?} Description et habitude $\rightarrow$ imparfait (\esp{era}, \esp{cultivaba}, \esp{había}). Événement unique et daté $\rightarrow$ prétérit indéfini (\esp{creó}, \esp{construyó}, \esp{hubo}). En version, l'imparfait espagnol se rend par l'imparfait français, l'indéfini par le passé composé. Avec la grille des cinq questions, le plan en quatre étapes et le lexique du monde rural, tu es désormais capable de raconter, par écrit comme à l'oral, une expérience vécue au village.
\end{aretenir}

\newpage

\coursec{Activité d'intégration}

\courssub{Objectifs de l'activité}

À l'issue de cette activité d'intégration, tu seras capable de :
\begin{itemize}
\item mobiliser le \cle{lexique du monde rural} (\esp{campo, agricultor, cosecha, éxodo rural, cooperativa, carretera, desarrollo local}...) dans une production authentique ;
\item expliquer les \cle{causes} et les \cle{conséquences} de l'exode rural et de la mobilité ;
\item employer correctement et de manière combinée l'\cle{imparfait} (description du passé) et le \cle{prétérit indéfini} (actions ponctuelles) ;
\item produire oralement et par écrit un \cle{reportage radiophonique} de deux minutes, avec un rôle précis dans un groupe ;
\item écouter attentivement la production des autres groupes et remplir une \cle{grille d'écoute}.
\end{itemize}

\courssub{Situation}

La radio communautaire de ta région, \emph{Radio Mvolyé}, lance un concours scolaire intitulé \esp{« Mi pueblo ayer y hoy »} (« Mon village d'hier et d'aujourd'hui »). Chaque classe de Première doit envoyer un reportage de deux minutes sur un village camerounais, réel ou imaginé. Le meilleur reportage sera diffusé le samedi dans l'émission en espagnol de la radio, très écoutée par les élèves qui apprennent la langue de la Guinée équatoriale voisine.

Voici l'annonce officielle du concours, affichée dans la salle de classe :

\begin{textebox}[Anuncio del concurso — Radio Mvolyé]
¡ATENCIÓN, JÓVENES REPORTEROS!\\
Radio Mvolyé organiza el concurso « Mi pueblo ayer y hoy ».\\
Envíanos un reportaje de dos minutos sobre un pueblo camerunés.\\
Tu reportaje debe incluir :\\
— la descripción del pueblo de antes ;\\
— el testimonio de un joven que se fue a la ciudad ;\\
— una iniciativa de desarrollo local (una cooperativa, una carretera nueva, un proyecto agrícola).\\
¡Usa el pasado correctamente y seis palabras del vocabulario del mundo rural!\\
Fecha límite : el viernes a las 14 horas.\\
¡Buena suerte!
\end{textebox}

\textbf{Traduction de l'annonce :} « Attention, jeunes reporters ! Radio Mvolyé organise le concours "Mon village d'hier et d'aujourd'hui". Envoyez-nous un reportage de deux minutes sur un village camerounais. Votre reportage doit inclure : la description du village d'autrefois ; le témoignage d'un jeune parti à la ville ; une initiative de développement local (une coopérative, une route nouvelle, un projet agricole). Utilisez le passé correctement et six mots du vocabulaire du monde rural ! Date limite : vendredi à 14 heures. Bonne chance ! »

\courssub{Consignes}

Travail en \cle{groupes de 4 élèves}. Chaque membre du groupe joue un rôle obligatoire :

\begin{enumerate}
\item \textbf{Répartition des rôles} (5 minutes) : choisissez qui joue \esp{el periodista} (le journaliste), \esp{el anciano} (l'ancien qui décrit le village d'autrefois, à l'\textbf{imparfait}), \esp{el joven} (le jeune qui raconte son départ et son retour, au \textbf{prétérit indéfini}) et \esp{el representante de la cooperativa} (qui présente une initiative de développement local).
\item \textbf{Choix du village} (5 minutes) : choisissez un village camerounais réel (par exemple un village de la région du Centre, de l'Ouest ou du Nord) ou inventez un village avec un nom, une région et une culture principale (cacao, maïs, manioc...).
\item \textbf{Brainstorming lexical} (10 minutes) : sur une feuille, listez au moins \textbf{six mots du lexique de la leçon} que vous utiliserez : \esp{campo, agricultor, cosecha, éxodo rural, migración, transporte, carretera, desarrollo local, cooperativa}. Soulignez-les dans votre script.
\item \textbf{Rédaction du script} (25 minutes) : rédigez ensemble un script de \textbf{10 à 12 lignes}. Le journaliste pose au moins \textbf{trois questions} (avec ¿...?). L'ancien utilise au moins \textbf{trois verbes à l'imparfait}. Le jeune utilise au moins \textbf{trois verbes au prétérit indéfini}. Le représentant explique \textbf{une cause de l'exode rural et une solution}.
\item \textbf{Répétition} (10 minutes) : répétez le reportage en veillant à la prononciation (le \esp{j} se prononce « r » aspiré, le \esp{ñ} se prononce « gn » comme dans « gagner ») et à l'intonation des questions.
\item \textbf{Représentation} (2 minutes par groupe) : jouez votre reportage devant la classe, comme à la radio, avec une voix claire.
\item \textbf{Grille d'écoute} : pendant les reportages des autres groupes, remplissez la grille ci-dessous.
\end{enumerate}

\begin{center}
\begin{tabularx}{\linewidth}{|X|X|X|}
\hline
\rowcolor{popblueL} Critère d'écoute & Oui / Non & Détails relevés \\
\hline
Causes de l'exode rural mentionnées & & \\
\hline
Solutions de développement local proposées & & \\
\hline
Imparfait correctement employé (3 verbes) & & \\
\hline
Indéfini correctement employé (3 verbes) & & \\
\hline
6 mots du lexique utilisés & & \\
\hline
\end{tabularx}
\end{center}

\courssub{Ressources à mobiliser}

\begin{aretenir}[Boîte à outils de la leçon]
\textbf{Lexique :} \esp{el campo} (la campagne), \esp{el agricultor / la agricultora} (l'agriculteur), \esp{la cosecha} (la récolte), \esp{el éxodo rural} (l'exode rural), \esp{la migración} (la migration), \esp{el transporte} (le transport), \esp{la carretera} (la route), \esp{el desarrollo local} (le développement local), \esp{la cooperativa} (la coopérative).

\textbf{Imparfait} (description, habitudes, décor) : \esp{había} (il y avait), \esp{la gente trabajaba} (les gens travaillaient), \esp{los niños jugaban} (les enfants jouaient), \esp{la carretera era mala} (la route était mauvaise).

\textbf{Prétérit indéfini} (actions ponctuelles, accomplies) : \esp{me fui} (je suis parti), \esp{llegué a la ciudad} (je suis arrivé à la ville), \esp{volví} (je suis revenu), \esp{nació la cooperativa} (la coopérative est née), \esp{construyeron una carretera} (ils ont construit une route).

\textbf{Structures utiles :} \esp{¿Cómo era el pueblo antes?} (Comment était le village avant ?), \esp{¿Por qué se fue usted?} (Pourquoi êtes-vous parti ?), \esp{gracias a la cooperativa} (grâce à la coopérative), \esp{antes... pero ahora...} (avant... mais maintenant...).
\end{aretenir}

\begin{attention}[Pièges à éviter]
\begin{itemize}
\item Ne dis pas \esp{yo fue} : le prétérit indéfini de \esp{ir} à la première personne est \esp{yo fui} (sans accent).
\item \esp{Éxodo} prend un accent sur le \esp{é} ; \esp{migración} sur le \esp{ó}.
\item Attention au faux ami : \esp{la carretera} signifie « la route », pas « la charrette ».
\item Pour décrire le village d'autrefois, n'utilise pas l'indéfini : on dit \esp{el pueblo era pequeño y tranquilo} (imparfait), pas \esp{el pueblo fue pequeño}.
\end{itemize}
\end{attention}

\courssub{Proposition de production et corrigé commenté}

\begin{exoresolu}[Reportage modèle : « Nkolndongo, ayer y hoy »]
Énoncé : rédige et joue un reportage de 10 à 12 lignes respectant toutes les consignes du concours.
\tcblower
\textbf{Script modèle (production attendue) :}

\esp{— Periodista : Buenos días, oyentes de Radio Mvolyé. Hoy estamos en Nkolndongo, un pueblo de la región del Centro. Abuelo Mateo, ¿cómo era el pueblo antes?}

\esp{— Anciano : Antes el pueblo era pequeño y tranquilo. Los agricultores trabajaban en el campo y la cosecha de cacao era buena, pero la carretera era muy mala y no había electricidad.}

\esp{— Periodista : Joven Samuel, ¿por qué se fue usted a la ciudad?}

\esp{— Joven : Me fui a Yaoundé en 2019 porque no había trabajo aquí. Llegué a la ciudad, busqué trabajo y aprendí muchas cosas. Pero el año pasado volví a mi pueblo.}

\esp{— Periodista : ¿Por qué volvió?}

\esp{— Joven : Volví porque nació una cooperativa agrícola. Ahora trabajo con mis padres.}

\esp{— Representante : Sí, la cooperativa compra el cacao de los agricultores a un buen precio. Además, el Estado construyó una carretera nueva. Gracias a este desarrollo local, el éxodo rural disminuyó y los jóvenes regresan.}

\esp{— Periodista : ¡Un ejemplo magnífico! Mi pueblo ayer y hoy : Nkolndongo vuelve a vivir.}

\textbf{Traduction :} « — Journaliste : Bonjour, auditeurs de Radio Mvolyé. Aujourd'hui nous sommes à Nkolndongo, un village de la région du Centre. Grand-père Mateo, comment était le village avant ? — L'ancien : Avant, le village était petit et tranquille. Les agriculteurs travaillaient dans les champs et la récolte de cacao était bonne, mais la route était très mauvaise et il n'y avait pas d'électricité. — Journaliste : Jeune Samuel, pourquoi êtes-vous parti à la ville ? — Le jeune : Je suis parti à Yaoundé en 2019 parce qu'il n'y avait pas de travail ici. Je suis arrivé en ville, j'ai cherché du travail et j'ai appris beaucoup de choses. Mais l'année dernière, je suis revenu dans mon village. — Journaliste : Pourquoi êtes-vous revenu ? — Le jeune : Je suis revenu parce qu'une coopérative agricole est née. Maintenant je travaille avec mes parents. — Le représentant : Oui, la coopérative achète le cacao des agriculteurs à un bon prix. De plus, l'État a construit une route nouvelle. Grâce à ce développement local, l'exode rural a diminué et les jeunes reviennent. — Journaliste : Un exemple magnifique ! Mon village d'hier et d'aujourd'hui : Nkolndongo revit. »

\textbf{Commentaire des choix de langue :}
\begin{itemize}
\item \textbf{Imparfait} pour le décor et les habitudes : \esp{era, trabajaban, era, había} — le village d'autrefois est décrit comme un tableau, sans début ni fin précis.
\item \textbf{Indéfini} pour les actions ponctuelles du récit : \esp{me fui, llegué, busqué, aprendí, volví, nació, construyó, disminuyó} — chaque action est accomplie une fois, datée.
\item \textbf{Lexique obligatoire} : \esp{pueblo, agricultores, campo, cosecha, carretera, cooperativa, desarrollo local, éxodo rural} — huit mots, soit plus que les six exigés.
\item \textbf{Questions bien formées} avec ¿...? : \esp{¿Cómo era...? ¿Por qué se fue usted? ¿Por qué volvió?}
\item \textbf{Cause et conséquence} explicitées : absence de travail (cause de l'exode), coopérative et route (solutions), retour des jeunes (conséquence positive).
\end{itemize}
\end{exoresolu}

\courssub{Bilan de l'activité}

Cette activité t'a permis de mobiliser en situation réelle les trois savoir-faire de la leçon : lire et comprendre un texte sur le monde rural, expliquer les causes et les conséquences de l'exode rural, et raconter au passé une expérience rurale en combinant imparfait et indéfini. Tu as aussi découvert le genre du \cle{reportage radiophonique}, utile pour l'expression orale au baccalauréat. Pour t'auto-évaluer, vérifie : ai-je utilisé six mots du lexique ? ai-je distingué description (imparfait) et actions (indéfini) ? ai-je expliqué au moins une cause et une solution ? Si oui, tu es prêt à devenir reporter de \emph{Radio Mvolyé} !

\newpage

\coursec{Méthodes et exercices résolus}

\coursec{El mundo rural, la movilidad y el desarrollo local}

\courssub{Texte support et découverte}

\begin{textebox}[Vie rurale et développement local : le cas de Batchingou]
Batchingou es un pueblo del Oeste de Camerún, cerca de Dschang. Hace diez años, la vida allí era difícil. Los jóvenes se iban a Douala o Yaoundé porque no había trabajo. Las carreteras eran malas y las cosechas de café y cacao se pudrían en los campos. El pueblo envejecía.

En 2016, los agricultores crearon una cooperativa: «Cooperativa Batchingou Desarrollo». Con la ayuda de una ONG española, construyeron una carretera asfaltada y una unidad de transformación de cacao. Hoy, la cooperativa compra la producción a un precio justo. Los jóvenes vuelven. Pierre, 28 años, ha vuelto de Douala: «Aquí tengo un salario digno y estoy cerca de mi familia.» El turismo rural también empieza: los visitantes descubren las plantaciones y la cultura bamiléké.

El desarrollo local cambia la cara del pueblo. El éxodo rural ya no es una fatalidad.
\end{textebox}

\begin{definition}[Traduction et glossaire des mots clés]
\textbf{Traduction :} « Batchingou est un village de l'Ouest du Cameroun, près de Dschang. Il y a dix ans, la vie y était difficile. Les jeunes partaient à Douala ou Yaoundé parce qu'il n'y avait pas de travail. Les routes étaient mauvaises et les récoltes de café et de cacao pourrissaient dans les champs. Le village vieillissait.

En 2016, les agriculteurs ont créé une coopérative : "Coopérative Batchingou Développement". Avec l'aide d'une ONG espagnole, ils ont construit une route goudronnée et une unité de transformation du cacao. Aujourd'hui, la coopérative achète la production à un prix juste. Les jeunes reviennent. Pierre, 28 ans, est revenu de Douala : "Ici, j'ai un salaire décent et je suis près de ma famille." Le tourisme rural commence aussi : les visiteurs découvrent les plantations et la culture bamiléké.

Le développement local change le visage du village. L'exode rural n'est plus une fatalité. »

\textbf{Glossaire :} \cle{hace diez años} (il y a dix ans) ; \cle{se pudrían} (pourrissaient, imparfait) ; \cle{envejecía} (vieillissait) ; \cle{crearon} (créèrent, indéfini) ; \cle{asfaltada} (goudronnée) ; \cle{unidad de transformación} (unité de transformation) ; \cle{precio justo} (prix juste) ; \cle{salario digno} (salaire décent) ; \cle{turismo rural} (tourisme rural) ; \cle{fatalidad} (fatalité, destin inévitable).
\end{definition}

\courssub{Compréhension du texte}

\begin{methode}[Comprendre un texte sur le monde rural et le développement local]
Pour réussir les questions de compréhension au Baccalauréat, suis cette démarche en cinq étapes :
\begin{enumerate}
\item \textbf{Lis d'abord les questions}, puis le texte deux fois : une lecture globale (idée générale), une lecture attentive (souligne les mots-clés du lexique rural : \esp{campo}, \esp{cosecha}, \esp{éxodo rural}, \esp{cooperativa}, \cle{carretera}, \cle{desarrollo local}).
\item \textbf{Repère la structure du texte} : situation initiale (la vie au village), problème (l'exode rural), solutions (le développement local), résultat (retour des jeunes).
\item \textbf{Localise la réponse} dans le texte : chaque question renvoie à un ou deux paragraphes précis. Recopie le début de la phrase qui contient la réponse, puis reformule avec tes mots.
\item \textbf{Réponds en espagnol, par des phrases complètes}, avec sujet + verbe conjugué. Jamais de réponse en français, jamais d'un seul mot si la question commence par \esp{¿Por qué...?} ou \esp{¿Qué...?}.
\item \textbf{Vérifie} : accord (masculin/féminin, singulier/pluriel), accents, temps du verbe (le passé pour raconter, le présent pour décrire une situation actuelle).
\end{enumerate}
Réflexes : pour une question \esp{¿Por qué?}, commence par \esp{Porque...} ; pour \esp{¿Qué consecuencias tiene?}, utilise \esp{Como consecuencia, ...} ; pour \esp{¿Qué soluciones propone el texto?}, utilise \esp{El texto propone... / Según el texto, ...}.
\end{methode}

\begin{exoresolu}[Compréhension de texte : type Bac (texte court modèle)]
\textbf{Texte} : \esp{En el pueblo de Mbanga, antes casi todos los jóvenes se iban a Douala porque no había trabajo ni carreteras buenas. Pero en 2015, los agricultores crearon una cooperativa de cacao. Gracias a ella, construyeron una carretera nueva y una escuela. Hoy muchos jóvenes vuelven al pueblo porque la cooperativa les ofrece empleo y un futuro en el campo.}

\textbf{Questions} : 1) \esp{¿Por qué se iban los jóvenes de Mbanga?} 2) \esp{¿Qué crearon los agricultores en 2015?} 3) \esp{¿Qué consecuencias tuvo la cooperativa?} 4) \esp{¿Por qué vuelven hoy los jóvenes?}
\tcblower
\textbf{Raisonnement} : chaque question vise un segment précis du texte. On repère les mots-clés de la question (\esp{jóvenes}, \esp{2015}, \esp{cooperativa}, \esp{vuelven}) et on reformule la phrase correspondante.

\textbf{Réponses rédigées} :
\begin{enumerate}
\item \esp{Los jóvenes se iban de Mbanga porque no había trabajo ni carreteras buenas en el pueblo.} (Question avec \esp{¿Por qué?} $\rightarrow$ réponse avec \esp{porque} ; on garde l'imparfait \esp{había} car le texte décrit une situation passée.)
\item \esp{En 2015, los agricultores crearon una cooperativa de cacao.} (Prétérit indéfini \esp{crearon} : action ponctuelle, datée.)
\item \esp{Como consecuencia, construyeron una carretera nueva y una escuela, y la cooperativa ofrece empleo a los jóvenes.}
\item \esp{Hoy muchos jóvenes vuelven al pueblo porque la cooperativa les offre empleo y un futuro en el campo.} (Présent \esp{vuelven}, \esp{ofrece} : situation actuelle, signalée par \esp{hoy}.)
\end{enumerate}
\textbf{Conclusion méthodologique} : une bonne réponse reprend les termes de la question, cite l'information du texte et respecte le temps verbal du passage concerné.
\end{exoresolu}

\begin{exercice}[Questions sur le texte support « Batchingou »]
Répondez en espagnol, par des phrases complètes.
\begin{enumerate}
\item ¿Dónde está situado Batchingou?
\item ¿Cómo era la vida en Batchingou hace diez años? (Deux détails minimum)
\item ¿Por qué se iban los jóvenes?
\item ¿Qué crearon los agricultores en 2016?
\item ¿Qué dos infraestructuras construyeron gracias a la ONG?
\item ¿Qué hace la cooperativa hoy con la producción de cacao?
\item ¿Por qué ha vuelto Pierre de Douala? (Citez deux raisons)
\item ¿Qué nueva actividad económica aparece en el pueblo?
\item Según el texto, ¿el éxodo rural es inevitable? Justifiez.
\end{enumerate}
\end{exercice}

\courssub{Lexique thématique : Le monde rural}

\begin{aretenir}[Vocabulaire essentiel du module]
\begin{center}
\begin{tabularx}{\linewidth}{|X|X|X|}
\hline
\rowcolor{popblueL} \textbf{Español} & \textbf{Français} & \textbf{Exemple / Contexte} \\
\hline
\cle{el campo} & la campagne & \esp{Mi abuelo vive en el campo.} \\
\hline
\cle{el agricultor / la agricultora} & l'agriculteur & \esp{Los agricultores siembran maíz.} \\
\hline
\cle{la cosecha} & la récolte & \esp{La cosecha de cacao fue abundante.} \\
\hline
\cle{el éxodo rural} & l'exode rural & \esp{El éxodo rural vacía los pueblos.} \\
\hline
\cle{la migración} & la migration & \cle{la migración} campo-ciudad. \\
\hline
\cle{el transporte} & le transport & \esp{El transporte rural es deficiente.} \\
\hline
\cle{la carretera} & la route & \esp{Construyeron una carretera nueva.} \\
\hline
\cle{el desarrollo local} & le développement local & \esp{El desarrollo local frena el éxodo.} \\
\hline
\cle{la cooperativa} & la coopérative & \esp{La cooperativa ayuda a los agricultores.} \\
\hline
\end{tabularx}
\end{center}
\textbf{Autres mots utiles :} \cle{la siembra} (le semis), \cle{el cultivo} (la culture), \cle{el maíz} (le maïs), \cle{los frijoles / las judías} (les haricots), \cle{la ONG} (l'ONG), \cle{la unidad de transformación} (l'unité de transformation), \cle{el precio justo} (le prix juste), \cle{el salario digno} (le salaire décent), \cle{el turismo rural} (le tourisme rural), \cle{envejecer} (vieillir), \cle{pudrirse} (pourrir).
\end{aretenir}

\courssub{Grammaire : Prétérit indéfini et imparfait}

\begin{propriete}[Formation : Prétérit indéfini (Pretérito indefinido)]
\textbf{Verbes réguliers} (terminaisons accentuées aux 1\up{re} et 3\up{e} pers. sg) :
\begin{center}
\begin{tabular}{|l|c|c|c|c|c|c|}
\hline
\rowcolor{popblueL} & \textbf{yo} & \textbf{tú} & \textbf{él/ella} & \textbf{nosotros} & \textbf{vosotros} & \textbf{ellos} \\
\hline
\textbf{-AR} (hablar) & hablé & hablaste & habló & hablamos & hablasteis & hablaron \\
\hline
\textbf{-ER} (comer) & comí & comiste & comió & comimos & comisteis & comieron \\
\hline
\textbf{-IR} (vivir) & viví & viviste & vivió & vivimos & vivisteis & vivieron \\
\hline
\end{tabular}
\end{center}
\textbf{Verbes irréguliers fréquents (radical modifié + terminaisons spécifiques sans accent)} :
\begin{center}
\begin{tabular}{|l|c|c|c|c|c|c|}
\hline
\rowcolor{popblueL} & \textbf{yo} & \textbf{tú} & \textbf{él/ella} & \textbf{nosotros} & \textbf{vosotros} & \textbf{ellos} \\
\hline
\textbf{ser / ir} & fui & fuiste & fue & fuimos & fuisteis & fueron \\
\hline
\textbf{hacer} & hice & hiciste & hizo & hicimos & hicisteis & hicieron \\
\hline
\textbf{tener} & tuve & tuviste & tuvo & tuvimos & tuvisteis & tuvieron \\
\hline
\textbf{estar} & estuve & estuviste & estuvo & estuvimos & estuvisteis & estuvieron \\
\hline
\textbf{poder} & pude & pudiste & pudo & pudimos & pudisteis & pudieron \\
\hline
\textbf{decir} & dije & dijiste & dijo & dijimos & dijisteis & dijeron \\
\hline
\textbf{ver} & vi & viste & vio & vimos & visteis & vieron \\
\hline
\textbf{dar} & di & diste & dio & dimos & disteis & dieron \\
\hline
\end{tabular}
\end{center}
\emph{Attention :} \esp{hizo} (3\up{e} sg) change \cle{c $\rightarrow$ z} pour garder le son [k] ; \esp{dijo} (3\up{e} sg) change \cle{i $\rightarrow$ o}.
\end{propriete}

\begin{propriete}[Formation : Imparfait (Pretérito imperfecto)]
\textbf{Verbes réguliers} (seuls 3 verbes irréguliers : \esp{ser}, \esp{ir}, \esp{ver}) :
\begin{center}
\begin{tabular}{|l|c|c|c|c|c|c|}
\hline
\rowcolor{popblueL} & \textbf{yo} & \textbf{tú} & \textbf{él/ella} & \textbf{nosotros} & \textbf{vosotros} & \textbf{ellos} \\
\hline
\textbf{-AR} (hablar) & hablaba & hablabas & hablaba & \textbf{hablábamos} & hablabais & hablaban \\
\hline
\textbf{-ER} (comer) & comía & comías & comía & \textbf{comíamos} & comíais & comían \\
\hline
\textbf{-IR} (vivir) & vivía & vivías & vivía & \textbf{vivíamos} & vivíais & vivían \\
\hline
\end{tabular}
\end{center}
\textbf{Irréguliers} :
\begin{center}
\begin{tabular}{|l|c|c|c|c|c|c|}
\hline
\rowcolor{popblueL} & \textbf{yo} & \textbf{tú} & \textbf{él/ella} & \textbf{nosotros} & \textbf{vosotros} & \textbf{ellos} \\
\hline
\textbf{ser} & era & eras & era & \textbf{éramos} & erais & eran \\
\hline
\textbf{ir} & iba & ibas & iba & \textbf{íbamos} & ibais & iban \\
\hline
\textbf{ver} & veía & veías & veía & \textbf{veíamos} & veíais & veían \\
\hline
\end{tabular}
\end{center}
\emph{Note :} À la 1\up{re} pers. pl., l'accent tombe sur le \cle{á} (-AR), \cle{í} (-ER/-IR), \cle{é} (ser), \cle{í} (ir), \cle{í} (ver).
\end{propriete}

\begin{propriete}[Emploi : Choisir entre Indéfini et Imparfait]
\begin{itemize}
\item \textbf{Indéfini} = \textbf{Action ponctuelle, achevée, datée, qui fait avancer le récit}.
  \begin{itemize}
  \item Repères : \cle{ayer}, \cle{anoche}, \cle{el año pasado}, \cle{en 2016}, \cle{un día}, \cle{suddenly} (\esp{de repente}), \cle{primero}, \cle{luego}, \cle{al final}.
  \item Ex. : \esp{En 2016, los agricultores \textbf{crearon} una cooperativa.}
  \end{itemize}
\item \textbf{Imparfait} = \textbf{Description, décor, habitude, action en cours, état mental/physique dans le passé}.
  \begin{itemize}
  \item Repères : \cle{antes}, \cle{antes de}, \cle{cada día}, \cle{todos los días}, \cle{mientras}, \cle{cuando} (si description), \cle{era}, \cle{había}, \cle{tenía}, \cle{estaba}, \cle{quería}, \cle{sabía}.
  \item Ex. : \esp{Antes, la vida \textbf{era} difícil. \textbf{Había} mucha pobreza.}
  \end{itemize}
\item \textbf{Combinaison classique} : \esp{Mientras \textbf{trabajábamos} (imparfait, fond), \textbf{llegó} (indéfini, événement) la lluvia.}
\end{itemize}
\end{propriete}

\begin{attention}[Verbes à changement de voyelle (Présent) vs Passés]
Les verbes comme \cle{volver} (o$\rightarrow$ue), \cle{poder} (o$\rightarrow$ue), \cle{dormir} (o$\rightarrow$ue), \cle{pensar} (e$\rightarrow$ie), \cle{querer} (e$\rightarrow$ie), \cle{pedir} (e$\rightarrow$i) changent de radical \textbf{uniquement au présent} (et au subjonctif).
\begin{itemize}
\item Présent : \esp{vuelvo, vuelves, vuelve, volvemos, volvéis, vuelven}.
\item Indéfini : \esp{volví, volviste, volvió, volvimos, volvisteis, volvieron} (radical \emph{volv-}).
\item Imparfait : \esp{volvía, volvías, volvía, volvíamos, volvíais, volvían} (radical \emph{volv-}).
\end{itemize}
\textbf{Erreur fatale :} écrire \esp{*vuelvían} à l'imparfait.
\end{attention}

\courssub{Savoir-faire 1 : Lire et comprendre un texte sur le monde rural}

\begin{methode}[Comprendre un texte sur le monde rural et le développement local]
Pour réussir les questions de compréhension au Baccalauréat, suis cette démarche en cinq étapes :
\begin{enumerate}
\item \textbf{Lis d'abord les questions}, puis le texte deux fois : une lecture globale (idée générale), une lecture attentive (souligne les mots-clés du lexique rural : \esp{campo}, \esp{cosecha}, \esp{éxodo rural}, \esp{cooperativa}, \cle{carretera}, \cle{desarrollo local}).
\item \textbf{Repère la structure du texte} : situation initiale (la vie au village), problème (l'exode rural), solutions (le développement local), résultat (retour des jeunes).
\item \textbf{Localise la réponse} dans le texte : chaque question renvoie à un ou deux paragraphes précis. Recopie le début de la phrase qui contient la réponse, puis reformule avec tes mots.
\item \textbf{Réponds en espagnol, par des phrases complètes}, avec sujet + verbe conjugué. Jamais de réponse en français, jamais d'un seul mot si la question commence par \esp{¿Por qué...?} ou \esp{¿Qué...?}.
\item \textbf{Vérifie} : accord (masculin/féminin, singulier/pluriel), accents, temps du verbe (le passé pour raconter, le présent pour décrire une situation actuelle).
\end{enumerate}
Réflexes : pour une question \esp{¿Por qué?}, commence par \esp{Porque...} ; pour \esp{¿Qué consecuencias tiene?}, utilise \esp{Como consecuencia, ...} ; pour \esp{¿Qué soluciones propone el texto?}, utilise \esp{El texto propone... / Según el texto, ...}.
\end{methode}

\courssub{Savoir-faire 2 : Expliquer les causes et les conséquences de l'exode rural}

\begin{methode}[Exprimer la cause et la conséquence]
Pour expliquer un phénomène social comme l'exode rural, il faut maîtriser les connecteurs logiques :
\begin{enumerate}
\item \textbf{Identifier} dans le texte les causes (le « pourquoi ») et les conséquences (le « résultat »).
\item \textbf{Exprimer la cause} : \esp{porque} + phrase ; \esp{por} / \esp{a causa de} / \esp{debido a} + nom (\esp{debido a la falta de trabajo}).
\item \textbf{Exprimer la conséquence} : \esp{por eso}, \esp{entonces}, \esp{como consecuencia}, \esp{por lo tanto}, \esp{así que} + phrase.
\item \textbf{Enchaîner} au moins deux causes et deux conséquences avec \esp{además}, \esp{también}, \esp{sin embargo} pour nuancer.
\item \textbf{Conclure} par une solution possible : \esp{Para resolver este problema, hay que... / es necesario...} + infinitif.
\end{enumerate}
Structures à retenir : \esp{La falta de... provoca...} ; \esp{El éxodo rural tiene consecuencias graves: ...} ; \esp{Gracias a la cooperativa, el pueblo se desarrolla.}
\end{methode}

\begin{exoresolu}[Expression : causes et conséquences]
\textbf{Énoncé} : À partir du texte sur Batchingou, rédigez en espagnol un petit paragraphe (6 à 8 lignes) intitulé \esp{El éxodo rural: causas y consecuencias}. Utilisez au moins deux connecteurs de cause et deux connecteurs de conséquence.
\tcblower
\textbf{Raisonnement} : on classe d'abord les idées. Causes : pas de travail, mauvaises routes, pourriture des récoltes, vieillissement. Conséquences : les jeunes partent, le village se vide, perte de savoir-faire. Solutions : coopérative, route, transformation locale, tourisme. On choisit ensuite les connecteurs : \esp{porque}, \esp{a causa de}, \esp{por eso}, \esp{como consecuencia}, \esp{gracias a}.

\textbf{Production modèle} :

\esp{El éxodo rural es un problema grave en muchos pueblos de Camerún. Los jóvenes se van a las grandes ciudades como Douala porque no hay empleo en el campo y a causa de las malas carreteras. Además, las cosechas se pierden por falta de transformación. Como consecuencia, los pueblos se vacían y la población envejece. Por eso, es necesario crear cooperativas y mejorar el transporte. Gracias a estas iniciativas, como en Batchingou, los jóvenes vuelven y el mundo rural tiene futuro.}

« L'exode rural est un problème grave dans beaucoup de villages du Cameroun. Les jeunes partent vers les grandes villes comme Douala parce qu'il n'y a pas d'emploi à la campagne et à cause des mauvaises routes. De plus, les récoltes se perdent par manque de transformation. En conséquence, les villages se vident et la population vieillit. C'est pourquoi il faut créer des coopératives et améliorer le transport. Grâce à ces initiatives, comme à Batchingou, les jeunes reviennent et le monde rural a un avenir. »

\textbf{Conclusion} : le paragraphe suit le plan problème $\rightarrow$ causes $\rightarrow$ conséquences $\rightarrow$ solutions, avec cinq connecteurs variés (\esp{porque}, \esp{a causa de}, \esp{además}, \esp{como consecuencia}, \esp{por eso}, \esp{gracias a}) et des verbes correctement conjugués au présent (vérité générale).
\end{exoresolu}

\courssub{Savoir-faire 3 : Raconter au passé une expérience rurale}

\begin{methode}[Choisir entre prétérit indéfini et imparfait]
Pour raconter un souvenir (des vacances au village, une récolte), pose-toi deux questions :
\begin{enumerate}
\item \textbf{Décor, description, habitude, situation en arrière-plan} $\rightarrow$ \textbf{imparfait} : \esp{El pueblo era pequeño. Hacía calor. Todos los días íbamos al río.}
\item \textbf{Action ponctuelle, datée, qui fait avancer le récit} $\rightarrow$ \textbf{prétérit indéfini} : \esp{En 2019 visité a mis abuelos. Ayudé en la cosecha. Volví a Yaoundé en septiembre.}
\item \textbf{Combinaison} : l'imparfait décrit le fond, l'indéfini donne les événements : \esp{Mientras trabajábamos en el campo (fond), llegó la lluvia (événement).}
\item \textbf{Repères chronologiques} : \esp{primero}, \esp{luego}, \esp{después}, \esp{al final}, \esp{aquella noche}, \esp{de repente}.
\item \textbf{Attention aux irréguliers} : \esp{hacer $\rightarrow$ hice} ; \esp{ir/ser $\rightarrow$ fui} ; \esp{tener $\rightarrow$ tuve} ; \esp{estar $\rightarrow$ estuve} ; \esp{poder $\rightarrow$ pude} ; \esp{decir $\rightarrow$ dije} ; \esp{ver $\rightarrow$ vi} ; \esp{dar $\rightarrow$ di}.
\end{enumerate}
\end{methode}

\begin{exoresolu}[Rédaction : un souvenir au village]
\textbf{Énoncé} : Racontez en espagnol (8 à 10 lignes) des vacances passées dans un village camerounais. Commencez par : \esp{El verano pasado pasé dos semanas en el pueblo de mis abuelos...}
\tcblower
\textbf{Raisonnement} : on construit le récit en alternant décor (imparfait) et actions (indéfini). Plan : arrivée (indéfini), description du village (imparfait), activités quotidiennes (imparfait d'habitude), un événement marquant (indéfini), bilan (indéfini + impression).

\textbf{Production modèle} :

\esp{El verano pasado pasé dos semanas en el pueblo de mis abuelos, cerca de Bafoussam. Cuando llegué, el pueblo era tranquilo y las casas eran pequeñas. Todos los días me levantaba temprano y ayudaba a mi abuelo en el campo: él cultivaba maíz y frijoles. Una mañana, mientras trabajábamos, empezó a llover muy fuerte y corrimos a la casa. Aquella noche, toda la familia cenó junta y mi abuela contó historias del pasado. Al final de las vacaciones, volví a Yaoundé muy contento: aprendí mucho sobre la vida rural y entendí el valor del trabajo de los agricultores.}

« L'été dernier, j'ai passé deux semaines dans le village de mes grands-parents, près de Bafoussam. Quand je suis arrivé, le village était tranquille et les maisons étaient petites. Tous les jours, je me levais tôt et j'aidais mon grand-père aux champs : il cultivait du maïs et des haricots. Un matin, pendant que nous travaillions, il s'est mis à pleuvoir très fort et nous avons couru jusqu'à la maison. Ce soir-là, toute la famille a dîné ensemble et ma grand-mère a raconté des histoires du passé. À la fin des vacances, je suis rentré à Yaoundé très content : j'ai beaucoup appris sur la vie rurale et j'ai compris la valeur du travail des agriculteurs. »

\textbf{Justifications grammaticales} : \esp{era}, \esp{eran}, \esp{me levantaba}, \esp{ayudaba}, \esp{cultivaba}, \esp{trabajábamos} = imparfait (description et habitudes) ; \esp{pasé}, \esp{llegué}, \esp{empezó}, \esp{corrimos}, \esp{cenó}, \esp{contó}, \esp{volví}, \esp{aprendí}, \esp{entendí} = indéfini (actions ponctuelles). Noter \esp{cuando llegué} : l'arrivée est un événement unique, donc indéfini, et non imparfait.
\end{exoresolu}

\courssub{Méthode du commentaire de texte}

\begin{methode}[Structure du commentaire composé (Bac LV2)]
Le commentaire de texte suit un plan rigoureux en trois parties. Tu disposes de 1h30 à 2h.
\begin{enumerate}
\item \textbf{Introduction} (1 paragraphe) :
  \begin{itemize}
  \item \textbf{Accroche} : phrase d'ouverture sur le thème (monde rural, exode, développement).
  \item \textbf{Présentation du texte} : nature (article, témoignage, extrait de roman), auteur (si connu), source, date, titre, thème général.
  \item \textbf{Problématique} : question à laquelle le texte répond (ex. : \emph{Comment le développement local peut-il inverser la dynamique de l'exode rural ?}).
  \item \textbf{Annonce du plan} : \esp{En primer lugar... En segundo lugar... Finalmente...}
  \end{itemize}
\item \textbf{Développement} (2 ou 3 parties, 1 paragraphe chacune) :
  \begin{itemize}
  \item Chaque partie = une idée force + \textbf{citation courte} intégrée + \textbf{analyse} (explication, lien avec le thème, vocabulaire clé).
  \item Partie 1 : Le constat initial (l'exode, les causes). Partie 2 : L'initiative locale (la coopérative, la route). Partie 3 : Les résultats (retour, espoir, tourisme).
  \item Connecteurs : \esp{En primer lugar}, \esp{Por un lado... Por otro lado}, \esp{Además}, \esp{Sin embargo}, \esp{En conclusión}.
  \end{itemize}
\item \textbf{Conclusion} (1 paragraphe) :
  \begin{itemize}
  \item \textbf{Bilan} : réponse synthétique à la problématique.
  \item \textbf{Ouverture} : lien avec l'actualité (autres pays, Guinée équatoriale, rôle des ONG, jeunesse africaine).
  \end{itemize}
\end{enumerate}
\textbf{Règles d'or} : Écrire en espagnol (ou français selon consigne), phrases complètes, pas de paraphrase sans analyse, citations entre guillemets \esp{« ... »}.
\end{methode}

\begin{exoresolu}[Commentaire modèle sur « El renacimiento de Batchingou »]
\textbf{Problématique} : \emph{¿Cómo puede el desarrollo local revertir la dinámica del éxodo rural?}

\textbf{Introducción}
El texto « El renacimiento de Batchingou » presenta la transformación de un pueblo camerunés gracias a una iniciativa cooperativa. El tema central es la lucha contra el éxodo rural mediante el desarrollo local. La pregunta que guía nuestro análisis es: ¿cómo puede el desarrollo local revertir la dinámica del éxodo rural? Veremos primero el diagnóstico inicial, luego la solución cooperativa y finalmente los resultados esperanzadores.

\textbf{Desarrollo}
\textbf{1. Un pueblo amenazado por el éxodo.} En el primer párrafo, el autor describe una situación crítica: «la vida allí era difícil», «los jóvenes se iban», «las carreteras eran malas», «el pueblo envejecía». El uso del imperfecto (\esp{era}, \esp{eran}, \esp{envejecía}) instala un décor de abandono estructural. La causa principal es económica: \esp{no había trabajo}. La consecuencia demográfica es grave: la pérdida de la fuerza viva.

\textbf{2. La cooperativa como motor de cambio.} El segundo párrafo narra el punto de inflexión (indéfini: \esp{crearon}, \esp{construyeron}). La creación de la «Cooperativa Batchingou Desarrollo» en 2016, apoyada por una ONG española, trae infraestructuras concretas: «una carretera asfaltada» y «una unidad de transformación». El término \esp{precio justo} evoca el comercio justo, clave para la renta de los agricultores.

\textbf{3. El retorno de la juventud y nuevas perspectivas.} El tercer párrafo muestra los efectos en presente (\esp{vuelven}, \esp{tiene}, \esp{empieza}). El testimonio de Pierre (\esp{salario digno}, \esp{cerca de mi familia}) humaniza la estadística. Aparece el \esp{turismo rural}, diversificando la economía. La frase final, \esp{El éxodo rural ya no es una fatalidad}, resume la tesis: la fatalidad se vence con organización endógena y apoyo externo.

\textbf{Conclusión}
El texto demuestra que el éxodo rural no es irreversible. La clave está en la asociación de los productores (cooperativa), la infraestructura (carretera) y la valorización in situ (transformación, turismo). Este modelo de Batchingou, replicable en otras zonas de Camerún o de Guinea Ecuatorial, ofrece una vía africana al desarrollo sostenible, donde la juventud deja de ser «fuerza de partida» para ser «fuerza de quedada».
\end{exoresolu}

\courssub{Savoir-faire 4 : Traduire (thème) sur le monde rural}

\begin{methode}[Réussir le thème : du français vers l'espagnol]
\begin{enumerate}
\item \textbf{Identifier le temps} de la phrase française : passé composé ou événement unique $\rightarrow$ indéfini ; imparfait ou habitude $\rightarrow$ imparfait ; présent $\rightarrow$ présent.
\item \textbf{Traduire le vocabulaire du thème} : la campagne = \esp{el campo} ; un agriculteur = \esp{un agricultor} ; la récolte = \esp{la cosecha} ; l'exode rural = \esp{el éxodo rural} ; une route = \esp{una carretera} ; le développement local = \esp{el desarrollo local}.
\item \textbf{Construire la phrase espagnole} : sujet + verbe conjugué + compléments ; placer les adjectifs après le nom (\esp{una carretera nueva}) sauf valeur expressive (\esp{la nueva carretera}).
\item \textbf{Attention aux prépositions} : aller à la ville = \esp{ir a la ciudad} ; grâce à = \esp{gracias a} ; à cause de = \esp{a causa de} / \esp{debido a}.
\item \textbf{Relire} : accents (\esp{éxodo}, \esp{migración}, \esp{carretera}), ñ (\esp{España}, \esp{cañón}), accords (\esp{las carreteras buenas}).
\end{enumerate}
\end{methode}

\begin{exoresolu}[Thème guidé]
\textbf{Énoncé} : Traduisez en espagnol :
1) « Avant, beaucoup de jeunes quittaient le village parce qu'il n'y avait pas de travail. »
2) « En 2018, les agriculteurs ont créé une coopérative de cacao. »
3) « Grâce à la nouvelle route, le village se développe. »
4) « Tous les ans, nous allions vendre la récolte à la ville. »
5) « Un jour, mon père a acheté un camion et notre vie a changé. »
\tcblower
\textbf{Raisonnement et solutions} :
\begin{enumerate}
\item « Avant » + habitude passée (« quittaient ») et description (« il n'y avait pas ») $\rightarrow$ imparfait partout : \esp{Antes, muchos jóvenes se iban del pueblo porque no había trabajo.} Attention : « quitter » se traduit ici par \esp{irse de} (partir de), pas par un calque de « quitter ».
\item Action datée (2018) $\rightarrow$ indéfini : \esp{En 2018, los agricultores crearon una cooperativa de cacao.} \esp{Crear} est régulier : \esp{crearon} sans accent sur la troisième personne du pluriel.
\item Situation actuelle $\rightarrow$ présent : \esp{Gracias a la nueva carretera, el pueblo se desarrolla.} L'adjectif \esp{nueva} est placé avant le nom ici (ordre expressif admis) ; on peut aussi écrire \esp{la carretera nueva}. \esp{Desarrollarse} est pronominal : ne pas oublier \esp{se}.
\item Habitude passée (\esp{tous les ans}) $\rightarrow$ imparfait : \esp{Todos los años íbamos a la ciudad para vender la cosecha.} \esp{Ir} $\rightarrow$ \esp{íbamos} (irrégulier imparfait).
\item Actions ponctuelles successives $\rightarrow$ indéfini : \esp{Un día, mi padre compró un camión y nuestra vida cambió.} \esp{Comprar} et \esp{cambiar} réguliers.
\end{enumerate}
\textbf{Conclusion} : au thème, le choix du temps et la préposition valent autant de points que le vocabulaire.
\end{exoresolu}

\courssub{Savoir-faire 5 : Traduire (version) et transformer les temps}

\begin{methode}[Réussir la version : de l'espagnol vers le français]
\begin{enumerate}
\item \textbf{Repérer les temps espagnols} : indéfini $\rightarrow$ passé composé ou passé simple français ; imparfait $\rightarrow$ imparfait français.
\item \textbf{Traduire mot à mot d'abord, puis reformuler} en français naturel, sans calquer l'ordre espagnol (\esp{una carretera nueva} = « une nouvelle route », adjectif antéposé en français).
\item \textbf{Méfie-toi des faux amis} : \esp{éxito} = « succès » (pas « exode ») ; \esp{éxodo} = « exode » ; \esp{actual} = « actuel » (pas « réel ») ; \esp{asistir a} = « assister à / aller à » ; \esp{realizar} = « réaliser / accomplir » (pas « réaliser » au sens de comprendre).
\item \textbf{Pour les transformations de temps} (exercice type « mettez le verbe au temps demandé ») : conjugue d'abord au présent mental, puis applique la terminaison du temps visé et vérifie l'irrégularité éventuelle.
\end{enumerate}
\end{methode}

\begin{exoresolu}[Version et transformation]
\textbf{Énoncé} : A) Traduisez en français : \esp{Cuando era niño, mi familia vivía en un pueblo. Todos los años, en diciembre, íbamos a la ciudad para vender la cosecha. Un día, mi padre compró un camión y nuestra vida cambió.} B) Transformez : mettez \esp{los jóvenes (volver) al pueblo} d'abord au présent, puis au prétérit indéfini, puis à l'imparfait.
\tcblower
\textbf{A) Raisonnement} : \esp{era}, \esp{vivía}, \esp{íbamos} sont des imparfaits (habitude et description) $\rightarrow$ imparfait français ; \esp{compró}, \esp{cambió} sont des indéfinis (événements) $\rightarrow$ passé composé.

\textbf{Traduction} : « Quand j'étais enfant, ma famille vivait dans un village. Tous les ans, en décembre, nous allions à la ville pour vendre la récolte. Un jour, mon père a acheté un camion et notre vie a changé. »

\textbf{B) Transformations} :
\begin{center}
\begin{tabular}{|l|l|l|}
\hline
\rowcolor{popblueL} Temps & Forme & Traduction \\
\hline
Presente & \esp{los jóvenes vuelven} & les jeunes reviennent \\
\hline
Pretérito indefinido & \esp{los jóvenes volvieron} & les jeunes sont revenus / revinrent \\
\hline
Pretérito imperfecto & \esp{los jóvenes volvían} & les jeunes revenaient \\
\hline
\end{tabular}
\end{center}
\textbf{Justification} : \esp{volver} est un verbe à changement de voyelle (o $\rightarrow$ ue) au présent (\esp{vuelven}), mais ce changement disparaît à l'indéfini (\esp{volvieron}) et à l'imparfait (\esp{volvían}), qui se conjuguent sur le radical de l'infinitif. Erreur fréquente : écrire \esp{*vuelvían} à l'imparfait.
\end{exoresolu}

\courssub{Entraînement : Exercices types Bac}

\begin{exercice}[Exercice 1 : Conjugaison — Indéfini ou Imparfait ?]
Mettez les verbes entre parenthèses au temps qui convient (indéfini ou imparfait).
\begin{enumerate}
\item Cuando yo \esp{(ser)} niño, mi abuelo \esp{(tener)} una finca grande.
\item Todos los veranos nosotros \esp{(ir)} al campo para \esp{(ayudar)} en la cosecha.
\item En 2015, una ONG \esp{(construir)} una carretera nueva.
\item \esp{(Hacer)} buen tiempo aquel día y los agricultores \esp{(poder)} trabajar bien.
\item De repente, \esp{(empezar)} a llover y todos \esp{(correr)} a la casa.
\item Antes, no \esp{(haber)} electricidad en el pueblo.
\item Mi abuela \esp{(contar)} historias mientras nosotros \esp{(cenar)}.
\item Al final, yo \esp{(volver)} a Yaoundé con muchos recuerdos.
\end{enumerate}
\end{exercice}

\begin{exercice}[Exercice 2 : Thème (Français $\rightarrow$ Espagnol)]
Traduisez en espagnol.
\begin{enumerate}
\item L'exode rural vide les villages. (Utilisez \esp{vaciar})
\item À cause du manque de routes, les récoltes pourrissent. (\esp{pudrirse})
\item En 2020, les jeunes ont créé une association pour le développement local.
\item Grâce à la coopérative, les agriculteurs vendent leur cacao à un prix juste.
\item Avant, il n'y avait pas d'école dans le village. (\esp{había})
\item Pierre est revenu de la ville l'année dernière. (\esp{volver})
\item Maintenant, il travaille dans l'unité de transformation. (\esp{trabajar})
\item Le tourisme rural commence à se développer. (\esp{desarrollarse})
\end{enumerate}
\end{exercice}

\begin{exercice}[Exercice 3 : Version (Espagnol $\rightarrow$ Français)]
Traduisez en français le paragraphe suivant :
\esp{« La cooperativa no solo compra el cacao. También organiza talleres de formación para los agricultores. Antes, ellos no sabían cómo mejorar la calidad. Ahora, gracias a los técnicos, la producción ha aumentado. Los jóvenes ven que hay futuro y no se van. Es un cambio radical para la comunidad. »}

\textit{Note :} \esp{ha aumentado} = passé composé (pretérito perfecto compuesto) $\rightarrow$ « a augmenté ».
\end{exercice}

\begin{exercice}[Exercice 4 : Expression écrite — Causes et conséquences]
Rédigez en espagnol un paragraphe d'une dizaine de lignes intitulé \esp{La movilidad rural en mi región}. Expliquez pourquoi les jeunes quittent votre région (causes) et ce qui se passe ensuite (conséquences). Proposez une solution. Utilisez au moins 4 connecteurs logiques (\esp{porque, a causa de, por eso, como consecuencia, además, sin embargo, gracias a}).
\end{exercice}

\begin{exercice}[Exercice 5 : Expression écrite — Récit de vacances]
Rédigez en espagnol (10-12 lignes) un récit intitulé \esp{Mis vacaciones en el campo}. Décrivez le village (imparfait), racontez vos activités quotidiennes (imparfait), puis un événement marquant (indéfini) et votre retour (indéfini). Utilisez au moins 5 verbes à l'indéfini et 5 à l'imparfait. Commencez par : \esp{El año pasado pasé quince días en el pueblo de...}
\end{exercice}

\courssub{Corrigés des exercices d'entraînement}

\begin{corrige}[Exercice 1 : Conjugaison]
\begin{enumerate}
\item \esp{era} (imparfait, description état), \esp{tenía} (imparfait, possession passé).
\item \esp{íbamos} (imparfait, habitude \emph{todos los veranos}), \esp{ayudar} (infinitif après \esp{ir a} ou but).
\item \esp{construyó} (indéfini, action datée \emph{En 2015}).
\item \esp{Hizo} (indéfini, événement météo ponctuel \emph{aquel día}), \esp{pudieron} (indéfini, capacité réalisée ce jour-là).
\item \esp{empezó} (indéfini, début soudain \emph{de repente}), \esp{corrieron} (indéfini, réaction immédiate).
\item \esp{había} (imparfait, existence passé \emph{antes}).
\item \esp{contaba} (imparfait, action en cours/description fond), \esp{cenábamos} (imparfait, action en cours simultanée \emph{mientras}).
\item \esp{volví} (indéfini, action ponctuelle finale \emph{al final}).
\end{enumerate}
\end{corrige}

\begin{corrige}[Exercice 2 : Thème]
\begin{enumerate}
\item \esp{El éxodo rural vacía los pueblos.}
\item \esp{A causa de la falta de carreteras, las cosechas se pudren.} (ou \esp{debido a la falta...})
\item \esp{En 2020, los jóvenes crearon una asociación para el desarrollo local.}
\item \esp{Gracias a la cooperativa, los agricultores venden su cacao a un precio justo.}
\item \esp{Antes, no había escuela en el pueblo.}
\item \esp{Pierre volvió de la ciudad el año pasado.}
\item \esp{Ahora, él trabaja en la unidad de transformación.}
\item \esp{El turismo rural empieza a desarrollarse.}
\end{enumerate}
\end{corrige}

\begin{corrige}[Exercice 3 : Version]
« La coopérative n'achète pas seulement le cacao. Elle organise aussi des ateliers de formation pour les agriculteurs. Avant, ils ne savaient pas comment améliorer la qualité. Maintenant, grâce aux techniciens, la production a augmenté. Les jeunes voient qu'il y a un avenir et ne partent pas. C'est un changement radical pour la communauté. »
\end{corrige}

\begin{corrige}[Exercice 4 : Expression — Causes et conséquences (Proposition)]
\esp{La movilidad rural en mi región es un fenómeno visible. Muchos jóvenes se van a Yaoundé o Douala porque no hay fábricas ni empresas en los pueblos y a causa de la mala calidad de las carreteras. Además, los servicios de salud y educación son insuficientes. Como consecuencia, los campos quedan abandonados y los pueblos pierden su vitalidad. Por eso, el gobierno y las ONG deben invertir en el desarrollo local. Gracias a la creación de cooperativas agrícolas y la mejora del transporte, la juventud podría quedarse y trabajar en su tierra.}
\end{corrige}

\begin{corrige}[Exercice 5 : Expression — Récit (Proposition)]
\esp{El año pasado pasé quince días en el pueblo de mis abuelos, cerca de Dschang. Cuando llegué, el pueblo era muy verde y las casas eran de ladrillo rojo. Hacía fresco por las mañanas. Todos los días me despertaba al canto del gallo y ayudaba a mi abuela en la huerta: regábamos los tomates y recogíamos la leña. Una tarde, mientras caminábamos por el sendero, vimos una serpiente grande y corrimos muy rápido. Aquella noche, mi abuelo mató un pollo y cenamos todos juntos alrededor del fuego. Me contó historias de la independencia de Camerún. Al final de la estancia, volví a Douala cansado pero feliz: entendí la dureza y la belleza de la vida en el campo.}
\textbf{Vérification temps :} Indéfini (\esp{pasé, llegué, vimos, corrimos, mató, cenamos, contó, volví, entendí} = 9) ; Imparfait (\esp{era, eran, hacía, me despertaba, ayudaba, regábamos, recogíamos, caminábamos} = 8). Conforme.
\end{corrige}

\begin{attention}[Les 5 erreurs qui coûtent des points au Bac]
\begin{enumerate}
\item \textbf{Confondre indéfini et imparfait} : \esp{Ayer iba al mercado} est faux pour une action unique ; écris \esp{Ayer fui al mercado}. Réflexe : date ou action unique = indéfini ; habitude ou décor = imparfait.
\item \textbf{Oublier les accents} : \esp{exodo}, \esp{migracion} sans accent sont fautifs : écris \esp{éxodo}, \esp{migración}. \esp{cooperativa} (mot grave terminé par voyelle) ne prend pas d'accent. Attention aussi à \esp{está} (verbe, il/elle est) contre \esp{esta} (adjectif démonstratif, ce/cette) et \esp{él} (lui) vs \esp{el} (le).
\item \textbf{Calquer le français} : « il n'y avait pas de travail » ne se traduit pas mot à mot ; dis \esp{no había trabajo}. De même, « quitter le village » = \esp{irse del pueblo} ou \esp{abandonar el pueblo}, jamais un verbe inventé. « Se rendre compte » = \esp{darse cuenta}, pas \esp{realizarse}.
\item \textbf{Négliger les accords} : \esp{las carretera buena} est doublement faux ; écris \esp{las carreteras buenas}. L'adjectif s'accorde en genre et en nombre avec le nom, même placé après. \esp{Unos jóvenes cansados}, \esp{unas carreteras nuevas}.
\item \textbf{Tomber dans les faux amis} : \esp{éxito} signifie « succès », pas « exode » (\esp{éxodo}) ; \esp{actualmente} signifie « actuellement », pas « en réalité » ; \esp{asistir a la escuela} = « aller à l'école », pas « assister l'école » ; \esp{realizar un sueño} = « réaliser un rêve (accomplir) », \esp{darse cuenta de} = « se rendre compte » ; \esp{la carpeta} = « le dossier/chemise », pas « la carpette » ; \esp{el éxito} = « le succès ».
\end{enumerate}
\end{attention}

\newpage

\coursec{Exercices}

\courssub{Je vérifie mes connaissances}

\begin{exercice}[QCM : le monde rural]
Entoure la bonne réponse.
\begin{enumerate}
\item Un \esp{agricultor} est une personne qui :
\begin{itemize}
\item[a)] travaille dans une usine
\item[b)] cultive la terre
\item[c)] conduit un taxi
\end{itemize}
\item Le \esp{éxodo rural} désigne :
\begin{itemize}
\item[a)] l'arrivée des citadins au village
\item[b)] le départ massif des habitants de la campagne vers les villes
\item[c)] la fête de la moisson
\end{itemize}
\item Une \esp{cooperativa} est :
\begin{itemize}
\item[a)] une association de producteurs qui travaillent ensemble
\item[b)] une grande autoroute
\item[c)] un impôt sur les récoltes
\end{itemize}
\item \esp{La cosecha} signifie :
\begin{itemize}
\item[a)] la semence
\item[b)] la récolte
\item[c)] la sécheresse
\end{itemize}
\end{enumerate}
\end{exercice}

\begin{exercice}[Vrai ou faux ? Justifie]
Réponds par \esp{verdadero} ou \esp{falso} et justifie ta réponse par une phrase complète en espagnol.
\begin{enumerate}
\item L'exode rural n'a que des conséquences positives pour les villages.
\item Les coopératives permettent aux agriculteurs de mieux vendre leurs récoltes.
\item Le manque de routes (\esp{carreteras}) est une cause de l'exode rural.
\item L'imparfait sert à raconter une action ponctuelle et terminée du passé.
\end{enumerate}
\end{exercice}

\begin{exercice}[Vocabulaire : à chaque mot sa définition]
Relie chaque mot espagnol à sa traduction française.

\begin{center}
\begin{tabularx}{\linewidth}{|X|X|}
\hline
\rowcolor{popblueL} \textbf{Español} & \textbf{Français} \\
\hline
1. el campo & a) la migration \\
\hline
2. la migración & b) la campagne \\
\hline
3. el transporte & c) le développement local \\
\hline
4. el desarrollo local & d) le transport \\
\hline
5. la carretera & e) la route \\
\hline
\end{tabularx}
\end{center}
\end{exercice}

\begin{exercice}[Conjugaison : indéfini ou imparfait ?]
Complète les phrases avec le verbe entre parenthèses au prétérit indéfini ou à l'imparfait.
\begin{enumerate}
\item Cuando yo \_\_\_\_\_\_\_\_ (ser) niño, \_\_\_\_\_\_\_\_ (pasar) las vacaciones en el pueblo de mis abuelos.
\item El año pasado, muchos jóvenes \_\_\_\_\_\_\_\_ (dejar) el campo para ir a la ciudad.
\item Ayer nosotros \_\_\_\_\_\_\_\_ (visitar) una cooperativa de cacao.
\item Antes, los agricultores \_\_\_\_\_\_\_\_ (vender) sus productos en el mercado del barrio.
\item En 2015, el gobierno \_\_\_\_\_\_\_\_ (construir) una carretera nueva hacia el pueblo.
\end{enumerate}
\end{exercice}

\courssub{Je m'entraîne}

\begin{exercice}[Thème : du français à l'espagnol]
Traduis les phrases suivantes en espagnol. Attention au choix entre imparfait et prétérit indéfini.
\begin{enumerate}
\item Quand j'étais petit, ma grand-mère cultivait le manioc au village.
\item L'année dernière, mon cousin a quitté le village pour travailler à Douala.
\item Tous les dimanches, nous allions au marché avec mon père.
\item Hier, les agriculteurs ont vendu toute leur récolte de cacao.
\item Avant, il n'y avait pas de route entre le village et la ville.
\end{enumerate}
\end{exercice}

\begin{exercice}[Texte à trous : une histoire au passé]
Mets les verbes entre parenthèses au temps correct (prétérit indéfini ou imparfait) et justifie ton choix (action ponctuelle, habitude ou description).

\begin{textebox}[Un recuerdo del pueblo]
Cuando yo (1. tener) \_\_\_\_\_\_\_\_ diez años, mi familia (2. vivir) \_\_\_\_\_\_\_\_ en un pueblo cerca de Bafoussam. Todos los días, mi abuelo (3. levantarse) \_\_\_\_\_\_\_\_ muy temprano y (4. ir) \_\_\_\_\_\_\_\_ al campo. Un día, mientras nosotros (5. trabajar) \_\_\_\_\_\_\_\_ en la cosecha, (6. llegar) \_\_\_\_\_\_\_\_ un camión de la cooperativa. Ese día, mi abuelo (7. vender) \_\_\_\_\_\_\_\_ todo su café y (8. estar) \_\_\_\_\_\_\_\_ muy contento.
\end{textebox}
\end{exercice}

\begin{exercice}[Appariement : causes, conséquences, solutions]
Complète le tableau avec les expressions suivantes : \esp{falta de carreteras} ; \esp{creación de cooperativas} ; \esp{envejecimiento del pueblo} ; \esp{búsqueda de trabajo} ; \esp{electrificación rural} ; \esp{pérdida de tradiciones}.

\begin{center}
\begin{tabularx}{\linewidth}{|X|X|X|}
\hline
\rowcolor{popblueL} \textbf{Causas} & \textbf{Consecuencias} & \textbf{Soluciones} \\
\hline
 & & \\
\hline
 & & \\
\hline
 & & \\
\hline
\end{tabularx}
\end{center}
\end{exercice}

\begin{exercice}[Transformation de phrases]
Transforme chaque phrase selon la consigne.
\begin{enumerate}
\item \esp{Los jóvenes dejan el campo.} $\rightarrow$ Mets la phrase au prétérit indéfini avec \esp{el año pasado}.
\item \esp{Mi abuelo cultiva maíz.} $\rightarrow$ Mets la phrase à l'imparfait avec \esp{antes}.
\item \esp{La cooperativa compra el cacao de los agricultores.} $\rightarrow$ Mets la phrase au prétérit indéfini avec \esp{ayer}.
\item \esp{Hay mucha gente en el mercado del pueblo.} $\rightarrow$ Mets la phrase à l'imparfait avec \esp{cuando era niño}.
\end{enumerate}
\end{exercice}

\courssub{Je me prépare au Bac}

\begin{exercice}[Compréhension de texte type Bac]
Lis le texte suivant, puis réponds aux questions par des phrases complètes en espagnol.

\begin{textebox}[Una cooperativa que cambia vidas]
En una pequeña región de los Andes, los campesinos de quince pueblos decidieron hace diez años crear una cooperativa de café. Antes, cada agricultor vendía solo su cosecha a intermediarios que pagaban precios muy bajos. Muchos jóvenes dejaban el campo y emigraban a la capital porque no veían futuro en sus pueblos.

Con la cooperativa, todo cambió. Los agricultores aprendieron nuevas técnicas de cultivo, compraron máquinas juntos y empezaron a vender su café directamente al extranjero. Con los beneficios, construyeron una escuela y un centro de salud. Además, repararon la carretera que une los pueblos con la ciudad más cercana.

Hoy, algunos jóvenes que habían emigrado han vuelto al pueblo. «Antes soñaba con la ciudad; ahora sueño con mi café», dice sonriendo uno de ellos. El desarrollo local, cuando es obra de todos, puede detener el éxodo rural.
\end{textebox}

\textbf{Glossaire} : \esp{intermediarios} = intermédiaires ; \esp{beneficios} = bénéfices ; \esp{centro de salud} = centre de santé ; \esp{habían emigrado} = avaient émigré ; \esp{detener} = arrêter, freiner.

\textbf{Questions de compréhension} (réponds en espagnol, par des phrases complètes) :
\begin{enumerate}
\item ¿Qué decidieron crear los campesinos hace diez años?
\item ¿Por qué emigraban los jóvenes antes de la cooperativa?
\item Menciona dos cosas que los agricultores hicieron gracias a la cooperativa.
\item ¿Qué construyeron con los beneficios?
\item Según el texto, ¿qué puede detener el éxodo rural?
\end{enumerate}

\textbf{Questions de langue} :
\begin{enumerate}
\item Relève dans le texte deux verbes au prétérit indéfini et deux verbes à l'imparfait, puis justifie l'emploi de chaque temps.
\item Trouve dans le texte un synonyme de \esp{agricultores} et un synonyme de \esp{empezaron}.
\end{enumerate}
\end{exercice}

\begin{exercice}[Version et thème type Bac]
\textbf{A. Version.} Traduis en français le passage suivant (environ 5 lignes).

\begin{textebox}[El éxodo rural]
Antes, la vida en el campo era dura: no había electricidad ni buenas carreteras, y los jóvenes no encontraban trabajo. Por eso, muchas familias dejaron sus pueblos y se instalaron en las grandes ciudades. Hoy, algunas cooperativas intentan desarrollar el campo para que los jóvenes vuelvan.
\end{textebox}

\textbf{B. Thème.} Traduis en espagnol les phrases suivantes.
\begin{enumerate}
\item Quand nous habitions au village, nous aidions nos parents pendant la récolte.
\item En 2018, ma sœur a quitté le village pour étudier à Yaoundé.
\item Les agriculteurs vendaient leur cacao à bas prix avant la création de la coopérative.
\item Hier, le chef du village a inauguré une nouvelle route.
\item Avant, il n'y avait ni école ni centre de santé dans mon village.
\end{enumerate}
\end{exercice}

\begin{exercice}[Expression écrite type Bac]
\textbf{Sujet} : \esp{Cuenta una experiencia en el campo o en tu pueblo.}

Rédige un texte de \textbf{80 à 100 mots} en espagnol. Tu raconteras un souvenir réel ou imaginaire : un séjour au village, une fête, une récolte, un travail dans une coopérative...

\textbf{Consignes} :
\begin{itemize}
\item Utilise l'imparfait pour les descriptions et les habitudes, le prétérit indéfini pour les actions ponctuelles.
\item Emploie au moins cinq mots du lexique de la leçon (\esp{campo, agricultor, cosecha, carretera, cooperativa, desarrollo local}...).
\item Organise ton texte : présentation du lieu et des personnes, déroulement du souvenir, impression finale.
\end{itemize}

\textbf{Grille critériée} (barème indicatif sur 10 points) :

\begin{center}
\begin{tabularx}{\linewidth}{|X|c|}
\hline
\rowcolor{popblueL} \textbf{Critère} & \textbf{Points} \\
\hline
Respect du sujet et de la longueur (80-100 mots) & 2 \\
\hline
Emploi correct de l'imparfait et du prétérit indéfini & 3 \\
\hline
Lexique du monde rural (au moins 5 mots de la leçon) & 2 \\
\hline
Organisation et cohérence du récit & 2 \\
\hline
Orthographe et accents & 1 \\
\hline
\end{tabularx}
\end{center}
\end{exercice}

\newpage

\coursec{Corrigés des exercices}

\coursec{Leçon EA10 : El mundo rural, la movilidad y el desarrollo local}

\courssub{Texte support et découverte}

\begin{textebox}[Texte support : La transformación de Nkolmekok — Lecture]
Hace diez años, Nkolmekok era un pueblo aislado cerca de Yaoundé. No había carretera asfaltada, ni electricidad, ni centro de salud. Los agricultores cultivaban cacao y mandioca, pero vendían su cosecha a intermediarios a precios muy bajos. Los jóvenes no veían futuro y emigraban a la capital para buscar trabajo. El pueblo envejecía y las tradiciones se perdían.

Un día, los campesinos de quince pueblos decidieron crear una cooperativa. Gracias a la cooperativa, aprendieron nuevas técnicas, compraron una máquina para secar el cacao y empezaron a vender directamente al extranjero. Con los beneficios, construyeron una escuela, un centro de salud y repararon la carretera. Hoy, algunos jóvenes que se habían ido están volviendo. El desarrollo local, cuando es obra de todos, puede detener el éxodo rural.
\end{textebox}

\begin{definition}[Glossaire du texte support]
\begin{itemize}
\item \esp{aislado} : isolé.
\item \esp{asfaltada} (route) : goudronnée.
\item \esp{intermediarios} : intermédiaires (négociants).
\item \esp{envejecía} (imparfait de \esp{envejecer}) : vieillissait.
\item \esp{campesinos} : paysans, cultivateurs (synonyme de \esp{agricultores}).
\item \esp{secár} : sécher (le cacao) ; \esp{una máquina para secar} : une sécheuse.
\item \esp{extranjero} : l'étranger ; \esp{al extranjero} : à l'étranger.
\item \esp{beneficios} : bénéfices, profits.
\item \esp{obra de todos} : œuvre de tous, affaire de tous.
\item \esp{detener} : arrêter, stopper.
\end{itemize}
\end{definition}

\courssub{Compréhension du texte}

\begin{exercice}[Questions de compréhension sur le texte support]
Répondez en espagnol par des phrases complètes.
\begin{enumerate}
\item ¿Cómo era Nkolmekok hace diez años? (Citez trois éléments.)
\item ¿Por qué emigraban los jóvenes a la capital?
\item ¿Qué decidieron hacer los campesinos de quince pueblos?
\item ¿Cuáles fueron tres resultados con los beneficios de la cooperativa?
\item Según la última frase, ¿qué puede hacer el desarrollo local?
\end{enumerate}
\end{exercice}

\begin{corrige}[Exercice 2 — Compréhension du texte support]
\begin{enumerate}
\item \esp{Hace diez años, Nkolmekok era un pueblo aislado, sin carretera asfaltada, sin electricidad y sin centro de salud.}
\item \esp{Los jóvenes emigraban porque no veían futuro en el pueblo y buscaban trabajo en la capital.}
\item \esp{Decidieron crear una cooperativa.}
\item \esp{Construyeron una escuela, un centro de salud y repararon la carretera.} (Deux éléments sur trois suffisent.)
\item \esp{El desarrollo local, cuando es obra de todos, puede detener el éxodo rural.}
\end{enumerate}
\textbf{Point de méthode} : Reprenez les mots de la question pour construire la réponse ; ne recopiez pas tout le paragraphe.
\end{corrige}

\courssub{Lexique thématique : le monde rural et le développement}

\begin{propriete}[Lexique de base — Mots du programme]
\begin{center}
\begin{tabularx}{\linewidth}{|X|X|X|}
\hline
\rowcolor{popblueL} \textbf{Español} & \textbf{Français} & \textbf{Exemple / Contexte} \\
\hline
el campo & la campagne, le champ & \esp{Viven en el campo.} / \esp{El campo de maíz.} \\
\hline
el agricultor / la agricultora & l'agriculteur, le paysan & \esp{El agricultor cultiva la tierra.} \\
\hline
el campesino / la campesina & le paysan (connotation sociale) & \esp{Los campesinos se organizan.} \\
\hline
la cosecha & la récolte (action et produit) & \esp{La cosecha de cacao fue buena.} \\
\hline
el éxodo rural & l'exode rural & \esp{El éxodo rural vacía los pueblos.} \\
\hline
la migración & la migration & \esp{La migración campo-ciudad.} \\
\hline
el transporte & le transport & \esp{El transporte de mercancías.} \\
\hline
la carretera & la route (goudronnée, interurbaine) & \esp{Construyeron una carretera nueva.} \\
\hline
el desarrollo local & le développement local & \esp{Proyectos de desarrollo local.} \\
\hline
la cooperativa & la coopérative & \esp{Una cooperativa de café.} \\
\hline
el intermediario & l'intermédiaire (commercial) & \esp{Venden sin intermediarios.} \\
\hline
la sequía & la sécheresse & \esp{La sequía arruinó la cosecha.} \\
\hline
la semilla & la semence, la graine & \esp{Comprar semillas mejoradas.} \\
\hline
\end{tabularx}
\end{center}
\end{propriete}

\begin{attention}[Faux amis et pièges fréquents]
\begin{itemize}
\item \esp{El campo} $\neq$ \emph{le camping} (camping = \esp{el camping} ou \esp{el campamento}). \esp{Campo} = campagne / champ.
\item \esp{La cosecha} (récolte) $\neq$ \esp{la semilla} (semence).
\item \esp{Un agricultor} cultive la terre (propriétaire ou fermier) ; ce n'est pas nécessairement un salarié (\esp{un jornalero} / \esp{un peón}).
\item \esp{La carretera} = route goudronnée ; \esp{el camino} = chemin, piste ; \esp{la calle} = rue (en ville).
\item Genre : \esp{la carretera, la migración, la cosecha, la cooperativa} (féminin) ; \esp{el campo, el transporte, el desarrollo, el éxodo, el agricultor} (masculin).
\item \esp{Búsqueda} (recherche) prend un accent sur le \emph{u} (\esp{buscar} $\rightarrow$ \esp{búsqueda}).
\end{itemize}
\end{attention}

\begin{corrige}[Exercice 1 — QCM : le monde rural]
\textbf{Réponses :}
\begin{enumerate}
\item \textbf{b)} cultive la terre. Un \esp{agricultor} est un cultivateur, un paysan. Attention au faux ami : \esp{agricultor} ne signifie pas « ouvrier agricole salarié » (\esp{jornalero}) mais toute personne qui cultive la terre.
\item \textbf{b)} le départ massif des habitants de la campagne vers les villes. \esp{Éxodo rural} = exode rural. On dit aussi \esp{la emigración del campo a la ciudad}.
\item \textbf{a)} une association de producteurs qui travaillent ensemble. Une \esp{cooperativa} met en commun les moyens (machines, vente, formation) de ses membres.
\item \textbf{b)} la récolte. \esp{La cosecha} désigne le moment où l'on récolte et le produit récolté : \esp{la cosecha de cacao} = la récolte de cacao.
\end{enumerate}
\textbf{Point de vigilance} : ne pas confondre \esp{la cosecha} (la récolte) et \esp{la semilla} (la semence, la graine) ; \esp{la sequía} = la sécheresse.
\end{corrige}

\begin{corrige}[Exercice 3 — Vocabulaire : à chaque mot sa définition]
\begin{center}
\begin{tabularx}{\linewidth}{|X|X|X|}
\hline
\rowcolor{popblueL} \textbf{Español} & \textbf{Français} & \textbf{Exemple} \\
\hline
1. el campo & b) la campagne & \esp{Vivo en el campo.} \\
\hline
2. la migración & a) la migration & \esp{La migración a la ciudad aumenta.} \\
\hline
3. el transporte & d) le transport & \esp{El transporte es difícil sin carreteras.} \\
\hline
4. el desarrollo local & c) le développement local & \esp{La cooperativa ayuda al desarrollo local.} \\
\hline
5. la carretera & e) la route & \esp{La carretera une el pueblo con la ciudad.} \\
\hline
\end{tabularx}
\end{center}
\textbf{Réponses} : 1-b ; 2-a ; 3-d ; 4-c ; 5-e.
\textbf{Point de vigilance} : \esp{el campo} signifie « la campagne » mais aussi « le champ » ; ne pas le confondre avec \esp{el camping}. Attention au genre : on dit \esp{la carretera}, \esp{la migración}, mais \esp{el transporte}, \esp{el desarrollo}.
\end{corrige}

\begin{corrige}[Exercice 7 — Appariement : causes, conséquences, solutions]
\begin{center}
\begin{tabularx}{\linewidth}{|X|X|X|}
\hline
\rowcolor{popblueL} \textbf{Causas} & \textbf{Consecuencias} & \textbf{Soluciones} \\
\hline
falta de carreteras & envejecimiento del pueblo & creación de cooperativas \\
\hline
búsqueda de trabajo & pérdida de tradiciones & electrificación rural \\
\hline
\end{tabularx}
\end{center}
\textbf{Explications :}
\begin{itemize}
\item \textbf{Causes} (pourquoi les gens partent) : \esp{falta de carreteras} (le manque de routes isole le village) ; \esp{búsqueda de trabajo} (les jeunes cherchent du travail en ville).
\item \textbf{Conséquences} (ce qui arrive ensuite) : \esp{envejecimiento del pueblo} (il ne reste que les anciens) ; \esp{pérdida de tradiciones} (les coutumes disparaissent).
\item \textbf{Soluciones} (ce qu'on peut faire) : \esp{creación de cooperativas} (vendre mieux, ensemble) ; \esp{electrificación rural} (l'électricité permet de vivre et de travailler au village).
\end{itemize}
\textbf{Point de vigilance} : \esp{búsqueda} (la recherche) prend un accent sur le u ; son verbe est \esp{buscar} (chercher).
\end{corrige}

\courssub{Grammaire : Prétérit indéfini et Imparfait — Raconter le passé}

\begin{propriete}[Choix entre Prétérit Indéfini et Imparfait]
\textbf{Le Prétérit Indéfini (\emph{pretérito indefinido})} exprime une action \textbf{ponctuelle, terminée, datée}, qui fait avancer le récit.
\begin{itemize}
\item Marqueurs fréquents : \esp{ayer, anteayer, el año pasado, en 2015, un día, ese día, hace dos días, la semana pasada}.
\item Exemple : \esp{Ayer vendí mi cosecha.} (Action unique, hier).
\end{itemize}

\textbf{L'Imparfait (\emph{pretérito imperfecto})} exprime : 
\begin{enumerate}
\item Une \textbf{description} (physique, météo, âge, état d'esprit, décor) dans le passé.
\item Une \textbf{habitude}, une action \textbf{répétée} (\emph{todos los días, antes, siempre, cuando era niño}).
\item Une action \textbf{en cours} (arrière-plan) interrompue par une autre (souvent introduite par \esp{mientras}).
\end{enumerate}
Marqueurs fréquents : \esp{antes, antes de, todos los días, cada semana, cuando era niño, siempre, normalmente, mientras}.

\textbf{Règle d'or narrative} : \esp{Mientras + Imparfait (arrière-plan), + Indéfini (événement ponctuel)}.
\esp{Ejemplo: Mientras trabajábamos (imparfait), llegó (indéfini) el camión.}
\end{propriete}

\begin{propriete}[Conjugaison — Rappel des formes]
\textbf{Indéfini (réguliers) :}
\begin{center}
\begin{tabular}{|l|l|l|l|}
\hline
\rowcolor{popblueL} \textbf{Persona} & \textbf{-AR (hablar)} & \textbf{-ER (vender)} & \textbf{-IR (vivir)} \\
\hline
yo & hablé & vendí & viví \\
\hline
tú & hablaste & vendiste & viviste \\
\hline
él/ella/usted & habló & vendió & vivió \\
\hline
nosotros/as & hablamos & vendimos & vivimos \\
\hline
vosotros/as & hablasteis & vendisteis & vivisteis \\
\hline
ellos/ellas/ustedes & hablaron & vendieron & vivieron \\
\hline
\end{tabular}
\end{center}

\textbf{Imparfait (réguliers) :}
\begin{center}
\begin{tabular}{|l|l|l|l|}
\hline
\rowcolor{popblueL} \textbf{Persona} & \textbf{-AR (cultivar)} & \textbf{-ER (vender)} & \textbf{-IR (vivir)} \\
\hline
yo & cultivaba & vendía & vivía \\
\hline
tú & cultivabas & vendías & vivías \\
\hline
él/ella/usted & cultivaba & vendía & vivía \\
\hline
nosotros/as & cultivábamos & vendíamos & vivíamos \\
\hline
vosotros/as & cultivabais & vendíais & vivíais \\
\hline
ellos/ellas/ustedes & cultivaban & vendían & vivían \\
\hline
\end{tabular}
\end{center}

\textbf{Irréguliers importants (Indéfini) :}
\begin{itemize}
\item \esp{Ser / Ir} : \esp{fui, fuiste, fue, fuimos, fuisteis, fueron} (mêmes formes).
\item \esp{Estar} : \esp{estuve, estuviste, estuvo, estuvimos, estuvisteis, estuvieron}.
\item \esp{Tener} : \esp{tuve, tuviste, tuvo, tuvimos, tuvisteis, tuvieron}.
\item \esp{Hacer} : \esp{hice, hiciste, hizo, hicimos, hicisteis, hicieron}.
\item \esp{Poder} : \esp{pude, pudiste, pudo, pudimos, pudisteis, pudieron}.
\item \esp{Querer} : \esp{quise, quisiste, quiso, quisimos, quisisteis, quisieron}.
\item \esp{Saber} : \esp{supe, supiste, supo, supimos, supisteis, supieron}.
\item \esp{Venir} : \esp{vine, viniste, vino, vinimos, vinisteis, vinieron}.
\item \esp{Decir} : \esp{dije, dijiste, dijo, dijimos, dijisteis, dijeron}.
\item \esp{Traer} : \esp{traje, trajiste, trajo, trajimos, trajisteis, trajeron}.
\item \esp{Construir / Leer / Oír / Caer} : changement \emph{i} $\rightarrow$ \emph{y} à la 3\ieme{} pers. sg/pl : \esp{construyó, construyeron ; leyó, leyeron ; oyó, oyeron ; cayó, cayeron}.
\end{itemize}

\textbf{Irréguliers (Imparfait) — Seulement 3 verbes :}
\begin{itemize}
\item \esp{Ser} : \esp{era, eras, era, éramos, erais, eran}.
\item \esp{Ir} : \esp{iba, ibas, iba, íbamos, ibais, iban} (accent sur \emph{í} à \emph{nosotros}).
\item \esp{Ver} : \esp{veía, veías, veía, veíamos, veíais, veían}.
\end{itemize}

\textbf{\esp{Haber} impersonnel (il y a / il y avait) :} \esp{Hay} (présent), \esp{Había} (imparfait), \esp{Hubo} (indéfini). \textbf{Invariable} (jamais \esp{*habían}).
\end{propriete}

\begin{corrige}[Exercice 4 — Conjugaison : indéfini ou imparfait ?]
\begin{enumerate}
\item \esp{Cuando yo \textbf{era} niño, \textbf{pasaba} las vacaciones en el pueblo de mis abuelos.} — Deux imparfaits : description (l'âge, \esp{era} de \esp{ser}) + habitude répétée chaque année (\esp{pasaba}).
\item \esp{El año pasado, muchos jóvenes \textbf{dejaron} el campo para ir a la ciudad.} — Indéfini : action ponctuelle datée (\esp{el año pasado}). \esp{Dejar} régulier.
\item \esp{Ayer nosotros \textbf{visitamos} una cooperativa de cacao.} — Indéfini : action unique et terminée (\esp{ayer}). Attention : \esp{visitamos} = 1\iere{} personne du pluriel de l'indéfini (identique au présent pour les verbes en -ar) ; c'est \esp{ayer} qui impose le passé.
\item \esp{Antes, los agricultores \textbf{vendían} sus productos en el mercado del barrio.} — Imparfait : habitude passée (\esp{antes}).
\item \esp{En 2015, el gobierno \textbf{construyó} una carretera nueva hacia el pueblo.} — Indéfini : action ponctuelle datée. Attention à la forme irrégulière : \esp{construir} $\rightarrow$ \esp{construyó} (i $\rightarrow$ y à la 3\ieme{} personne du singulier et du pluriel : \esp{construyeron}).
\end{enumerate}
\textbf{Astuce} : les marqueurs \esp{ayer, el año pasado, en 2015, un día, ese día} appellent l'indéfini ; \esp{antes, todos los días, cuando era niño, siempre, mientras} appellent l'imparfait.
\end{corrige}

\begin{corrige}[Exercice 6 — Texte à trous : une histoire au passé]
\begin{textebox}[Un recuerdo del pueblo — Corrigé]
Cuando yo (1) \textbf{tenía} diez años, mi familia (2) \textbf{vivía} en un pueblo cerca de Bafoussam. Todos los días, mi abuelo (3) \textbf{se levantaba} muy temprano y (4) \textbf{iba} al campo. Un día, mientras nosotros (5) \textbf{trabajábamos} en la cosecha, (6) \textbf{llegó} un camión de la cooperativa. Ese día, mi abuelo (7) \textbf{vendió} todo su café y (8) \textbf{estuvo} muy contento.
\end{textebox}
\textbf{Justifications :}
\begin{enumerate}
\item \esp{tenía} : imparfait, description (l'âge, le cadre).
\item \esp{vivía} : imparfait, situation durable dans le passé.
\item \esp{se levantaba} : imparfait, habitude (\esp{todos los días}). Verbe pronominal : le pronom \esp{se} se place devant le verbe conjugué.
\item \esp{iba} : imparfait de \esp{ir} (irrégulier), habitude.
\item \esp{trabajábamos} : imparfait, action en cours (\esp{mientras} = pendant que), arrière-plan.
\item \esp{llegó} : indéfini, action ponctuelle qui interrompt l'arrière-plan (\esp{un día}).
\item \esp{vendió} : indéfini, action ponctuelle et terminée (\esp{ese día}).
\item \esp{estuvo} : indéfini de \esp{estar} (irrégulier : \esp{estuve, estuviste, estuvo...}) : ici, réaction ponctuelle, bornée dans le temps (\esp{ese día}).
\end{enumerate}
\textbf{Traduction} : « Quand j'avais dix ans, ma famille vivait dans un village près de Bafoussam. Tous les jours, mon grand-père se levait très tôt et allait aux champs. Un jour, pendant que nous travaillions à la récolte, un camion de la coopérative arriva. Ce jour-là, mon grand-père vendit tout son café et fut très content. »
\textbf{Point de vigilance} : la structure \esp{mientras + imparfait, + indéfini} est un classique : l'imparfait peint le décor, l'indéfini fait avancer l'histoire.
\end{corrige}

\begin{corrige}[Exercice 8 — Transformation de phrases (choix du temps et conjugaison)]
\begin{enumerate}
\item \esp{El año pasado, los jóvenes \textbf{dejaron} el campo.} — Indéfini, 3\ieme{} personne du pluriel : \esp{dejar} régulier (\esp{dejaron}).
\item \esp{Antes, mi abuelo \textbf{cultivaba} maíz.} — Imparfait : habitude passée. Accent sur le \emph{a} : \esp{cultivaba}.
\item \esp{Ayer, la cooperativa \textbf{compró} el cacao de los agricultores.} — Indéfini : action datée. Accent obligatoire sur le \emph{ó} : \esp{compró} (sinon on lit le présent \esp{compro} = j'achète).
\item \esp{Cuando era niño, \textbf{había} mucha gente en el mercado del pueblo.} — Imparfait : description. \esp{había} est invariable (forme impersonnelle de \esp{haber}) même devant un pluriel (\esp{mucha gente} est singulier, mais on dirait aussi \esp{había muchas personas}).
\end{enumerate}
\textbf{Points de vigilance} : l'accent final de la 3\ieme{} personne du singulier de l'indéfini (\esp{compró, vendió, llegó, construyó}) distingue le passé du présent ; l'oublier est une faute éliminatoire au Bac. Le changement de temps exige souvent un marqueur temporel cohérent (\esp{el año pasado, antes, ayer, cuando era niño}).
\end{corrige}

\courssub{Méthode du commentaire et commentaire modèle}

\begin{methode}[Commentaire de texte type Bac — 3 étapes]
\begin{enumerate}
\item \textbf{Analyse globale (Contexte) :} Identifier le thème (monde rural, exode, développement), le type de texte (récit, article, témoignage), la source, l'intention (informer, sensibiliser). Noter la structure (introduction, développement, conclusion).
\item \textbf{Analyse détaillée (Contenu + Langue) :} 
\begin{itemize}
\item \emph{Contenu :} Résumer les idées principales par paragraphes. Relever les arguments (causes, conséquences, solutions). Citer des exemples précis du texte.
\item \emph{Langue :} Relever les temps du récit (Indéfini / Imparfait) et justifier leur emploi (ponctuel vs description/habitude). Relever le lexique thématique (champ sémantique). Repérer les connecteurs logiques (\esp{por eso, gracias a, pero, hoy}).
\end{itemize}
\item \textbf{Synthèse / Appréciation personnelle :} Répondre à la question de synthèse (souvent : « Que pensez-vous de... ? », « Selon le texte, comment... ? »). Exprimer un avis argumenté en réutilisant le vocabulaire de la leçon.
\end{enumerate}
\end{methode}

\begin{exemplebox}[Commentaire modèle sur le texte support « La transformación de Nkolmekok »]
\textbf{1. Présentation} : Le texte « La transformación de Nkolmekok » est un récit qui illustre le thème du développement local face à l'exode rural. Il oppose la situation passée (il y a 10 ans) à la situation actuelle.

\textbf{2. Analyse du contenu} : 
\begin{itemize}
\item \textbf{Paragraphe 1 (Passé / Imparfait dominant) :} Description d'un village isolé (\esp{aislado, no había carretera, no había electricidad}). Causes de l'exode : prix bas (\esp{precios muy bajos}) à cause des intermédiaires, manque d'avenir pour les jeunes (\esp{no veían futuro}),
\end{itemize}
\end{exemplebox}


\newpage

\coursec{Fiche bilan}

\begin{popfigure}
\begin{tikzpicture}[mindmap, grow cyclic, every node/.style={concept, font=\small},
root concept/.append style={concept color=popblue, font=\bfseries\small, minimum size=3.2cm},
level 1 concept/.append style={sibling angle=60, level distance=4.1cm, minimum size=2.6cm}]
\node[root concept]{El mundo rural, la movilidad y el desarrollo local}
child[concept color=popgreen]{node[align=center]{Lexique\\ \esp{campo, agricultor, cosecha}}}
child[concept color=poporange]{node[align=center]{Éxodo rural\\ \esp{causas y consecuencias}}}
child[concept color=popgold]{node[align=center]{Movilidad\\ \esp{transporte, carretera, migración}}}
child[concept color=poppurple]{node[align=center]{Desarrollo local\\ \esp{cooperativa, proyectos}}}
child[concept color=poppink]{node[align=center]{Gramática\\ pretérito indefinido}}
child[concept color=popblue!60]{node[align=center]{Gramática\\ pretérito imperfecto}};
\end{tikzpicture}
\legende{Carte mentale de la leçon : le monde rural, la mobilité et le développement local}
\end{popfigure}

\begin{bilan}
\textbf{1. Lexique clé à connaître par cœur}
\begin{center}
\begin{tabularx}{\linewidth}{|X|X|}
\hline
\rowcolor{popblueL} \textbf{Español} & \textbf{Français} \\
\hline
\esp{el campo} & la campagne \\
\hline
\esp{el/la agricultor(a)} & l'agriculteur / la cultivatrice \\
\hline
\esp{la cosecha} & la récolte \\
\hline
\esp{sembrar / cultivar / recoger} & semer / cultiver / récolter \\
\hline
\esp{el éxodo rural} & l'exode rural \\
\hline
\esp{la migración / emigrar} & la migration / émigrer \\
\hline
\esp{el transporte} & le transport \\
\hline
\esp{la carretera} & la route \\
\hline
\esp{el desarrollo local} & le développement local \\
\hline
\esp{la cooperativa} & la coopérative \\
\hline
\esp{el mercado / el pueblo} & le marché / le village \\
\hline
\esp{la tierra fértil} & la terre fertile \\
\hline
\end{tabularx}
\end{center}

\textbf{2. Grammaire : raconter le passé}
\begin{center}
\begin{tabularx}{\linewidth}{|X|X|X|}
\hline
\rowcolor{popgreenL} \textbf{Temps} & \textbf{Emploi} & \textbf{Exemple} \\
\hline
Pretérito indefinido & action ponctuelle, achevée, datée & \esp{Ayer llegué al pueblo.} « Hier, je suis arrivé au village. » \\
\hline
Pretérito imperfecto & description, habitude, décor, cause & \esp{El campo era verde y la gente trabajaba.} « La campagne était verte et les gens travaillaient. » \\
\hline
\end{tabularx}
\end{center}
\begin{itemize}
\item Indéfini régulier : \esp{hablar} $\to$ \esp{hablé, hablaste, habló, hablamos, hablasteis, hablaron} ; \esp{vivir} $\to$ \esp{viví, viviste, vivió, vivimos, vivisteis, vivieron}.
\item Imparfait régulier : \esp{-aba} (verbes en \esp{-ar}) ; \esp{-ía} (verbes en \esp{-er/-ir}). Seuls 3 irréguliers : \esp{ir} $\to$ \esp{iba}, \esp{ser} $\to$ \esp{era}, \esp{ver} $\to$ \esp{veía}.
\item Indéfini irréguliers fréquents : \esp{ser/ir} $\to$ \esp{fui, fue} ; \esp{hacer} $\to$ \esp{hice, hizo} ; \esp{tener} $\to$ \esp{tuve} ; \esp{estar} $\to$ \esp{estuve} ; \esp{poder} $\to$ \esp{pude}.
\item Signal horaire : \esp{ayer, el año pasado, de repente} $\to$ indéfini ; \esp{todos los días, mientras, antes} $\to$ imparfait.
\end{itemize}

\textbf{3. Méthodes express}
\begin{itemize}
\item \textbf{Lire un texte sur le monde rural} : repérer le lieu, les personnages, les problèmes (exode, routes) et les solutions (coopératives, projets locaux).
\item \textbf{Expliquer causes et conséquences} : utiliser \esp{porque, por eso, así que, como consecuencia} ; une cause = manque d'emploi, d'écoles, de routes ; une conséquence = vieillissement des villages, croissance des villes.
\item \textbf{Raconter au passé} : décor et habitudes à l'imparfait, actions principales à l'indéfini.
\end{itemize}
\end{bilan}

\begin{aretenir}[Checklist avant le Bac]
\begin{itemize}
\item[$\square$] Je connais le lexique : \esp{campo, agricultor, cosecha, éxodo rural, migración, transporte, carretera, desarrollo local, cooperativa}.
\item[$\square$] Je sais définir l'exode rural et citer au moins deux causes (emploi, éducation, routes) et deux conséquences.
\item[$\square$] Je conjugue sans faute l'indéfini des verbes réguliers en \esp{-ar, -er, -ir}.
\item[$\square$] Je connais les indéfinis irréguliers : \esp{fui, hice, hizo, tuve, estuve, pude, dije}.
\item[$\square$] Je conjugue l'imparfait : \esp{-aba / -ía}, et les irréguliers \esp{iba, era, veía}.
\item[$\square$] Je choisis correctement entre indéfini (action achevée) et imparfait (description, habitude).
\item[$\square$] Je sais traduire « il y a dix ans » par \esp{hace diez años} et « quand j'étais enfant » par \esp{cuando era niño}.
\item[$\square$] Je n'oublie pas les accents : \esp{habló, llegó, vivió, migración, carretera}.
\item[$\square$] Je distingue \esp{ser} et \esp{estar} au passé : \esp{era joven} (identité) / \esp{estaba cansado} (état).
\item[$\square$] Je sais rédiger un court récit au passé sur une expérience rurale (8 à 10 lignes, connecteurs : \esp{primero, luego, después, finalmente}).
\item[$\square$] Je peux citer une initiative de développement local au Cameroun (coopérative de cacao, marché rural, piste rurale).
\item[$\square$] Je relis ma production : accord sujet-verbe, accents, ponctuation espagnole ¿ ? ¡ !
\end{itemize}
\end{aretenir}

\findecours

\end{document}
